% La dissertation philosophique — Philosophie 1re Tech — Cours signé OnBuch+
% Programme officiel MINESEC. Compiler avec : tectonic N18-dissertation-philosophique.tex
\newcommand{\DOCMATIERE}{Philosophie}
\newcommand{\DOCNIVEAU}{1re Tech}
\newcommand{\DOCMODULE}{Philosophie – classe de Première technique}
\newcommand{\DOCLECON}{N18}
\newcommand{\DOCTITRE}{La dissertation philosophique}
% OnBuch+ — préambule « cours signés » ENRICHI (compilé avec tectonic / XeLaTeX)
% Système visuel « Pop » d'OnBuch+ : fond crème (jamais blanc), bordures encre
% épaisses, ombres « dures » décalées, titres Archivo Black en capitales.
% Version MATHÉMATIQUES : graphes (pgfplots/tikz), boîtes pédagogiques
% (définition, propriété, méthode, attention, le savais-tu, exercice résolu,
% exercice, TP, bilan), page de garde et signature OnBuch+ ancrée en bas.
%
% Macros attendues dans main.tex AVANT \input{preamble} :
%   \DOCMATIERE  \DOCNIVEAU  \DOCMODULE  \DOCLECON (numéro)  \DOCTITRE
\documentclass[11pt]{article}

\usepackage{fontspec}
\usepackage[french]{babel}
\usepackage[a4paper,margin=2cm,top=3.4cm,bottom=2.8cm,headheight=1.4cm,footskip=1.3cm]{geometry}
\usepackage{amsmath,amssymb,amsfonts}
\usepackage{mathtools} % \xmapsto, \coloneqq... souvent utilisés par les modèles
\usepackage{cancel} % simplifications barrées (bilans de forces)
% \abs{z} (module, valeur absolue) et \norm{u} (norme), utilisés par les modèles
\providecommand{\abs}{}\let\abs\relax\DeclarePairedDelimiter{\abs}{\lvert}{\rvert}
\providecommand{\norm}{}\let\norm\relax\DeclarePairedDelimiter{\norm}{\lVert}{\rVert}
\usepackage[table]{xcolor} % [table] : fournit aussi \rowcolors (alternance de lignes)
\usepackage{pagecolor}
\usepackage{tikz}
\usetikzlibrary{arrows.meta,positioning,calc,shapes.geometric,shapes.misc,shapes.arrows,decorations.pathmorphing,decorations.pathreplacing,decorations.markings,patterns,shadows,mindmap,trees,fit,backgrounds,matrix,intersections,angles,quotes}
\usepackage{pgfplots}
\usepackage[version=4]{mhchem} % équations nucléaires \ce{^{235}_{92}U}
\usepackage[european,siunitx]{circuitikz} % schémas électriques
\pgfplotsset{compat=1.18}
\usepgfplotslibrary{fillbetween,groupplots} % aires entre courbes (calcul intégral)
\usepackage{enumitem}
\usepackage{fancyhdr}
% [most] charge toutes les bibliothèques tcolorbox (listings, minted, raster,
% documentation...) et gonfle inutilement la mémoire TeX sur un document long
% (~300 boîtes) : on ne charge que ce qui est réellement utilisé.
\usepackage[skins,breakable]{tcolorbox}
\usepackage{graphicx}
\usepackage{colortbl}
\usepackage{booktabs}
\usepackage{tabularx}
\usepackage{multirow}
\usepackage{diagbox} % case d'en-tête diagonale des tableaux à double entrée
% Colonnes extensibles centrée / gauche / droite (C, L, R), souvent utilisées par les modèles
\newcolumntype{C}{>{\centering\arraybackslash}X}
\newcolumntype{L}{>{\raggedright\arraybackslash}X}
\newcolumntype{R}{>{\raggedleft\arraybackslash}X}
\usepackage{array}
\usepackage{multicol}
\usepackage{siunitx}
% parse-numbers=false : accepte aussi \SI{2,5 \times 10^{-3}}{...}
\sisetup{locale=FR,output-decimal-marker={,},group-minimum-digits=4,parse-numbers=false}
\DeclareSIUnit\ohm{\ensuremath{\symup{\Omega}}} % U+2126 absent de la police
\DeclareSIUnit\division{div} % réglages d'oscilloscope (ms/div, V/div)
\DeclareSIUnit\tr{tr} % tours (vitesse de rotation, ex. \SI{1500}{\tr\per\minute})
\DeclareSIUnit\molar{M} % concentration molaire (ex. \SI{3}{\molar}, \SI{1}{\micro\molar})
\DeclareSIUnit\an{an} % année (le modèle confond parfois avec \year, un compteur TeX) — ex. \SI{5}{\kilo\meter\per\an}
\DeclareSIUnit\FCFA{FCFA} % franc CFA (ex. \SI{500}{\FCFA\per\kilo\gram})
\DeclareSIUnit\degreeBrix{\textdegree Brix} % teneur en sucre d'un sirop/jus (réfractométrie)
\DeclareSIUnit\month{mois} % le modèle utilise parfois \month, un compteur TeX, comme unité de durée
\DeclareSIUnit\microliter{\micro\liter} % le modèle écrit parfois \microliter en un seul token
\DeclareSIUnit\million{millions} % dénombrement (ex. \SI{15}{\million\per\mL} spermatozoïdes)
\usepackage{qrcode}
% \degree : commande fréquemment utilisée par les modèles pour un angle ou une
% température en dehors de \SI{}{} (non fournie par les paquets chargés).
\providecommand{\degree}{^{\circ}}
% \qed (fin de démonstration), utilisé par les modèles sans amsthm
\providecommand{\qed}{\hfill$\square$}
% \barre : double barre des valeurs interdites dans un tableau de variations (tkz-tab)
\providecommand{\barre}{\ensuremath{\|}}
% environnement proof (amsthm non chargé) utilisé par les modèles
\ifdefined\proof\else\newenvironment{proof}[1][Démonstration]{\par\noindent\textit{#1.}\ }{\qed\par}\fi
\usepackage{lastpage}
\usepackage{hyperref}
\hypersetup{hidelinks}

% ============================================================
% Palette « Pop » OnBuch+ — mêmes hex que la classe P (plus_ui.dart)
% ============================================================
\definecolor{poporange}{HTML}{F59321}
\definecolor{popdark}{HTML}{E07A0C}
\definecolor{popcream}{HTML}{FFF6E8}
\definecolor{popink}{HTML}{1C1714}
\definecolor{poppurple}{HTML}{7A5AE0}
\definecolor{popgreen}{HTML}{2E9E6A}
\definecolor{poppink}{HTML}{E0457B}
\definecolor{popblue}{HTML}{2F7DE1}
\definecolor{popgold}{HTML}{C98A12}
\definecolor{popmuted}{HTML}{8A7F76}
\definecolor{popline}{HTML}{EADFD2}
\definecolor{popgray}{HTML}{8A7F76}
\colorlet{popgrey}{popgray}
% Alias tolérants (le modèle nomme parfois les couleurs différemment)
\colorlet{popwhite}{white}
\colorlet{popblack}{popink}
\colorlet{popred}{poppink}
\colorlet{popyellow}{popgold}
% Teintes claires (fonds de tableaux/encadrés) — le modèle nomme parfois
% les couleurs avec un suffixe L (ex. popblueL) sans les avoir définies.
\colorlet{poporangeL}{poporange!20}
\colorlet{popdarkL}{popdark!20}
\colorlet{poppurpleL}{poppurple!20}
\colorlet{popgreenL}{popgreen!20}
\colorlet{poppinkL}{poppink!20}
\colorlet{popblueL}{popblue!20}
\colorlet{popgoldL}{popgold!20}
\colorlet{popredL}{popred!20}
\colorlet{popgrayL}{popgray!20}
% Teintes claires pour fonds de tableaux / zones
\colorlet{popblueL}{popblue!10!white}
\colorlet{poporangeL}{poporange!14!white}
\colorlet{popgreenL}{popgreen!12!white}
\colorlet{poppurpleL}{poppurple!12!white}
\colorlet{poppinkL}{poppink!12!white}
\colorlet{popgoldL}{popgold!14!white}

\pagecolor{popcream}

% ============================================================
% Typographie
% ============================================================
\defaultfontfeatures{Ligatures=TeX}
\setmainfont{PlusJakartaSans-Medium.ttf}[Path=../../../fonts/,BoldFont=PlusJakartaSans-Bold.ttf]
% Maths en sans-serif (Fira Math) avec chiffres et lettres droites de Plus
% Jakarta, pour que les formules se fondent dans le texte.
\usepackage[math-style=ISO,bold-style=ISO]{unicode-math}
\setmathfont{FiraMath-Regular.otf}[Path=../../../fonts/]
\setmathfont{PlusJakartaSans-Medium.ttf}[Path=../../../fonts/,range={up/num,up/{Latin,latin},\mathup}]
\usepackage{icomma} % virgule décimale sans espace en mode math
% Caractères mathématiques Unicode absents des polices de texte (Plus Jakarta,
% Archivo Black) : rendus par leur équivalent mathématique, en texte comme en
% formule — sinon ils s'affichent en carré vide (ex. « Arithmétique dans ℤ »).
\usepackage{newunicodechar}
\newunicodechar{ℝ}{\ensuremath{\mathbb{R}}}
\newunicodechar{ℤ}{\ensuremath{\mathbb{Z}}}
\newunicodechar{ℂ}{\ensuremath{\mathbb{C}}}
\newunicodechar{ℕ}{\ensuremath{\mathbb{N}}}
\newunicodechar{ℚ}{\ensuremath{\mathbb{Q}}}
\newunicodechar{𝔻}{\ensuremath{\mathbb{D}}}
\newunicodechar{α}{\ensuremath{\alpha}}
\newunicodechar{β}{\ensuremath{\beta}}
\newunicodechar{γ}{\ensuremath{\gamma}}
\newunicodechar{δ}{\ensuremath{\delta}}
\newunicodechar{ε}{\ensuremath{\varepsilon}}
\newunicodechar{θ}{\ensuremath{\theta}}
\newunicodechar{λ}{\ensuremath{\lambda}}
\newunicodechar{μ}{\ensuremath{\mu}}
\newunicodechar{σ}{\ensuremath{\sigma}}
\newunicodechar{τ}{\ensuremath{\tau}}
\newunicodechar{φ}{\ensuremath{\varphi}}
\newunicodechar{ψ}{\ensuremath{\psi}}
\newunicodechar{ω}{\ensuremath{\omega}}
\newunicodechar{Σ}{\ensuremath{\Sigma}}
% majuscules produites par \MakeUppercase dans les titres (α-aminés -> Α)
\newunicodechar{Α}{\ensuremath{\alpha}}
\newunicodechar{Β}{\ensuremath{\beta}}
\newunicodechar{Γ}{\ensuremath{\Gamma}}
\newunicodechar{Δ}{\ensuremath{\Delta}}
\newunicodechar{Ε}{\ensuremath{\varepsilon}}
\newunicodechar{Θ}{\ensuremath{\Theta}}
\newunicodechar{Λ}{\ensuremath{\Lambda}}
\newunicodechar{Μ}{\ensuremath{\mu}}
\newunicodechar{Τ}{\ensuremath{\tau}}
\newunicodechar{Φ}{\ensuremath{\Phi}}
\newunicodechar{Ψ}{\ensuremath{\Psi}}
\newunicodechar{Ω}{\ensuremath{\Omega}}
\newunicodechar{↦}{\ensuremath{\mapsto}}
\newunicodechar{∈}{\ensuremath{\in}}
\newunicodechar{∉}{\ensuremath{\notin}}
\newunicodechar{∧}{\ensuremath{\wedge}}
\newunicodechar{⇔}{\ensuremath{\Leftrightarrow}}
\newunicodechar{⟺}{\ensuremath{\Longleftrightarrow}}
\newunicodechar{⇒}{\ensuremath{\Rightarrow}}
\newunicodechar{⟹}{\ensuremath{\Longrightarrow}}
\newunicodechar{∘}{\ensuremath{\circ}}
\newunicodechar{∪}{\ensuremath{\cup}}
\newunicodechar{∩}{\ensuremath{\cap}}
\newunicodechar{≡}{\ensuremath{\equiv}}
\newunicodechar{⊂}{\ensuremath{\subset}}
\newunicodechar{⊥}{\ensuremath{\perp}}
\newunicodechar{∃}{\ensuremath{\exists}}
\newunicodechar{∀}{\ensuremath{\forall}}
\newunicodechar{∛}{\ensuremath{\sqrt[3]{\,}}}
\newunicodechar{∜}{\ensuremath{\sqrt[4]{\,}}}
\newunicodechar{⃗}{\ensuremath{{}^{\to}}}
\newunicodechar{ⁿ}{\ifmmode^{n}\else\textsuperscript{\ensuremath{n}}\fi}
\newunicodechar{ˣ}{\ifmmode^{x}\else\textsuperscript{\ensuremath{x}}\fi}
\newunicodechar{ᵇ}{\ifmmode^{b}\else\textsuperscript{\ensuremath{b}}\fi}
\newunicodechar{ʳ}{\ifmmode^{r}\else\textsuperscript{\ensuremath{r}}\fi}
\newunicodechar{ᵘ}{\ifmmode^{u}\else\textsuperscript{\ensuremath{u}}\fi}
\newunicodechar{ʸ}{\ifmmode^{y}\else\textsuperscript{\ensuremath{y}}\fi}
\newunicodechar{ᵃ}{\ifmmode^{a}\else\textsuperscript{\ensuremath{a}}\fi}
\newunicodechar{ᵉ}{\ifmmode^{e}\else\textsuperscript{\ensuremath{e}}\fi}
\newunicodechar{ᵏ}{\ifmmode^{k}\else\textsuperscript{\ensuremath{k}}\fi}
\newunicodechar{ᵖ}{\ifmmode^{p}\else\textsuperscript{\ensuremath{p}}\fi}
\newunicodechar{ᴺ}{\ifmmode^{N}\else\textsuperscript{\ensuremath{N}}\fi}
\newunicodechar{ᶜ}{\ifmmode^{c}\else\textsuperscript{\ensuremath{c}}\fi}
\newunicodechar{ʰ}{\ifmmode^{h}\else\textsuperscript{\ensuremath{h}}\fi}
\newunicodechar{ᵠ}{\ifmmode^{q}\else\textsuperscript{\ensuremath{q}}\fi}
\newunicodechar{ₙ}{\ifmmode_{n}\else\textsubscript{\ensuremath{n}}\fi}
\newunicodechar{ₐ}{\ifmmode_{a}\else\textsubscript{\ensuremath{a}}\fi}
\newunicodechar{ᵢ}{\ifmmode_{i}\else\textsubscript{\ensuremath{i}}\fi}
\newunicodechar{ᵣ}{\ifmmode_{r}\else\textsubscript{\ensuremath{r}}\fi}
\newunicodechar{ₖ}{\ifmmode_{k}\else\textsubscript{\ensuremath{k}}\fi}
\newunicodechar{ᵦ}{\ifmmode_{\beta}\else\textsubscript{\ensuremath{\beta}}\fi}
\newunicodechar{ₓ}{\ifmmode_{x}\else\textsubscript{\ensuremath{x}}\fi}
\newfontfamily\popdisplay{ArchivoBlack-Regular.ttf}[Path=../../../fonts/]
\newfontfamily\poplogo{SpaceGrotesk-Bold.ttf}[Path=../../../fonts/]
\newfontfamily\popsemi{PlusJakartaSans-SemiBold.ttf}[Path=../../../fonts/]
\newfontfamily\popxbold{PlusJakartaSans-ExtraBold.ttf}[Path=../../../fonts/]
\newfontfamily\popxboldls{PlusJakartaSans-ExtraBold.ttf}[Path=../../../fonts/,LetterSpace=3.0]

\setlist[enumerate]{leftmargin=1.7em,itemsep=5pt,topsep=4pt}
\setlist[itemize]{leftmargin=1.7em,itemsep=4pt,topsep=4pt,label={\color{poporange}\textbullet}}
\setlength{\parindent}{0pt}
\setlength{\parskip}{5pt}
\renewcommand{\arraystretch}{1.25}

% pgfplots : style Pop par défaut
\pgfplotsset{
  every axis/.append style={axis line style={popink,line width=1pt},tick style={popink},
    grid style={popline,line width=0.5pt},label style={font=\small},tick label style={font=\footnotesize}},
  popcurve/.style={poporange,line width=1.8pt,no markers,smooth},
}
% Même style disponible dans un \draw TikZ simple (hors pgfplots)
\tikzset{popcurve/.style={poporange,line width=1.8pt}}
\tikzset{popgrid/.style={popline,very thin}}
\colorlet{popcurve}{poporange} % le nom du style est parfois utilisé comme couleur

% ============================================================
% Pill / squiggle
% ============================================================
\newcommand{\poppill}[3]{%
  \tikz[baseline=-0.55ex]\node[draw=popink,line width=1.2pt,rounded corners=2.6pt,%
    fill=#2,inner xsep=6.5pt,inner ysep=3pt]{%
    \popxboldls\fontsize{7}{7}\selectfont\color{#3}\MakeUppercase{#1}};%
}
\newcommand{\popsquiggle}[1]{%
  \tikz[baseline=-0.4ex]{
    \draw[#1,line width=2.6pt,line cap=round]
      plot [smooth] coordinates {(0,0.09) (0.34,0.01) (0.68,0.21) (1.02,0.01) (1.36,0.10)};
  }%
}
% Mot-clé surligné (marqueur orange sous le texte)
\newcommand{\cle}[1]{{\popxbold\color{popdark}#1}}

% ============================================================
% En-tête / pied
% ============================================================
\pagestyle{fancy}
\fancyhf{}
\renewcommand{\headrulewidth}{0pt}
\renewcommand{\footrulewidth}{0pt}
\lhead{\raisebox{-0.15ex}{\poplogo\fontsize{13}{13}\selectfont\color{popink}OnBuch\color{poporange}+}}
\rhead{\poppill{\DOCMATIERE\ \textbullet\ \DOCNIVEAU\ \textbullet\ Leçon \DOCLECON}{white}{popink}}
\lfoot{\popxbold\fontsize{7.6}{7.6}\selectfont\color{popmuted}\MakeUppercase{Cours signé OnBuch+}\ \textbullet\ {\popsemi\DOCTITRE}}
\rfoot{\tikz[baseline=-0.6ex]\node[rounded corners=8.5pt,fill=popink,minimum height=17pt,minimum width=17pt,inner xsep=5pt,inner sep=0]{\popxbold\fontsize{8}{8}\selectfont\color{popcream}\thepage\,/\,\pageref*{LastPage}};}

% ============================================================
% Titres
% ============================================================
\newcommand{\chaptitre}[2]{%
  \begin{tcolorbox}[enhanced,colback=poporange,colframe=popink,boxrule=2.6pt,arc=11pt,
      drop shadow=popink,left=15pt,right=15pt,top=12pt,bottom=11pt,before skip=3pt,after skip=18pt]
    \poppill{#2}{popink}{popcream}\par\vspace{9pt}
    {\raggedright\hyphenpenalty=10000\exhyphenpenalty=10000\popdisplay\fontsize{22}{24}\selectfont\color{popcream}\MakeUppercase{#1}\par}
  \end{tcolorbox}
}
% Section numérotée : pastille orange avec le numéro + titre Archivo
\newcounter{popsec}
\newcommand{\coursec}[1]{%
  \stepcounter{popsec}\setcounter{popsub}{0}%
  \par\vspace{16pt}\noindent
  \begin{minipage}{\linewidth}
  \tikz[baseline=-0.7ex]\node[draw=popink,line width=1.6pt,fill=poporange,rounded corners=4pt,minimum size=22pt,inner sep=0]{\popdisplay\fontsize{12}{12}\selectfont\color{popcream}\thepopsec};%
  \hspace{8pt}{\popdisplay\fontsize{15}{17}\selectfont\color{popink}\MakeUppercase{#1}}\par
  \vspace{4pt}\noindent{\color{poporange}\rule{36pt}{3.6pt}}
  \end{minipage}\par\nopagebreak\vspace{8pt}
}
\newcounter{popsub}
\newcommand{\courssub}[1]{%
  \stepcounter{popsub}%
  \par\vspace{9pt}\noindent{\popxbold\fontsize{11.5}{13}\selectfont\color{popink}{\color{poporange}\thepopsec.\thepopsub}\hspace{6pt}#1}\par\nopagebreak\vspace{4pt}
}

% ============================================================
% Boîtes pédagogiques — même recette : fond blanc, bordure épaisse colorée,
% ombre dure encre, badge plein. Titre optionnel [#1].
% ============================================================
\tcbset{popbase/.style={breakable,enhanced,colback=white,boxrule=2.3pt,arc=9pt,
  shadow={3pt}{-3pt}{0pt}{popink},left=14pt,right=14pt,top=11pt,bottom=11pt,before skip=11pt,after skip=15pt,
  fontupper=\popsemi\fontsize{10.6}{14.5}\selectfont\color{popink},
  fontlower=\popsemi\fontsize{10.6}{14.5}\selectfont\color{popink}}}
% Coche dessinée (le glyphe ✓ n'existe pas dans les polices du document)
\newcommand{\popcheck}[1]{\tikz[baseline=-0.5ex]\draw[#1,line width=1.6pt,line cap=round,line join=round](-0.13,0.0)--(-0.04,-0.09)--(0.13,0.1);}
\providecommand{\coche}{\popcheck{popgreen}}
\providecommand{\resultat}[1]{\par\smallskip\noindent\fcolorbox{popgreen}{popgreenL}{\parbox{\dimexpr\linewidth-2\fboxsep-2\fboxrule\relax}{\textbf{Résultat.} #1}}\par\smallskip}
\newcommand{\popboxhead}[3]{\poppill{#1}{#2}{white}\ifx\relax#3\relax\else\hspace{6pt}{\popxbold\fontsize{10.6}{12}\selectfont\color{popink}#3}\fi\par\vspace{7pt}}

\newtcolorbox{aretenir}[1][]{popbase,colframe=popgreen,before upper={\popboxhead{À retenir}{popgreen}{#1}}}
\newtcolorbox{exemplebox}[1][]{popbase,colframe=poporange,before upper={\popboxhead{Exemple}{poporange}{#1}}}
\newtcolorbox{definition}[1][]{popbase,colframe=popblue,before upper={\popboxhead{Définition}{popblue}{#1}}}
\newtcolorbox{propriete}[1][]{popbase,colframe=poppurple,before upper={\popboxhead{Propriété}{poppurple}{#1}}}
\newtcolorbox{methode}[1][]{popbase,colframe=popgold,before upper={\popboxhead{Méthode}{popgold}{#1}}}
\newtcolorbox{attention}[1][]{popbase,colframe=poppink,before upper={\popboxhead{Attention}{poppink}{#1}}}
\newtcolorbox{experience}[1][]{popbase,colframe=popblue,colback=popblueL,before upper={\popboxhead{Expérience}{popblue}{#1}}}
% « Le savais-tu ? » : carte violette pleine (contexte, culture, Cameroun)
\newtcolorbox{savaistu}[1][]{popbase,colback=poppurple,colframe=popink,
  fontupper=\popsemi\fontsize{10.4}{14.2}\selectfont\color{white},
  before upper={\poppill{Le savais-tu\,?}{white}{poppurple}\ifx\relax#1\relax\else\hspace{6pt}{\popxbold\color{white}#1}\fi\par\vspace{7pt}}}
% Exercice résolu : énoncé en haut, \tcblower puis la solution sur fond crème
\newcounter{popexo}\newcounter{popexr}
\newtcolorbox{exoresolu}[1][]{popbase,colframe=poporange,colbacklower=popcream,
  segmentation style={popink,line width=1.2pt,dashed},
  before upper={\stepcounter{popexr}\popboxhead{Exercice résolu \thepopexr}{poporange}{#1}},
  before lower={\poppill{Solution}{popink}{popcream}\par\vspace{6pt}}}
% Exercice (non résolu en place) — la correction est dans la partie Corrigés
\newtcolorbox{exercice}[1][]{popbase,colframe=popink,
  before upper={\stepcounter{popexo}\popboxhead{Exercice \thepopexo}{popink}{#1}}}
% Corrigé
\newtcolorbox{corrige}[1][]{popbase,colframe=popgreen,colback=popgreenL,
  before upper={\popboxhead{Corrigé}{popgreen}{#1}}}
% Objectifs / prérequis / bilan
\newtcolorbox{objectifs}{popbase,colframe=popink,colback=poporangeL,
  before upper={\popboxhead{Objectifs de la leçon}{popink}{}}}
\newtcolorbox{prerequis}{popbase,colframe=popink,colback=popblueL,
  before upper={\popboxhead{Prérequis}{popblue}{}}}
\newtcolorbox{bilan}{popbase,colframe=popink,colback=white,boxrule=2.8pt,
  before upper={\popboxhead{Fiche bilan}{poporange}{L'essentiel à connaître pour l'examen}}}
% Figure encadrée avec légende
\newtcolorbox{popfigure}[1][]{enhanced,breakable=false,colback=white,colframe=popline,boxrule=1.4pt,arc=8pt,
  left=10pt,right=10pt,top=10pt,bottom=8pt,before skip=10pt,after skip=12pt,halign=center,
  lower separated=false,halign lower=center,
  fontlower=\popsemi\fontsize{9}{11.5}\selectfont\color{popmuted},
  #1}
\newcommand{\legende}[1]{\par\vspace{6pt}{\popsemi\fontsize{9}{11.5}\selectfont\color{popmuted}#1}}

% ============================================================
% Signature OnBuch+
% ============================================================
\newcommand{\poplogoinline}[1]{{\poplogo\fontsize{#1}{#1}\selectfont\color{popink}OnBuch\color{poporange}+}}
\newcommand{\signatureonbuch}{%
  \begin{tcolorbox}[enhanced,colback=popink,colframe=popink,boxrule=2.2pt,arc=10pt,
      drop shadow=poporange,left=14pt,right=14pt,top=10pt,bottom=10pt,before skip=0pt,after skip=0pt]
    \tikz[baseline=-0.7ex]\node[circle,fill=popgreen,draw=popcream,line width=1.3pt,minimum size=22pt,inner sep=0]{\popcheck{white}};%
    \hspace{9pt}\begin{minipage}[c]{0.62\linewidth}
      {\popxbold\fontsize{12}{13}\selectfont\color{popcream}Signé\ \textbullet\ {\poplogo OnBuch\color{poporange}+}}\par\vspace{2pt}
      {\raggedright\popsemi\fontsize{8}{10.5}\selectfont\color{popline}\MakeUppercase{Équipe pédagogique OnBuch+}\\\MakeUppercase{Conforme au programme officiel MINESEC}\par}
    \end{minipage}\hfill
    {\poplogo\fontsize{22}{22}\selectfont\color{popcream}OnBuch\color{poporange}+}
  \end{tcolorbox}%
}

% ============================================================
% Page de garde
% #1 titre · #2 matière · #3 niveau · #4 module · #5 n° leçon · #6 durée ·
% #7 description / objectifs (texte ou itemize)
% ============================================================
\newcommand{\pagegarde}[7]{%
  \thispagestyle{empty}
  \newgeometry{a4paper,margin=1.7cm,top=1.6cm,bottom=1.4cm}%
  \begin{tikzpicture}[remember picture,overlay]
    \fill[poporange!18] ([xshift=-3.2cm,yshift=-3.6cm]current page.north east) circle (4.6cm);
    \fill[poppurple!14] ([xshift=2.4cm,yshift=7.6cm]current page.south west) circle (3.8cm);
  \end{tikzpicture}%
  {\poplogo\fontsize{30}{30}\selectfont\color{popink}OnBuch\color{poporange}+}\hfill
  \raisebox{0.6ex}{\poppill{Cours signé \textbullet\ Premium}{popink}{popcream}}\par
  \vspace{1pt}\popsquiggle{poppurple}\par
  \vspace{8pt}
  \begin{tcolorbox}[enhanced,colback=poporange,colframe=popink,boxrule=2.8pt,arc=13pt,
      drop shadow=popink,left=18pt,right=18pt,top=12pt,bottom=13pt,after skip=10pt]
    \poppill{#2 \textbullet\ #3}{popink}{popcream}\hspace{5pt}\poppill{#4}{white}{popink}\par\vspace{8pt}
    {\popdisplay\fontsize{14}{14}\selectfont\color{popink}LEÇON #5}\par\vspace{4pt}
    {\raggedright\hyphenpenalty=10000\exhyphenpenalty=10000\popdisplay\fontsize{25}{27}\selectfont\color{popcream}\MakeUppercase{#1}\par}
    \vspace{8pt}
    {\popxbold\fontsize{9.5}{9.5}\selectfont\color{popink}\MakeUppercase{Durée conseillée : #6}}
  \end{tcolorbox}
  \begin{tcolorbox}[enhanced,colback=white,colframe=popink,boxrule=2.2pt,arc=10pt,
      drop shadow=popink,left=16pt,right=16pt,top=10pt,bottom=10pt,after skip=8pt]
    \poppill{Dans cette leçon}{poppurple}{white}\par\vspace{5pt}
    \popsemi\fontsize{9.3}{12.6}\selectfont\color{popink}#7
  \end{tcolorbox}
  \noindent\begin{minipage}[t]{0.487\linewidth}
    \begin{tcolorbox}[enhanced,colback=popgreen,colframe=popink,boxrule=2.2pt,arc=10pt,
        drop shadow=popink,left=11pt,right=11pt,top=10pt,bottom=10pt,height=3.1cm,valign=center]
      \tcbox[colback=white,colframe=popink,boxrule=1.3pt,arc=5pt,boxsep=0pt,nobeforeafter,
        left=3pt,right=3pt,top=3pt,bottom=3pt,valign=center]{\qrcode[height=1.9cm]{https://play.google.com/store/apps/details?id=com.luvvix.onbuch&hl=fr}}%
      \hspace{8pt}\begin{minipage}[c]{0.5\linewidth}
        {\popdisplay\fontsize{11}{12.5}\selectfont\color{white}\MakeUppercase{Télécharge OnBuch}}\par\vspace{4pt}
        {\raggedright\popsemi\fontsize{8.4}{11}\selectfont\color{white}Gratuit \textbullet\ Google Play\\Tuteur IA, annales, concours, résultats\par}
      \end{minipage}
    \end{tcolorbox}
  \end{minipage}\hfill
  \begin{minipage}[t]{0.487\linewidth}
    \begin{tcolorbox}[enhanced,colback=poppurple,colframe=popink,boxrule=2.2pt,arc=10pt,
        drop shadow=popink,left=11pt,right=11pt,top=10pt,bottom=10pt,height=3.1cm,valign=center]
      \tikz[baseline=-0.7ex]\node[circle,fill=white,draw=popink,line width=1.3pt,minimum size=1.6cm,inner sep=0]{\popdisplay\fontsize{24}{24}\selectfont\color{poppurple}+};%
      \hspace{8pt}\begin{minipage}[c]{0.6\linewidth}
        {\popdisplay\fontsize{11}{12.5}\selectfont\color{white}\MakeUppercase{Passe à OnBuch+}}\par\vspace{4pt}
        {\raggedright\popsemi\fontsize{8.4}{11}\selectfont\color{white}Tous les cours signés, Tuteur IA illimité, fiches PDF — onglet OnBuch+ de l'app\par}
      \end{minipage}
    \end{tcolorbox}
  \end{minipage}
  \vspace{8pt}
  \popcontenu
  \vfill
  \signatureonbuch
  \restoregeometry
  \newpage
}

% Rangée « Ce cours contient » (page de garde)
\newcommand{\poptile}[3]{% #1 couleur #2 gros texte #3 légende
  \begin{tcolorbox}[enhanced,colback=#1,colframe=popink,boxrule=2pt,arc=9pt,shadow={2.5pt}{-2.5pt}{0pt}{popink},
      left=6pt,right=6pt,top=9pt,bottom=9pt,halign=center,valign=center,height=1.75cm,width=\linewidth,nobeforeafter]
    {\popdisplay\fontsize{12.5}{13}\selectfont\color{white}\MakeUppercase{#2}}\par\vspace{3pt}
    {\centering\popsemi\fontsize{7.8}{9.5}\selectfont\color{white}#3\par}
  \end{tcolorbox}}
\newcommand{\popcontenu}{%
  \par\noindent{\popxboldls\fontsize{7.5}{7.5}\selectfont\color{popmuted}CE COURS SIGNÉ CONTIENT}\par\vspace{7pt}
  \noindent\begin{minipage}[t]{0.235\linewidth}\poptile{popblue}{Cours}{complet, illustré, pas à pas}\end{minipage}\hfill
  \begin{minipage}[t]{0.235\linewidth}\poptile{poporange}{Méthodes}{et exercices résolus}\end{minipage}\hfill
  \begin{minipage}[t]{0.235\linewidth}\poptile{poppink}{Exercices}{type examen + corrigés}\end{minipage}\hfill
  \begin{minipage}[t]{0.235\linewidth}\poptile{popgreen}{Fiche bilan}{l'essentiel à retenir}\end{minipage}\par}

% Sommaire
\newtcolorbox{sommaire}{enhanced,colback=white,colframe=popink,boxrule=2.2pt,arc=10pt,
  drop shadow=popink,left=18pt,right=18pt,top=13pt,bottom=13pt,before skip=6pt,after skip=20pt,
  before upper={\poppill{Sommaire}{poppurple}{white}\par\vspace{9pt}\popsemi\fontsize{10.6}{14.5}\selectfont\color{popink}}}

% Signature de fin de cours — poussée tout en bas de la dernière page
\newcommand{\findecours}{%
  \par\vfill\nopagebreak
  \begin{center}\popsquiggle{poporange}\end{center}\vspace{4pt}
  \signatureonbuch
}
\AtBeginDocument{\def\llbracket{\mathopen{[\mkern-3.5mu[}}\def\rrbracket{\mathclose{]\mkern-3.5mu]}}} % intervalles d'entiers (glyphe ⟦ absent de la police)
% Glyphes absents de Fira Math : redéfinis à partir de glyphes disponibles
\AtBeginDocument{%
  \def\setminus{\mathbin{\backslash}}\let\smallsetminus\setminus
  \def\vdots{\mathord{\vbox{\baselineskip3.2pt\lineskiplimit0pt\kern4pt\hbox{.}\hbox{.}\hbox{.}}}}%
  \def\ddots{\mathinner{\mkern1mu\raise6.5pt\hbox{.}\mkern2mu\raise3.6pt\hbox{.}\mkern2mu\raise0.7pt\hbox{.}\mkern1mu}}%
  \def\triangle{\mathord{\scalebox{0.9}{$\symup{\Delta}$}}}%
  \def\nabla{\mathord{\raisebox{\depth}{\scalebox{1}[-1]{$\symup{\Delta}$}}}}%
  \def\checkmark{\popcheck{popdark}}%
  \def\star{\mathbin{\ast}}\def\bigstar{\mathord{\ast}}%
  \def\diamond{\mathbin{\scalebox{0.6}{$\Diamond$}}}%
  \def\bigcup{\mathop{\mathchoice{\raisebox{-0.3em}{\scalebox{1.7}{$\cup$}}}{\scalebox{1.25}{$\cup$}}{\cup}{\cup}}\displaylimits}%
  \def\bigcap{\mathop{\mathchoice{\raisebox{-0.3em}{\scalebox{1.7}{$\cap$}}}{\scalebox{1.25}{$\cap$}}{\cap}{\cap}}\displaylimits}%
  \def\lesseqqgtr{\mathrel{\lessgtr}}\def\bigoplus{\mathop{\oplus}\displaylimits}%
}
\providecommand{\ssi}{si et seulement si} % abréviation fréquente
\AtBeginDocument{\providecommand{\wideparen}{\overparen}} % arc de cercle

%% FIGFIX-BEGIN v5 — correctifs globaux des figures (docs/onbuch-generation/tools/figfix)
\usepackage{adjustbox}\usepackage{etoolbox}\tcbuselibrary{raster}
% toute figure TikZ/pgfplots plus large que la ligne est réduite à la largeur disponible
\BeforeBeginEnvironment{tikzpicture}{\begin{adjustbox}{max width=\linewidth}}
\AfterEndEnvironment{tikzpicture}{\end{adjustbox}}
% légendes pgfplots lisibles par-dessus les courbes
\pgfplotsset{every axis/.append style={legend style={fill=white,fill opacity=0.92,text opacity=1,draw=black!15,inner sep=3pt}}}
% étiquettes d'arêtes (syntaxe edge["..."]) avec fond blanc
\tikzset{every edge quotes/.append style={fill=white,inner sep=1.2pt}}
%% FIGFIX-END
\begin{document}

\pagegarde{La dissertation philosophique}{Philosophie}{1re Tech}{Philosophie – classe de Première technique}{N18}{5 séances (fiche de progression harmonisée)}{Cette leçon présente la dissertation philosophique comme exercice méthodique de réflexion : de l'analyse du sujet à la conclusion, en passant par la confrontation dialectique de la thèse et de l'antithèse. Elle outille l'élève de Première technique pour rédiger une argumentation cohérente, exigence centrale de l'épreuve de philosophie au Baccalauréat.
\vspace{4pt}
\begin{multicols}{2}\small
\begin{itemize}
\item À la fin de cette leçon, je sais définir la dissertation philosophique et distinguer ses types (explication de texte, dissertation de sujet)
\item Analyser un sujet de dissertation en dégageant la problématique et les enjeux
\item Rédiger une introduction complète (accroche, analyse des termes, problématique, annonce du plan)
\item Construire une thèse argumentée et illustrée par des exemples concrets
\item Formuler une antithèse qui remet en cause la thèse de manière rigoureuse
\item Élaborer une synthèse qui dépasse le conflit des positions et rédiger une conclusion
\end{itemize}\end{multicols}}

\begin{sommaire}
\begin{multicols}{2}
\begin{enumerate}
\item Généralités sur la dissertation philosophique
\item L'introduction
\item La thèse
\item L'antithèse
\item La synthèse et la conclusion
\item Application guidée : rédiger une dissertation complète
\item Activité d'intégration
\item Méthodes et exercices résolus
\item Exercices
\item Corrigés des exercices
\item Fiche bilan
\end{enumerate}
\end{multicols}
\end{sommaire}

\begin{objectifs}
\begin{itemize}
\item À la fin de cette leçon, je sais définir la dissertation philosophique et distinguer ses types (explication de texte, dissertation de sujet)
\item Analyser un sujet de dissertation en dégageant la problématique et les enjeux
\item Rédiger une introduction complète (accroche, analyse des termes, problématique, annonce du plan)
\item Construire une thèse argumentée et illustrée par des exemples concrets
\item Formuler une antithèse qui remet en cause la thèse de manière rigoureuse
\item Élaborer une synthèse qui dépasse le conflit des positions et rédiger une conclusion
\item Appliquer la méthode dialectique à des situations concrètes de la vie courante camerounaise
\item Rédiger et justifier une démarche argumentative complète en temps limité
\end{itemize}
\end{objectifs}

\begin{prerequis}
\begin{itemize}
\item Notion d'argument et de raisonnement (français, classes de Seconde)
\item Distinction entre opinion et vérité (leçons antérieures de philosophie)
\item Capacité à rédiger un paragraphe argumenté (français)
\item Notions de logique élémentaire : cohérence, contradiction
\end{itemize}
\end{prerequis}

\newpage

\coursec{Généralités sur la dissertation philosophique}

\courssub{Situation d'accroche et problématique}

Lors d'un débat au lycée de Ngaoundéré, deux élèves s'affrontent : l'un affirme que \og le téléphone portable détruit la jeunesse camerounaise \fg{}, l'autre soutient le contraire. Le proviseur, qui les écoute, leur demande alors : \og Avez-vous examiné le problème sous toutes ses faces avant de conclure ? \fg{} Cette scène quotidienne illustre l'exigence fondamentale de la dissertation philosophique : ne pas se contenter d'une opinion, mais construire une réflexion rigoureuse en trois moments.

Chacun de nous ressemble à ces deux élèves : nous avons des opinions, nous les défendons avec passion, mais souvent sans méthode. Or, au Baccalauréat technique, le correcteur n'attend ni votre opinion ni votre passion : il attend une \cle{réflexion organisée}. La problématique de cette leçon peut donc se formuler ainsi : \emph{comment passer d'une simple dispute d'opinions à une argumentation méthodique digne d'un philosophe ?} Pour répondre à cette question, nous devons d'abord définir la dissertation, situer sa place dans l'épreuve de philosophie, distinguer les deux types d'épreuves proposés au Baccalauréat, identifier les qualités exigées du candidat, présenter la structure canonique du devoir et enfin comprendre les critères d'évaluation.

\courssub{Définition de la dissertation philosophique}

Commençons par une intuition simple. Quand deux personnes se disputent, chacune répète son point de vue sans écouter l'autre : c'est un dialogue de sourds. La dissertation, au contraire, oblige à \emph{faire dialoguer les points de vue opposés en soi-même}. Celui qui disserte doit être à la fois l'avocat du \og pour \fg{} et l'avocat du \og contre \fg{}, avant de devenir le juge qui tranche en connaissance de cause.

\begin{definition}[La dissertation philosophique]
La \cle{dissertation philosophique} est un exercice écrit de réflexion personnelle et argumentée sur un problème philosophique, dans lequel l'élève examine méthodiquement une question en confrontant des points de vue opposés (la thèse et l'antithèse) pour parvenir, par une synthèse, à une position réfléchie et justifiée.
\end{definition}

Trois éléments de cette définition méritent d'être soulignés :

\begin{itemize}
\item C'est une réflexion \cle{personnelle} : le candidat ne récite pas un cours, il pense par lui-même en mobilisant les notions et les auteurs étudiés.
\item C'est une réflexion \cle{argumentée} : chaque affirmation doit être soutenue par une raison, un exemple ou une référence philosophique.
\item C'est une réflexion \cle{méthodique} : elle suit un plan rigoureux, le plan dialectique (thèse, antithèse, synthèse), qui reproduit le mouvement naturel de la pensée qui cherche la vérité.
\end{itemize}

L'étymologie elle-même est éclairante : le mot \emph{dissertation} vient du latin \emph{dissertare}, qui signifie \og discuter, examiner en détail \fg{}. Disserter, c'est donc d'abord \emph{examiner}, non pas affirmer. Comme le disait le philosophe français Alain, la philosophie n'est pas un savoir à réciter mais une manière de penser : \og penser, c'est dire non \fg{}, c'est-à-dire refuser les évidences toutes faites pour les soumettre à l'examen critique.

\courssub{Origine et place de la dissertation dans l'épreuve de philosophie}

La dissertation est un exercice d'origine scolaire française, hérité de la tradition rhétorique gréco-latine : dans l'Antiquité, la \emph{disputatio} médiévale consistait déjà à débattre d'une question en examinant les arguments contraires avant de conclure. Introduite dans l'enseignement secondaire, elle est devenue l'exercice roi de la classe de philosophie, car elle vérifie à la fois la culture, la logique et l'expression écrite de l'élève.

Au Cameroun, l'épreuve de philosophie au Probatoire et au Baccalauréat technique évalue la capacité du candidat à problématiser, à argumenter et à rédiger. La dissertation y occupe une place centrale : c'est l'épreuve la plus fréquemment choisie par les candidats, car elle porte directement sur les notions du programme (le devoir, la liberté, la conscience, le bonheur, la technique, l'État, etc.).

\begin{savaistu}[La philosophie au Baccalauréat technique camerounais]
Au Baccalauréat technique camerounais, la philosophie possède un coefficient élevé, généralement de 2 à 4 selon les séries. Une bonne note en philosophie peut donc faire basculer un résultat. Or la dissertation est l'épreuve la plus fréquente et la plus choisie : maîtriser sa méthode, c'est sécuriser des points précieux pour l'obtention du diplôme.
\end{savaistu}

\courssub{Les deux types d'épreuves : dissertation de sujet et explication de texte}

Au Baccalauréat, le candidat choisit en général entre deux sujets : une \cle{dissertation de sujet} (une question ou une citation à discuter) et une \cle{explication de texte} (un court passage d'un philosophe à commenter). Il est utile de connaître les deux, même si cette leçon se concentre sur la dissertation.

La \emph{dissertation de sujet} part d'une question (par exemple : \og La technique est-elle un danger pour l'homme ? \fg{}) ou d'une citation à discuter. Le candidat construit librement sa réflexion en mobilisant sa culture philosophique.

L'\emph{explication de texte} part d'un extrait d'œuvre (par exemple un passage de Platon ou de Descartes). Le candidat doit dégager la thèse de l'auteur, expliquer sa démarche et la discuter. Elle exige moins de culture personnelle mais plus de fidélité au texte.

\begin{center}
\begin{tabularx}{\linewidth}{|l|X|X|}
\hline
\rowcolor{popblueL} \textbf{Critère} & \textbf{Dissertation de sujet} & \textbf{Explication de texte} \\
\hline
Point de départ & Une question ou une citation à discuter & Un extrait d'œuvre philosophique \\
\hline
Travail principal & Problématiser, construire un plan dialectique & Dégager la thèse de l'auteur, suivre son argumentation \\
\hline
Culture exigée & Forte (notions, auteurs, exemples personnels) & Modérée (le texte fournit la matière) \\
\hline
Liberté du candidat & Grande liberté de réflexion & Fidélité obligée au texte \\
\hline
Piège principal & Réciter le cours sans répondre au sujet & Paraphraser le texte sans l'expliquer \\
\hline
\end{tabularx}
\end{center}

\begin{attention}[Ne pas confondre la dissertation avec un exposé ou un commentaire littéraire]
L'erreur la plus fréquente consiste à confondre la dissertation philosophique avec le commentaire littéraire ou le simple exposé de connaissances. Réciter tout son cours sur la liberté, sans répondre précisément à la question posée, ne vaut presque rien : le correcteur sanctionne le \og hors-sujet \fg{}. De même, la dissertation n'est ni un récit, ni une déclamation oratoire, ni un sermon moralisateur. Elle est une \emph{discussion argumentée d'un problème}. Chaque phrase doit servir la démonstration.
\end{attention}

\courssub{Les qualités exigées du candidat}

Pour réussir, quatre qualités sont indispensables :

\begin{enumerate}
\item La \cle{clarté} : le devoir doit se lire sans effort. Une phrase claire est courte, précise, sans mots inutiles. Comme le voulait Descartes dans sa méthode, il faut \og conduire par ordre ses pensées \fg{}.
\item La \cle{cohérence} : les parties du devoir doivent s'enchaîner logiquement. La thèse prépare l'antithèse, l'antithèse appelle la synthèse, la synthèse répond à la problématique de l'introduction. Il ne doit y avoir ni contradiction interne ni rupture.
\item La \cle{rigueur logique} : chaque argument doit être valide et chaque exemple doit illustrer réellement l'idée défendue. On évite les affirmations gratuites (\og tout le monde sait que... \fg{}).
\item La \cle{culture philosophique} : le candidat doit citer brièvement et à bon escient les auteurs du programme (Platon, Aristote, Descartes, Kant, Sartre, etc.), sans transformer le devoir en catalogue de noms.
\end{enumerate}

\courssub{La structure canonique de la dissertation}

La dissertation philosophique obéit à une architecture immuable, que l'on peut comparer à un sablier : on part du large (l'accroche), on resserre vers le problème, puis on déploie la discussion avant de refermer sur la réponse.

\begin{popfigure}
\begin{center}
\begin{tikzpicture}[font=\small, scale=0.9]
% Structure en sablier : 5 trapèzes empilés
% Introduction (large)
\fill[popblueL] (0,5.5) -- (7,5.5) -- (6.2,4.5) -- (0.8,4.5) -- cycle;
\node[align=center, text width=5.5cm] at (3.5,5.0) {\textbf{Introduction} : accroche, définition des termes, problématique, annonce du plan};
% Thèse
\fill[popgreenL] (0.8,4.2) -- (6.2,4.2) -- (5.6,3.2) -- (1.4,3.2) -- cycle;
\node[align=center, text width=4.5cm] at (3.5,3.7) {\textbf{Thèse} : première réponse argumentée (arguments + exemples)};
% Antithèse
\fill[poporangeL] (1.4,2.9) -- (5.6,2.9) -- (5.0,1.9) -- (2.0,1.9) -- cycle;
\node[align=center, text width=3.5cm] at (3.5,2.4) {\textbf{Antithèse} : réponse opposée, tout aussi argumentée};
% Synthèse
\fill[poppurpleL] (2.0,1.6) -- (5.0,1.6) -- (4.4,0.6) -- (2.6,0.6) -- cycle;
\node[align=center, text width=2.5cm] at (3.5,1.1) {\textbf{Synthèse} : dépassement du conflit, position nuancée};
% Conclusion
\fill[popgoldL] (2.6,0.3) -- (4.4,0.3) -- (4.0,-0.5) -- (3.0,-0.5) -- cycle;
\node[align=center, text width=1.5cm] at (3.5,-0.1) {\textbf{Conclusion} : bilan, réponse, ouverture};
% Flèche indicatrice
\draw[->,thick,popdark] (7.5,5.0) -- (7.5,-0.5) node[midway,right,align=left] {Resserrement progressif\\vers la réponse finale};
\end{tikzpicture}
\end{center}
\legende{Structure canonique de la dissertation philosophique : le devoir se resserre progressivement, de la mise en route de la réflexion (introduction) jusqu'à la réponse finale (conclusion), en passant par la confrontation de la thèse et de l'antithèse dépassée dans la synthèse.}
\end{popfigure}

\begin{aretenir}[Les cinq moments du devoir]
\begin{enumerate}
\item \textbf{Introduction} : accroche, définition des termes, problématique, annonce du plan.
\item \textbf{Thèse} : première réponse au problème, avec arguments et exemples.
\item \textbf{Antithèse} : réponse opposée, tout aussi bien argumentée.
\item \textbf{Synthèse} : dépassement des deux positions, réponse personnelle nuancée.
\item \textbf{Conclusion} : bilan clair de la réflexion, réponse à la problématique, éventuelle ouverture.
\end{enumerate}
\end{aretenir}

\courssub{Les critères d'évaluation au Baccalauréat}

Le correcteur du Baccalauréat technique juge le devoir selon trois grands critères :

\begin{itemize}
\item La \cle{pertinence de la problématique} : le candidat a-t-il compris le sujet ? A-t-il dégagé le vrai problème philosophique ? Un devoir hors sujet, même bien écrit, est lourdement sanctionné.
\item La \cle{qualité de l'argumentation} : les idées sont-elles organisées selon le plan dialectique ? Les arguments sont-ils solides, les exemples convaincants, les références exactes ?
\item La \cle{correction de la langue} : orthographe, grammaire, syntaxe, présentation. Une langue négligée brouille la pensée et fait perdre des points précieux.
\end{itemize}

\begin{methode}[Gérer son temps pendant l'épreuve (3 heures)]
\begin{enumerate}
\item \textbf{30 minutes} : lecture attentive du sujet, analyse des termes, formulation de la problématique, construction du plan détaillé au brouillon.
\item \textbf{2 heures} : rédaction complète du devoir au propre, en suivant le plan sans s'en écarter.
\item \textbf{15 minutes} : relecture attentive (orthographe, transitions, cohérence) et dernières corrections.
\end{enumerate}
Ne jamais rédiger sans plan : dix minutes de plan en sauvent trente de rédaction.
\end{methode}

\begin{exoresolu}[Identifier les qualités d'une dissertation]
\textbf{Énoncé.} Un élève de Première technique écrit, sur le sujet \og Faut-il toujours dire la vérité ? \fg{}, le passage suivant : \og La vérité est très importante dans la vie. Au village, les anciens disent que la vérité libère. Moi je pense qu'il faut toujours dire la vérité car le mensonge est mauvais. \fg{} Relevez deux qualités présentes et deux défauts de ce passage au regard des exigences de la dissertation.
\tcblower
\textbf{Solution.} \emph{Qualités} : (1) le passage est clair et se comprend facilement ; (2) il contient un début d'argument (le mensonge est mauvais) et une référence culturelle (la sagesse des anciens). \emph{Défauts} : (1) l'élève affirme son opinion personnelle (\og moi je pense \fg{}) sans examiner le point de vue opposé : il n'y a ni antithèse ni problématisation, alors que la dissertation exige la confrontation des thèses ; (2) l'argument est gratuit : pourquoi le mensonge est-il mauvais ? Aucune raison philosophique, aucun auteur (par exemple Kant, qui condamne le mensonge, ou le cas du mensonge pour sauver une vie) n'est mobilisé. Ce passage relève donc de la simple dispute d'opinions, comme au lycée de Ngaoundéré, et non encore de la dissertation philosophique.
\end{exoresolu}

En résumé, la dissertation philosophique est un exercice codifié mais libérateur : codifié par sa structure en cinq moments, libérateur parce qu'il apprend à penser par soi-même. Les sections suivantes détailleront chacun de ces moments, en commençant par l'introduction, la porte d'entrée du devoir.

\coursec{L'introduction}

Nous avons vu, dans la section précédente, que la dissertation philosophique est un exercice de réflexion problématisée, organisée en trois moments (thèse, antithèse, synthèse) et précédée d'une introduction. Mais avant même d'argumenter, il faut \emph{ouvrir} la réflexion. C'est le rôle de l'introduction : elle est la porte d'entrée du devoir, le premier contact entre le candidat et le correcteur. Une dissertation peut être comparée à un voyage : l'introduction est le moment où l'on présente la destination, où l'on explique pourquoi ce voyage vaut la peine, et où l'on annonce les étapes du parcours. Négliger l'introduction, c'est partir en voyage sans dire où l'on va.

\courssub{Fonction de l'introduction : présenter le problème et engager la réflexion}

Commençons par une intuition simple. Lorsqu'un chef de village réunit les notables pour trancher un différend, il ne commence jamais par donner son avis. Il rappelle d'abord les faits, il définit les termes du conflit, il pose la question qui divise, puis il invite chacun à s'exprimer. L'introduction philosophique procède exactement de la même manière : elle met en place le débat avant de le conduire.

\begin{definition}[L'introduction philosophique]
L'\cle{introduction} est la première partie de la dissertation. Elle présente le sujet, définit ses termes, dégage la \cle{problématique} (la question centrale) et annonce le plan du développement. Elle comporte classiquement \textbf{quatre étapes} : l'amorce, l'analyse des termes, la problématique et l'annonce du plan.
\end{definition}

L'introduction remplit deux fonctions inséparables :
\begin{itemize}
\item une fonction \textbf{expositive} : elle présente le problème au lecteur, montre que le sujet n'est pas une question banale mais une véritable difficulté de pensée ;
\item une fonction \textbf{engagée} : elle engage la réflexion, c'est-à-dire qu'elle crée l'attente, donne au lecteur l'envie de suivre l'argumentation et annonce la direction que prendra le devoir.
\end{itemize}

\begin{attention}[Ce que l'introduction n'est pas]
L'introduction n'est ni un résumé du cours, ni un exposé de toutes ses connaissances, ni la réponse au sujet. Celui qui donne dès l'introduction sa réponse définitive (« La technique libère l'homme, et voici pourquoi... ») supprime le problème au lieu de le poser : il n'a plus rien à démontrer, et son devoir perd sa dimension argumentative.
\end{attention}

\courssub{Première étape : l'amorce ou l'accroche}

L'amorce est la première phrase du devoir. Son rôle est de \emph{conduire naturellement} le lecteur vers le sujet, en partant de quelque chose de plus large que lui : un fait d'expérience commune, un constat, une référence culturelle ou historique.

\begin{definition}[L'amorce]
L'\cle{amorce} (ou accroche) est la mise en route de l'introduction. Elle part d'un fait général, d'une évidence partagée ou d'une référence, pour aboutir au sujet. Elle doit être brève (deux à quatre phrases) et rester en lien direct avec le sujet.
\end{definition}

Il existe trois grands types d'amorces :
\begin{enumerate}
\item \textbf{Le fait d'expérience ou le constat} : on part d'une réalité observable par tous. Exemple : « Chaque matin, des millions de Camerounais consultent leur téléphone portable avant même de saluer leur voisin. »
\item \textbf{La référence culturelle ou historique} : on part d'un auteur, d'un mythe, d'un événement. Exemple : « Depuis Prométhée dérobant le feu aux dieux, la technique apparaît comme le propre de l'homme. »
\item \textbf{Le lieu commun à interroger} : on part d'une opinion répandue que l'on va soumettre à l'examen. Exemple : « Il est communément admis que les machines facilitent la vie de l'homme. »
\end{enumerate}

\begin{exemplebox}[Trois amorces pour un même sujet]
Sujet : « La technique libère-t-elle l'homme ? »
\begin{itemize}
\item Constat : « Le téléphone portable, la moto-taxi, la machine à laver : la technique envahit chaque jour davantage notre existence quotidienne. »
\item Référence : « Selon le mythe grec, Prométhée donna le feu aux hommes pour compenser leur faiblesse naturelle : la technique serait donc née d'un manque. »
\item Lieu commun : « On entend souvent dire que le progrès technique a rendu la vie plus facile et l'homme plus libre. »
\end{itemize}
\end{exemplebox}

\begin{attention}[Piège de l'amorce]
Deux erreurs sont fréquentes. Premièrement, l'amorce \emph{trop éloignée} du sujet (« Depuis la nuit des temps, l'homme... ») qui sonne creux et inquiète le correcteur. Deuxièmement, l'amorce qui \emph{recopie le sujet} : « La technique libère-t-elle l'homme ? C'est la question que nous allons nous poser. » Ce n'est pas une amorce, c'est une répétition. L'amorce doit \emph{amener} le sujet, non le dupliquer.
\end{attention}

\courssub{Deuxième étape : l'analyse des termes du sujet}

Après l'amorce, il faut s'arrêter sur les mots du sujet. En philosophie, on ne peut pas discuter sérieusement d'une question si l'on ne sait pas exactement ce que signifient les mots qui la composent. Un mot mal compris, c'est tout le devoir qui dévie.

\begin{methode}[Analyser les termes du sujet]
\begin{enumerate}
\item \textbf{Repérer les mots-clés} : souligner les deux ou trois termes essentiels du sujet (les noms, le verbe, la relation exprimée).
\item \textbf{Définir chaque mot-clé} par soi-même, avec ses propres mots : donner le sens courant, puis préciser le sens philosophique. Éviter les définitions copiées du dictionnaire.
\item \textbf{Repérer les présupposés} : demander ce que le sujet \emph{admet sans le dire}. Par exemple, « La technique libère-t-elle l'homme ? » présuppose que l'homme pourrait être enchaîné ou aliéné par quelque chose.
\item \textbf{Dégager la tension ou le paradoxe} : montrer que les termes entretiennent un rapport problématique. Ici : la technique est faite par l'homme pour le servir, mais ne risque-t-elle pas de le dominer ? C'est cette tension qui fera naître la problématique.
\end{enumerate}
\end{methode}

Appliquons cette méthode au sujet « La technique libère-t-elle l'homme ? » :
\begin{itemize}
\item \textbf{Technique} : ensemble des procédés, outils et machines inventés par l'homme pour transformer la nature et faciliter son action. Sens philosophique : la technique n'est pas seulement l'outil, c'est une manière d'être au monde.
\item \textbf{Libérer} : rendre libre, délivrer d'une contrainte. Mais libérer de quoi ? Du travail pénible ? De la nature ? Ou de la liberté elle-même ?
\item \textbf{Homme} : être naturel mais aussi être de culture, capable de se dominer lui-même.
\item \textbf{Présupposé} : le sujet suppose que l'homme n'est pas pleinement libre et que la technique pourrait y changer quelque chose.
\item \textbf{Tension} : l'outil est fait pour servir ; or celui qui dépend totalement de ses outils n'est-il pas devenu leur serviteur ?
\end{itemize}

\begin{attention}[Définitions dictionnairiques plaquées]
Recopier « Selon le Petit Robert, la technique est... » est une erreur classique. Le correcteur attend \emph{ta} compréhension du terme, construite par la réflexion, et non une définition plaquée qui n'engage pas ta pensée. Définis avec tes mots, simplement et précisément.
\end{attention}

\courssub{Troisième étape : le dégagement de la problématique}

C'est le cœur de l'introduction, et souvent le moment décisif du devoir.

\begin{definition}[La problématique]
La \cle{problématique} est la question centrale qui naît de l'analyse du sujet et que tout le développement devra traiter. Elle reformule le sujet en creusant sa difficulté : elle met au jour le \emph{problème} que le sujet cache. Elle se formule généralement par une ou deux questions, souvent introduites par « on peut alors se demander », « dès lors » ou « mais ».
\end{definition}

Une bonne problématique possède trois qualités :
\begin{enumerate}
\item elle \textbf{découle} de l'analyse des termes (elle n'arrive pas comme un cheveu sur la soupe) ;
\item elle \textbf{reformule} le sujet en l'approfondissant, souvent sous la forme d'une alternative (ou bien... ou bien...) ;
\item elle \textbf{annonce} l'ensemble du devoir : chaque partie du développement devra y répondre partiellement.
\end{enumerate}

\begin{exemplebox}[Du sujet à la problématique]
Sujet : « Faut-il toujours dire la vérité ? »
\begin{itemize}
\item Analyse : « faut-il » renvoie au devoir moral ; « toujours » pose la question de l'absolu ; « dire la vérité » engage la parole et ses conséquences sur autrui.
\item Problématique : « La vérité est-elle une valeur absolue, ou son dire dépend-il des circonstances ? »
\end{itemize}
On voit que la problématique ne répète pas le sujet : elle en dégage l'enjeu (l'absolu contre les circonstances).
\end{exemplebox}

\begin{attention}[Erreur fréquente : recopier le sujet]
Beaucoup de candidats écrivent : « Le problème qui se pose est le suivant : la technique libère-t-elle l'homme ? » Ce n'est pas une problématique, c'est une recopie. La problématique doit \emph{interroger} le sujet, c'est-à-dire montrer en quoi il est problématique. Mauvaise formulation : répéter le sujet. Bonne formulation : « La technique, faite pour servir l'homme, ne risque-t-elle pas de l'asservir à son tour ? Autrement dit, la libération technique est-elle une libération véritable ? »
\end{attention}

\courssub{Quatrième étape : l'annonce du plan}

L'introduction se termine par l'annonce des grandes parties du développement. C'est une promesse faite au lecteur : voici le chemin que nous allons suivre.

\begin{methode}[Annoncer le plan]
\begin{enumerate}
\item Annoncer les parties dans l'ordre où elles seront traitées (thèse, antithèse, synthèse), en une ou deux phrases.
\item Rester sobre : ne pas détailler les arguments, seulement indiquer la direction.
\item Utiliser des formules de liaison : « Nous verrons d'abord que..., nous montrerons ensuite que..., avant de nous demander enfin si... »
\end{enumerate}
\end{methode}

Exemple : « Nous examinerons d'abord la manière dont la technique semble libérer l'homme des contraintes naturelles ; nous verrons ensuite qu'elle peut créer de nouvelles dépendances ; nous nous demanderons enfin si une liberté véritable ne suppose pas que l'homme demeure maître de ses outils. »

\begin{attention}[Plan annoncé, plan respecté]
Le plan annoncé dans l'introduction doit être \emph{exactement} celui du développement. Annoncer trois parties et n'en traiter que deux, ou les traiter dans un autre ordre, est une faute de construction sanctionnée par le correcteur.
\end{attention}

\courssub{Schéma récapitulatif et exemple guidé}

\begin{popfigure}
\begin{center}
\begin{tikzpicture}[font=\small]
\fill[popblue!25] (0,0) rectangle (3.2,1);
\fill[popgreen!25] (1.1,1) rectangle (4.3,2);
\fill[poporange!25] (2.2,2) rectangle (5.4,3);
\fill[poppurple!25] (3.3,3) rectangle (6.5,4);
\node[anchor=west] at (0.15,0.5) {\textbf{1. Amorce}};
\node[anchor=west] at (1.25,1.5) {\textbf{2. Analyse des termes}};
\node[anchor=west] at (2.35,2.5) {\textbf{3. Problématique}};
\node[anchor=west] at (3.45,3.5) {\textbf{4. Annonce du plan}};
\draw[-{Stealth},thick,popdark] (7,0.5) -- (7,3.8);
\node[rotate=90,anchor=center] at (7.45,2.15) {\footnotesize approfondissement progressif};
\node[anchor=west,text width=5.6cm] at (8,3.5) {\footnotesize partir d'un fait général pour amener le sujet};
\node[anchor=west,text width=5.6cm] at (8,2.5) {\footnotesize définir les mots-clés, repérer présupposés et tension};
\node[anchor=west,text width=5.6cm] at (8,1.5) {\footnotesize formuler la question centrale du devoir};
\node[anchor=west,text width=5.6cm] at (8,0.5) {\footnotesize présenter les grandes parties};
\end{tikzpicture}
\end{center}
\legende{Les quatre étapes de l'introduction : chaque marche s'appuie sur la précédente. On ne peut formuler la problématique qu'après avoir analysé les termes, ni annoncer le plan qu'après avoir posé le problème.}
\end{popfigure}

Construisons maintenant, pas à pas, une introduction complète sur le sujet « La technique libère-t-elle l'homme ? »

\begin{exoresolu}[Rédiger une introduction complète]
\textbf{Énoncé.} Rédiger l'introduction d'une dissertation sur le sujet : « La technique libère-t-elle l'homme ? » en suivant les quatre étapes.
\tcblower
\textbf{Solution détaillée.}

\emph{Étape 1 — Amorce (constat).} « Le téléphone portable, la moto-taxi, la machine à laver : la technique envahit chaque jour davantage notre existence quotidienne, au point qu'il semble impossible de vivre sans elle. »

\emph{Étape 2 — Analyse des termes.} « Par technique, on entend l'ensemble des outils et des procédés inventés par l'homme pour transformer la nature et faciliter son action. Libérer signifie délivrer d'une contrainte : mais de quelle contrainte s'agit-il ? Le sujet présuppose que l'homme n'est pas pleinement libre et que ses propres inventions pourraient l'y aider. Or une tension apparaît : l'outil est fait pour servir, mais celui qui ne peut plus se passer de ses outils n'est-il pas devenu leur dépendant ? »

\emph{Étape 3 — Problématique.} « On peut alors se demander : la technique, faite pour servir l'homme, ne risque-t-elle pas de l'asservir à son tour ? La libération qu'elle promet est-elle une liberté véritable ? »

\emph{Étape 4 — Annonce du plan.} « Nous verrons d'abord que la technique libère l'homme des contraintes naturelles ; nous montrerons ensuite qu'elle engendre de nouvelles dépendances ; nous nous demanderons enfin si la véritable liberté ne consiste pas à demeurer maître de ses outils. »
\end{exoresolu}

\begin{popfigure}
\begin{center}
\begin{tikzpicture}[font=\footnotesize]
\node[draw=popline,fill=popcream,rounded corners,text width=9.2cm,align=justify] (intro) at (0,0) {
Le téléphone portable envahit chaque jour notre existence, au point qu'il semble impossible de vivre sans lui. Par \emph{technique}, on entend l'ensemble des outils inventés pour transformer la nature ; \emph{libérer} signifie délivrer d'une contrainte. Or l'outil est fait pour servir : celui qui ne peut plus s'en passer n'est-il pas devenu son dépendant ? La technique, faite pour servir l'homme, ne risque-t-elle pas de l'asservir ? Nous verrons d'abord qu'elle libère des contraintes naturelles, puis qu'elle crée de nouvelles dépendances, avant de demander si l'homme peut demeurer maître de ses outils.};
\node[draw=popblue,fill=popblue!15,rounded corners] (a) at (6.8,1.8) {amorce};
\node[draw=popgreen,fill=popgreen!15,rounded corners] (b) at (6.8,0.6) {analyse};
\node[draw=poporange,fill=poporange!15,rounded corners] (c) at (6.8,-0.6) {problématique};
\node[draw=poppurple,fill=poppurple!15,rounded corners] (d) at (6.8,-1.8) {plan};
\draw[-{Stealth},popblue] (a.west) -- (4.7,1.5);
\draw[-{Stealth},popgreen] (b.west) -- (4.7,0.4);
\draw[-{Stealth},poporange] (c.west) -- (4.7,-0.5);
\draw[-{Stealth},poppurple] (d.west) -- (4.7,-1.4);
\end{tikzpicture}
\end{center}
\legende{Introduction annotée : chaque flèche renvoie à l'une des quatre étapes. Remarquer les enchaînements logiques (« or », « nous verrons ») qui assurent la fluidité du texte.}
\end{popfigure}

\courssub{Erreurs à éviter et bilan}

Récapitulons les trois fautes qui ruinent une introduction :
\begin{itemize}
\item \textbf{Le hors-sujet} : une amorce trop générale ou une analyse qui s'écarte des mots du sujet fait dériver tout le devoir. Reste toujours au plus près des termes donnés.
\item \textbf{Les définitions plaquées} : les définitions de dictionnaire, copiées sans réflexion, montrent que le candidat ne pense pas par lui-même.
\item \textbf{La problématique absente} : une introduction sans question centrale est une introduction sans problème ; or sans problème, il n'y a rien à discuter, et le développement devient un simple exposé.
\end{itemize}

\begin{aretenir}[L'essentiel]
L'introduction comporte quatre étapes enchaînées : l'\cle{amorce} (fait, constat ou référence qui amène le sujet), l'\cle{analyse des termes} (définition des mots-clés, présupposés, tension), la \cle{problématique} (question centrale qui reformule et approfondit le sujet) et l'\cle{annonce du plan} (présentation sobre des parties). Chaque étape prépare la suivante : on ne peut poser un bon problème qu'après avoir bien compris les mots du sujet.
\end{aretenir}

\begin{savaistu}[L'introduction, première impression du correcteur]
Au Probatoire et au Baccalauréat technique, les correcteurs lisent l'introduction en premier. Une problématique claire et bien formulée conditionne souvent la moitié de l'appréciation globale du devoir : elle prouve que le candidat a compris le sujet et sait où il va. Inversement, une introduction brouillonne décourage le lecteur avant même le développement. Soigner l'introduction, c'est donc soigner sa copie tout entière.
\end{savaistu}

\begin{exercice}[À toi de jouer]
Pour chacun des sujets suivants, rédige une amorce, analyse les termes et formule une problématique : 1) « Faut-il toujours dire la vérité ? » ; 2) « L'État doit-il tout décider pour les citoyens ? » ; 3) « Le travail est-il une contrainte ou une liberté ? »
\end{exercice}

\begin{corrige}[Exercice (sujet 3)]
Éléments de correction pour le sujet 3. \emph{Amorce} : constat — « À Douala comme à Yaoundé, des millions de travailleurs se lèvent chaque matin pour gagner leur pain : le travail occupe l'essentiel de la vie humaine. » \emph{Analyse} : travail = activité par laquelle l'homme transforme la nature et produit ; contrainte = obligation subie ; liberté = pouvoir d'agir selon sa volonté. Présupposé : le travail pourrait n'être qu'une nécessité imposée. Tension : le travail fatigue et oblige, mais c'est aussi par lui que l'homme se réalise et s'affranchit de la nature. \emph{Problématique} : « Le travail est-il seulement une nécessité subie, ou est-il aussi le moyen par lequel l'homme se libère et s'accomplit ? »
\end{corrige}

L'introduction ainsi construite ouvre naturellement sur le développement : c'est dans la première partie, la thèse, que nous donnerons sa première réponse à la problématique — réponse que la section suivante nous apprendra à construire et à argumenter.

\coursec{La thèse}

Après avoir maîtrisé l'art de l'introduction — cette porte d'entrée qui pose le problème, l'annonce et le plan —, nous entrons maintenant dans le cœur même de la dissertation : le développement. La première grande partie du développement est la \textbf{thèse}. C'est elle qui donne la première réponse argumentée à la problématique. Rappelons-le : une dissertation n'est pas une suite d'opinions, c'est un \emph{parcours de la raison}. La thèse en est la première étape structurée.

\courssub{Définition et rôle de la thèse dans l'économie de la dissertation}

\begin{definition}[La thèse]
La \cle{thèse} est la première position argumentée que le candidat soutient dans le développement de sa dissertation. Elle répond une première fois à la problématique, généralement de manière affirmative ou selon un sens précis (ex. : « Oui, l'État est nécessaire à la liberté » ou « La liberté ne s'exerce pleinement qu'à travers l'État »). Elle n'est pas une simple opinion : elle doit être \textbf{justifiée}, \textbf{structurée} et \textbf{illustrée}.
\end{definition}

Dans la dialectique classique (thèse — antithèse — synthèse), la thèse joue le rôle de \textbf{position initiale}. Elle prend au sérieux un aspect du problème et tente de le démontrer. Au Probatoire comme au Baccalauréat technique, l'examinateur attend que cette partie soit \emph{rigoureuse} : une thèse mal construite, c'est tout l'édifice qui vacille.

\begin{aretenir}[Place de la thèse]
\begin{itemize}
    \item Elle constitue la \textbf{première grande partie} du développement (souvent notée I dans le plan détaillé).
    \item Elle comprend généralement \textbf{2 à 3 arguments distincts}, chacun développé dans un paragraphe séparé.
    \item Elle se termine par une \textbf{transition} qui en montre les limites et appelle l'antithèse.
\end{itemize}
\end{aretenir}

\courssub{Architecture d'un paragraphe de thèse : la méthode du « paragraphe unique »}

Un paragraphe de dissertation n'est pas un bloc de texte compact. Il obéit à une \textbf{structure logique immuable} que tout candidat doit maîtriser comme une seconde nature.

\begin{methode}[Construire un paragraphe argumentatif — 5 étapes]
\begin{enumerate}
    \item \textbf{Idée directrice (ou phrase-topic)} : Annoncer en une phrase claire l'argument principal du paragraphe. C'est la \cle{thèse partielle} que ce paragraphe va défendre.
    \item \textbf{Argument(s)} : Expliquer, justifier, démontrer l'idée directrice. Mobiliser la raison (logique), l'autorité (philosophes) ou l'expérience (faits vécus, historiques, sociaux).
    \item \textbf{Exemple(s) concret(s)} : Illustrer l'argument par une situation précise, datée, localisée si possible (Cameroun, actualité, histoire). L'exemple \emph{n'est pas une preuve}, il \emph{rend sensible} l'argument.
    \item \textbf{Analyse de l'exemple} : Montrer \emph{en quoi} l'exemple confirme l'argument. C'est l'étape la plus souvent oubliée et la plus valorisée à l'examen.
    \item \textbf{Mini-transition (ou phrase de liaison)} : Refermer le paragraphe en annonçant l'argument suivant ou en préparant la transition vers l'antithèse.
\end{enumerate}
\end{methode}

\begin{popfigure}
\begin{tikzpicture}[
    block/.style={draw=popblue, thick, rounded corners, fill=popblue!10, minimum height=1.2cm, text width=3.2cm, align=center, font=\small},
    arrow/.style={-Stealth, thick, popdark},
    node distance=0.8cm and 1.5cm
]
\node[block, align=center] (idee){\textbf{1. Idée directrice}\\(phrase-topic)\\\textit{« L'État garantit la sécurité sans laquelle pas de liberté »}};
\node[block, below=of idee, align=center] (arg){\textbf{2. Argument(s)}\\\textit{Logique : contrat social}\\\textit{Autorité : Hobbes, Locke}\\\textit{Expérience : État de nature}};
\node[block, below=of arg, align=center] (ex){\textbf{3. Exemple concret}\\\textit{Cameroun : forces de l'ordre,}\\\textit{justice, administration territoriale}\\\textit{assurent l'ordre public}};
\node[block, below=of ex, align=center] (ana){\textbf{4. Analyse}\\\textit{Sans police ni justice,}\\\textit{la loi du plus fort règne :}\\\textit{liberté = 0 pour le faible}};
\node[block, below=of ana, align=center] (trans){\textbf{5. Transition}\\\textit{Mais l'État peut aussi}\\\textit{devenir oppresseur...}};
\draw[arrow] (idee) -- (arg);
\draw[arrow] (arg) -- (ex);
\draw[arrow] (ex) -- (ana);
\draw[arrow] (ana) -- (trans);
% Étiquettes latérales
\node[left=0.3cm of idee, text width=2.5cm, align=right, font=\small\itshape, popmuted] {Annoncer};
\node[left=0.3cm of arg, text width=2.5cm, align=right, font=\small\itshape, popmuted] {Prouver};
\node[left=0.3cm of ex, text width=2.5cm, align=right, font=\small\itshape, popmuted] {Illustrer};
\node[left=0.3cm of ana, text width=2.5cm, align=right, font=\small\itshape, popmuted] {Analyser};
\node[left=0.3cm of trans, text width=2.5cm, align=right, font=\small\itshape, popmuted] {Relier};
\end{tikzpicture}
\legende{Architecture type d'un paragraphe argumentatif de thèse : chaque étage repose sur le précédent. L'analyse (4) est l'étage critique qui transforme l'exemple en preuve.}
\end{popfigure}

\courssub{Les trois types d'arguments : varier les registres pour convaincre}

Une thèse solide ne repose pas sur un seul type d'argument. Le candidat avisé alterne les \textbf{trois registres} pour montrer l'étendue de sa culture philosophique et de sa capacité de raisonnement.

\begin{popfigure}
\begin{tikzpicture}[
    grow'=right, level distance=5.5cm,
    level 1/.style={sibling distance=3.5cm},
    level 2/.style={sibling distance=3.2cm},
    every node/.style={draw=popdark, thick, rounded corners, fill=popcream, minimum height=1cm, text width=3cm, align=center, font=\small},
    edge from parent/.style={draw=popblue, thick, -Stealth},
    edge from parent path={(\tikzparentnode.east) -- ++(1cm,0) |- (\tikzchildnode.west)}
]
\node[align=center]{\textbf{Types d'arguments}\\(registres de preuve)}
    child { node[align=center]{\textbf{1. Argument d'autorité}\\\textit{(Référence philosophique)}}
        child { node[align=center]{\textbf{Exemple :}\\\textit{Kant : « L'État de droit}\\\textit{transforme la liberté}\\\textit{sauvage en liberté civile »}} }
    }
    child { node[align=center]{\textbf{2. Argument logique}\\\textit{(Raisonnement pur)}}
        child { node[align=center]{\textbf{Exemple :}\\\textit{Si liberté = absence de contrainte,}\\\textit{alors État (contrainte) $\sim$ liberté.}\\\textit{Mais si liberté = autonomie,}\\\textit{État (loi) = condition.}} }
    }
    child { node[align=center]{\textbf{3. Argument par l'expérience}\\\textit{(Fait observable, histoire, société)}}
        child { node[align=center]{\textbf{Exemple :}\\\textit{Somalie (1991--2012) :}\\\textit{effondrement de l'État}\\\textit{$\to$ insécurité, milices,}\\\textit{liberté nulle pour les civils}} }
    };
\end{tikzpicture}
\legende{Les trois registres argumentatifs et un exemple d'application au sujet « L'État est-il nécessaire à la liberté ? ». Une bonne thèse en mobilise au moins deux.}
\end{popfigure}

\begin{definition}[Argument d'autorité]
Citer un philosophe \textbf{pertinent} et \textbf{expliqué}. Dire « Kant dit que... » ne suffit pas : il faut restituer \emph{la pensée} de l'auteur et montrer \emph{pourquoi} elle éclaire le problème. \textbf{Attention} : l'autorité n'est pas une preuve absolue, elle est un \emph{appui}.
\end{definition}

\begin{definition}[Argument logique]
Un raisonnement \emph{à priori}, indépendant de l'expérience : déduction, analogie, réduction à l'absurde, mise en évidence d'une contradiction. Ex. : « Si l'État n'est pas nécessaire, alors l'anarchie garantirait la liberté. Or l'anarchie est la loi du plus fort. Donc... »
\end{definition}

\begin{definition}[Argument par l'expérience]
S'appuyer sur des \textbf{faits réels} : histoire (Révolution française, décolonisation), actualité (crise anglophone au Cameroun, guerre en Ukraine), sociologie (enquêtes, statistiques), vie quotidienne. C'est l'argument le plus accessible aux élèves de l'enseignement technique.
\end{definition}

\begin{savaistu}[Les grands noms qui renforcent une thèse]
\begin{itemize}
    \item \textbf{Platon} : l'État juste, gouverné par les philosophes-rois (\emph{République}).
    \item \textbf{Hobbes} : l'État (Léviathan) comme rempart contre la « guerre de tous contre tous ».
    \item \textbf{Locke} : l'État protège la vie, la liberté, la propriété (droits naturels).
    \item \textbf{Rousseau} : l'État légitime = volonté générale ; « obéir à la loi qu'on s'est prescrite, c'est la liberté ».
    \item \textbf{Kant} : l'État de droit, condition de la liberté extérieure (coexistence selon des lois universelles).
    \item \textbf{Senghor} : l'État comme instrument de la « civilisation de l'universel », au service de la personne humaine.
    \item \textbf{Arendt} : la liberté comme action politique \emph{dans} l'espace public (polis/État).
\end{itemize}
\textbf{Règle d'or} : ne citez que ce que vous comprenez et pouvez expliquer. Une citation mal comprise décrédibilise la copie.
\end{savaistu}

\courssub{Le choix des exemples : ancrer la philosophie dans le réel camerounais}

\begin{attention}[Piège n°1 : l'accumulation stérile]
L'erreur la plus fréquente : aligner trois exemples (Cameroun, France, Rome antique) \textbf{sans les analyser}. L'examinateur note : « Exemples non exploités ». \textbf{Un seul exemple bien analysé vaut mieux que trois survolés.}
\end{attention}

\begin{propriete}[Critères d'un bon exemple]
\begin{enumerate}
    \item \textbf{Concret} : daté, localisé, nommé (pas « dans certains pays » mais « au Cameroun, lors de la crise de 2016--2019 dans le Nord-Ouest et le Sud-Ouest »).
    \item \textbf{Pertinent} : il illustre \emph{précisément} l'argument en cours, pas un autre.
    \item \textbf{Analysé} : le candidat explicite le lien : « Cet exemple montre que... », « On voit ici que... », « Cela confirme l'idée selon laquelle... ».
    \item \textbf{Varié} : alterner exemple historique, exemple d'actualité, exemple de la vie locale (quartier, école, marché, administration).
\end{enumerate}
\end{propriete}

\begin{exemplebox}[Exemples camerounais pour « L'État est-il nécessaire à la liberté ? »]
\begin{itemize}
    \item \textbf{Positif (thèse)} : La carte d'identité biométrique, le passeport, le permis de conduire — documents d'État qui \emph{permettent} la circulation, l'identification, l'accès aux services bancaires, donc l'exercice de libertés concrètes.
    \item \textbf{Positif} : L'école publique, l'université d'État — l'État finance l'instruction qui \emph{rend capable} de liberté (autonomie intellectuelle).
    \item \textbf{Nuancé (transition)} : Les contrôles routiers abusifs, les arrestations arbitraires, la loi anti-terrorisme de 2014 utilisée contre des opposants — l'État \emph{restreint} la liberté au nom de la sécurité.
\end{itemize}
\end{exemplebox}

\courssub{La transition vers l'antithèse : l'art de la pivot}

La thèse ne s'arrête pas brutalement. Elle se termine par une \textbf{transition} qui remplit deux fonctions :
\begin{enumerate}
    \item \textbf{Bilan} : résumer en une phrase ce que la thèse a établi (ex. : « Ainsi, l'État apparaît comme la condition institutionnelle sans laquelle la liberté reste une abstraction vide »).
    \item \textbf{Ouverture critique} : montrer \emph{les limites} de cette thèse, \emph{l'insuffisance} de cette réponse, pour justifier le passage à l'antithèse (ex. : « Pourtant, l'histoire montre que l'État peut devenir le premier ennemi de la liberté. C'est ce que l'antithèse examinera »).
\end{enumerate}

\begin{methode}[Rédiger la transition thèse $\to$ antithèse]
\begin{enumerate}
    \item Reprendre le mot-clé de la thèse (ex. : « condition », « garantie », « cadre »).
    \item Introduire un opérateur de concession : \textbf{« Toutefois »}, \textbf{« Cependant »}, \textbf{« Mais »}, \textbf{« Certes... toutefois »}.
    \item Pointer la limite : excès de l'État, dérive totalitaire, confusion ordre/liberté, négligence des libertés sociales, etc.
    \item Annoncer l'antithèse : « C'est pourquoi il faut se demander si l'État ne menace pas la liberté qu'il prétend protéger. »
\end{enumerate}
\end{methode}

\courssub{Exemple guidé : rédaction complète d'une thèse sur « L'État est-il nécessaire à la liberté ? »}

Voici une thèse modèle, structurée en \textbf{trois arguments} (donc trois paragraphes), suivie de sa transition. Chaque paragraphe respecte la méthode en 5 étapes.

\begin{exoresolu}[Thèse complète : « L'État est-il nécessaire à la liberté ? »]
\textbf{Sujet :} L'État est-il nécessaire à la liberté ? \\
\textbf{Problématique :} L'État, par sa contrainte légale, semble limiter la liberté ; mais sans lui, la liberté n'est-elle pas livrée à l'arbitraire ? \\
\textbf{Thèse (I) :} \textbf{Oui, l'État est la condition institutionnelle indispensable à l'exercice effectif de la liberté.}
\tcblower
\textbf{Argument 1 : L'État instaure l'ordre juridique sans lequel la liberté n'est que la loi du plus fort.}
\begin{itemize}
    \item \textit{Idée directrice :} La liberté suppose un cadre de règles communes, garanties par une puissance publique impartiale.
    \item \textit{Argument logique :} Dans l'état de nature (Hobbes, \emph{Léviathan}, 1651), « l'homme est un loup pour l'homme » : chacun juge sa propre cause, la force prime le droit. La liberté y est illusoire car précaire. L'État, par le contrat social, substitue la \emph{loi} à la \emph{force}.
    \item \textit{Argument d'autorité :} Pour Locke (\emph{Second traité du gouvernement civil}, 1690), l'État naît pour protéger les « droits naturels » (vie, liberté, propriété) que l'individu ne peut garantir seul.
    \item \textit{Exemple concret :} Au Cameroun, le Code pénal et le Code de procédure pénale, appliqués par les tribunaux (TPI, Cour d'appel, Cour suprême), permettent à un citoyen spolié de ses terres à Douala ou à Bamenda de saisir la justice au lieu de recourir à l'autodéfense violente.
    \item \textit{Analyse :} Cet exemple montre que l'État, par son monopole de la violence légitime (Weber), transforme un conflit de force en procédure de droit : le faible a une chance face au fort. C'est la \emph{condition de possibilité} de la liberté civile.
    \item \textit{Mini-transition :} Mais l'ordre juridique ne suffit pas : il faut encore que l'État fournisse les moyens concrets d'exercer sa liberté.
\end{itemize}

\medskip
\textbf{Argument 2 : L'État crée les conditions matérielles et éducatives de l'autonomie (liberté positive).}
\begin{itemize}
    \item \textit{Idée directrice :} La liberté ne se réduit pas à l'absence d'entraves (liberté négative) ; elle requiert des capacités réelles (liberté positive), que l'État doit permettre.
    \item \textit{Argument logique :} Un analphabète, un malade sans soins, un chômeur sans qualification sont \emph{formellement} libres mais \emph{réellement} incapables de choisir leur vie. L'État, par la redistribution (impôts $\to$ services publics), égalise les chances.
    \item \textit{Argument d'autorité :} Pour Rousseau (\emph{Du contrat social}, 1762), « l'obéissance à la loi qu'on s'est prescrite, c'est la liberté ». L'État-providence (éducation, santé, formation professionnelle) rend les citoyens \emph{capables} d'autonomie.
    \item \textit{Exemple concret :} Le programme « Éducation de base » et les lycées techniques publics au Cameroun (ex. : Lycée technique de Douala-Bassa, Lycée polyvalent de Bonabéri) forment gratuitement des milliers de jeunes aux métiers de l'industrie, du BTP, de l'informatique. Sans cet État éducateur, ces jeunes — souvent issus de milieux modestes — n'auraient pas les compétences pour accéder à l'emploi salarié et à l'indépendance économique.
    \item \textit{Analyse :} Ici, l'État ne se contente pas de « ne pas empêcher » : il \emph{donne les moyens} d'être libre. La liberté technique et économique est une \emph{construction étatique}.
    \item \textit{Mini-transition :} Cependant, cette puissance étatique, si elle n'est pas contrôlée, peut basculer dans l'oppression.
\end{itemize}

\medskip
\textbf{Argument 3 : L'État garantit l'espace public où s'exerce la liberté politique (citoyenneté).}
\begin{itemize}
    \item \textit{Idée directrice :} La liberté suprême, selon les Anciens et Arendt, est la participation à la chose publique ; l'État en est le cadre institutionnel.
    \item \textit{Argument d'autorité :} Hannah Arendt (\emph{La crise de la culture}, 1961) distingue la \emph{liberté} (action concertée dans l'espace public) de la \emph{libération} (délivrance des chaînes). L'État de droit organise cet espace : élections, parlement, presse, associations.
    \item \textit{Argument par l'expérience :} Les pays sans État fonctionnel (Somalie 1991--2012, Libye post-2011) voient disparaître toute liberté politique : pas d'élections, pas de débat public, justice des milices.
    \item \textit{Exemple concret :} Au Cameroun, les élections municipales, législatives et présidentielle de 2020, malgré leurs imperfections, ont permis à des partis d'opposition (MRC, SDF, UDC, etc.) de présenter des candidats, de faire campagne, d'obtenir des sièges (ex. : 139 conseillers municipaux MRC en 2020). L'État électoral a \emph{rendu possible} cette expression politique.
    \item \textit{Analyse :} Certes, le jeu politique camerounais reste perfectible (inégalités d'accès aux médias, contentieux post-électoraux). Mais \emph{sans} Commission électorale (ELECAM), \emph{sans} Code électoral, \emph{sans} Conseil constitutionnel — tous organes d'État —, \emph{aucune} compétition politique pacifique n'aurait lieu. L'État est le \emph{théâtre} de la liberté politique.
    \item \textit{Mini-transition :} Cette thèse établie, il faut maintenant en tester la solidité face aux contre-exemples historiques et théoriques.
\end{itemize}

\medskip
\textbf{Transition vers l'antithèse (II) :}
\begin{quote}
\textit{« Au terme de cette thèse, l'État apparaît comme le grand architecte de la liberté : il la fonde juridiquement, la rend matériellement possible, l'institue politiquement. \textbf{Toutefois}, l'histoire et la philosophie nous avertissent : le même État qui protège peut aussi asservir. La puissance publique, une fois instituée, tend à s'autonomiser, à confondre ordre public et conservation du pouvoir, à transformer la protection en tutelle. \textbf{C'est pourquoi il faut maintenant examiner la thèse inverse : l'État n'est-il pas, par sa nature même, une menace pour la liberté ?} »}
\end{quote}
\end{exoresolu}

\begin{aretenir}[Bilan chiffré pour le Probatoire / Bac Tech]
\begin{itemize}
    \item Une thèse réussie = \textbf{2 à 3 paragraphes} (arguments distincts).
    \item Chaque paragraphe = \textbf{1 idée directrice + 1 à 2 arguments + 1 exemple analysé + 1 transition}.
    \item Total arguments dans la dissertation = \textbf{6 à 9} (2--3 par partie $\times$ 3 parties).
    \item Temps conseillé à l'examen : \textbf{35--45 minutes} pour rédiger la thèse complète (sur 4h).
\end{itemize}
\end{aretenir}

\courssub{Exercice d'entraînement : à vous de jouer}

\begin{exercice}[Thèse à construire]
\textbf{Sujet :} « La technique libère-t-elle l'homme ? » \\
\textbf{Consigne :} Rédigez \textbf{un seul paragraphe} de thèse (argument 1) en respectant scrupuleusement la méthode en 5 étapes. Choisissez l'angle : « La technique libère l'homme de la pénibilité et de la nécessité naturelle. »
\begin{enumerate}
    \item Écrivez l'idée directrice.
    \item Développez l'argument (logique + autorité : Marx, \emph{Le Capital}, ou Ellul, \emph{La technique ou l'enjeu du siècle}).
    \item Donnez un exemple concret camerounais (ex. : pompe à eau solaire dans un village du Nord, téléphone mobile pour paysan, machine à décortiquer le riz à Ndop).
    \item Analysez l'exemple en 2--3 lignes.
    \item Terminez par une mini-transition vers l'argument 2.
\end{enumerate}
\end{exercice}

\begin{corrige}[Exercice 1]
\textbf{Idée directrice :} La technique, en automatisant les tâches pénibles et répétitives, affranchit l'homme de la contrainte naturelle immédiate et lui donne du temps pour des activités choisies. \\
\textbf{Argument :} Marx montre que le développement des forces productives (machines, énergie) réduit le « royaume de la nécessité » (travail imposé par la survie) pour étendre le « royaume de la liberté » (création, culture, loisirs). La technique est le levier historique de cette émancipation. \\
\textbf{Exemple :} Dans la plaine de Ndop (Région du Nord-Ouest), l'introduction de décortiqueuses motorisées pour le riz a libéré les femmes de la tâche épuisante du pilage manuel (3 à 4 heures/jour). Elles consacrent ce temps gagné à l'alphabétisation, au petit commerce ou à l'encadrement scolaire des enfants. \\
\textbf{Analyse :} Ici, la machine ne remplace pas l'humain pour l'exclure, elle le \emph{redéploie} vers des activités à plus forte valeur humaine. La technique a été \emph{libératrice} parce qu'elle a ciblé une contrainte biologique (fatigue, temps) sans aliéner le travailleur (gestion coopérative de la machine). \\
\textbf{Mini-transition :} Mais cette libération par la technique suppose que l'homme en reste le maître ; or, la technique peut aussi s'imposer à lui comme un système autonome...
\end{corrige}

\begin{attention}[Derniers pièges à éviter dans la thèse]
\begin{itemize}
    \item \textbf{Confondre thèse et introduction} : ne ré-annoncez pas le plan dans le premier paragraphe de thèse.
    \item \textbf{Écrire « Je pense que... »} : la dissertation est impersonnelle (on utilise « Il apparaît que... », « Force est de constater que... »).
    \item \textbf{Mélanger thèse et antithèse} : un paragraphe = un sens unique. Les nuances viennent en synthèse.
    \item \textbf{Oublier la transition finale} : sans elle, la dissertation semble s'arrêter, ne pas \emph{progresser}.
\end{itemize}
\end{attention}

La thèse est maintenant solidement construite. Nous verrons dans la section suivante comment l'\textbf{antithèse} vient la challenger sans la détruire, pour préparer la \textbf{synthèse} qui en tirera la vérité philosophique.

\coursec{L'antithèse}

Dans la section précédente, nous avons appris à construire la \cle{thèse} : une première réponse au problème, soutenue par des arguments et des exemples. Mais une dissertation philosophique qui s'arrêterait à la thèse ressemblerait à un procès où l'on n'entendrait que l'accusation, jamais la défense. Or le juge — et le correcteur du Probatoire ou du Baccalauréat technique — n'accepte de rendre un verdict qu'après avoir écouté les deux parties. C'est précisément le rôle de l'\cle{antithèse} : donner la parole à la position adverse.

\courssub{Définition de l'antithèse}

Commençons par une intuition. Deux amis, à Yaoundé, discutent du téléphone portable. Le premier affirme : « Le portable rapproche les gens, car il permet d'appeler sa famille restée au village. » Le second réplique : « Au contraire, le portable éloigne les gens : à table, chacun regarde son écran et plus personne ne se parle. » Le premier a énoncé une thèse ; le second vient de formuler une \cle{antithèse}. Remarquons que les deux parlent du même objet, avec des arguments sérieux, mais aboutissent à des conclusions opposées.

\begin{definition}[L'antithèse]
L'\cle{antithèse} est la position \textbf{contraire} à la thèse : elle conteste la réponse d'abord donnée au problème et montre que ce problème admet une \textbf{autre réponse légitime}, elle aussi fondée sur des arguments. Elle constitue la deuxième partie du développement d'une dissertation dialectique (plan thèse -- antithèse -- synthèse).
\end{definition}

L'antithèse n'est donc pas un simple « non » jeté à la figure de la thèse. C'est une \textbf{position construite}, avec sa propre logique, ses propres arguments et ses propres exemples. Elle révèle la \textbf{complexité du problème} : si la question posée était simple, elle n'aurait pas fait l'objet d'un sujet de dissertation.

\begin{aretenir}[L'essentiel]
Thèse et antithèse répondent à la \textbf{même question} mais de manière \textbf{opposée}. L'antithèse doit être défendue avec la \textbf{même rigueur} que la thèse : arguments, exemples, références d'auteurs. Une antithèse faible ou ridicule détruit toute la dissertation.
\end{aretenir}

\courssub{Le rôle dialectique de l'antithèse}

Pourquoi la philosophie exige-t-elle d'examiner le point de vue adverse ? Parce que la philosophie refuse le \cle{dogmatisme}, c'est-à-dire l'attitude de celui qui affirme sa vérité sans jamais la soumettre à l'épreuve de la contradiction.

\begin{propriete}[Rôle dialectique de l'antithèse]
L'antithèse remplit trois fonctions dans la dissertation :
\begin{enumerate}
\item \textbf{Éviter le dogmatisme} : en examinant sérieusement le point de vue opposé, le candidat prouve qu'il ne confond pas sa propre opinion avec la vérité.
\item \textbf{Approfondir le problème} : la confrontation des deux positions fait apparaître des aspects que chacune, prise isolément, avait négligés.
\item \textbf{Préparer la synthèse} : c'est parce que la thèse et l'antithèse se contredisent que naît le besoin d'une position supérieure capable de les dépasser.
\end{enumerate}
\end{propriete}

\begin{savaistu}[D'où vient la méthode thèse -- antithèse -- synthèse ?]
La démarche qui va de la thèse à l'antithèse puis à la synthèse est traditionnellement associée à la \cle{dialectique} du philosophe allemand \textbf{Hegel} (1770--1831), pour qui la vérité se construit par la confrontation des contraires. Cette méthode a été reprise et transformée par \textbf{Karl Marx} (1818--1883), qui l'a appliquée à l'histoire et à la société : selon lui, les contradictions sociales (par exemple entre patrons et ouvriers) font progresser l'histoire. En philosophie, discuter sérieusement l'opinion contraire est donc une tradition vieille de plus de deux siècles : Socrate déjà, dans les rues d'Athènes, faisait accoucher les esprits en opposant leurs certitudes à leurs propres contradictions.
\end{savaistu}

\begin{popfigure}
\begin{center}
\begin{tikzpicture}[scale=1.0]
% Poteau central
\draw[line width=2pt,popdark] (0,0) -- (0,3.2);
\draw[line width=2pt,popdark] (-1.1,0) -- (1.1,0);
\draw[line width=2pt,popdark] (0,0) -- (-1.1,0) (0,0) -- (1.1,0);
% Flèche du haut
\draw[-{Stealth},line width=2pt,popdark] (0,3.2) -- (0,3.9);
\node[above,popdark,font=\bfseries] at (0,3.9) {SYNTHÈSE};
% Flèche vers synthèse depuis les plateaux
\draw[-{Stealth},line width=1.2pt,popmuted,dashed] (-3.4,1.6) .. controls (-2.5,3.4) .. (-0.3,3.8);
\draw[-{Stealth},line width=1.2pt,popmuted,dashed] (3.4,1.6) .. controls (2.5,3.4) .. (0.3,3.8);
% Flèche de confrontation
\draw[{Stealth}-{Stealth},line width=1.5pt,popink] (-2.2,2.35) -- (2.2,2.35);
\node[above,popink,font=\small\itshape] at (0,2.4) {confrontation};
% Plateau gauche : THÈSE
\draw[line width=1.5pt,popblue] (-3.4,2.3) -- (-4.4,1.3);
\draw[line width=1.5pt,popblue] (-3.4,2.3) -- (-2.4,1.3);
\fill[popblueL,draw=popblue,line width=1.5pt] (-4.6,1.3) arc (180:360:1.2) -- cycle;
\node[popblue,font=\bfseries] at (-3.4,0.75) {THÈSE};
\node[popdark,font=\small,align=center] at (-3.4,-0.5) {« Le travail\\ aliène l'homme »};
% Plateau droit : ANTITHÈSE
\draw[line width=1.5pt,poporange] (3.4,2.3) -- (2.4,1.3);
\draw[line width=1.5pt,poporange] (3.4,2.3) -- (4.4,1.3);
\fill[poporangeL,draw=poporange,line width=1.5pt] (2.2,1.3) arc (180:360:1.2) -- cycle;
\node[poporange,font=\bfseries] at (3.4,0.75) {ANTITHÈSE};
\node[popdark,font=\small,align=center] at (3.4,-0.5) {« Le travail\\ humanise l'homme »};
\end{tikzpicture}
\end{center}
\legende{La balance dialectique : la thèse et l'antithèse se font face avec un poids d'arguments comparable ; leur confrontation équilibrée élève la réflexion vers la synthèse.}
\end{popfigure}

\courssub{Comment construire l'antithèse}

L'antithèse ne tombe pas du ciel : elle se construit méthodiquement, en quatre étapes.

\begin{methode}[Construire une antithèse solide]
\begin{enumerate}
\item \textbf{Reprendre les faiblesses signalées dans la transition.} À la fin de la thèse, on a formulé une objection ou une limite (« Cependant, cette position se heurte à une difficulté... »). L'antithèse développe précisément cette difficulté.
\item \textbf{Retourner la question.} Si la thèse répond « oui », l'antithèse répond « non », en montrant que ce « non » peut lui aussi se justifier.
\item \textbf{Argumenter avec la même rigueur que pour la thèse} : deux ou trois arguments, chacun illustré par un exemple concret et, si possible, appuyé sur un auteur.
\item \textbf{Utiliser des connecteurs d'opposition} pour marquer clairement le renversement : \emph{cependant}, \emph{néanmoins}, \emph{mais}, \emph{en revanche}, \emph{au contraire}, \emph{pourtant}, \emph{toutefois}.
\end{enumerate}
\end{methode}

\begin{attention}[Le piège de l'homme de paille]
L'erreur la plus fréquente consiste à présenter une antithèse \textbf{faible ou caricaturale}, facile à réfuter, pour donner l'impression que la thèse a gagné d'avance. Ce procédé porte un nom : l'\cle{homme de paille}. Par exemple, face à la thèse « les réseaux sociaux unissent la jeunesse », écrire « certains prétendent que les réseaux sociaux divisent parce qu'ils sont tous stupides » est une caricature, pas un argument. Le correcteur le remarque immédiatement et sanctionne. \textbf{Règle d'or} : défends l'antithèse comme si tu y croyais sincèrement, avec les meilleurs arguments possibles.
\end{attention}

\courssub{L'articulation logique entre thèse et antithèse}

Thèse et antithèse doivent s'opposer \textbf{réellement}, sur le \textbf{même terrain}, et non artificiellement. Deux défauts sont à éviter :

\begin{itemize}
\item \textbf{Le faux débat} : les deux parties disent la même chose avec des mots différents. Il n'y a alors aucune contradiction, donc aucun progrès de la pensée.
\item \textbf{Le dialogue de sourds} : les deux parties parlent de sujets différents (l'une du travail comme salaire, l'autre du travail comme création artistique) sans jamais se répondre. L'opposition doit porter sur le \textbf{même point précis} du problème.
\end{itemize}

\begin{exemplebox}[Thèse et antithèse sur le travail]
Sujet : \emph{Le travail libère-t-il l'homme ?}
\begin{itemize}
\item \textbf{Thèse} : le travail \textbf{aliène} l'homme. Selon \textbf{Marx}, dans le système capitaliste, l'ouvrier ne possède ni l'outil ni le produit de son travail : il se sent étranger à ce qu'il fabrique. Exemple : l'employé d'une usine de Douala qui répète le même geste huit heures par jour sans jamais voir le produit fini.
\item \textbf{Antithèse} : le travail \textbf{humanise et libère} l'homme. Pour \textbf{Hegel}, en transformant la nature par son travail, l'homme se transforme lui-même, prend conscience de ses capacités et conquiert sa dignité. Exemple : le menuisier de Bafoussam qui fabrique un meuble de ses mains éprouve la fierté de voir son œuvre accomplie et reconnue par ses clients.
\end{itemize}
L'opposition est réelle : les deux positions parlent du même objet (le rapport de l'homme à son travail) et aboutissent à des conclusions contradictoires (aliénation contre libération).
\end{exemplebox}

Voici comment se présente, sous forme de tableau, la confrontation complète sur un sujet type Baccalauréat technique :

\begin{center}
\begin{tabularx}{\linewidth}{|l|X|X|}
\hline
\rowcolor{popblueL} & \textbf{THÈSE : les réseaux sociaux unissent} & \textbf{ANTITHÈSE : les réseaux sociaux divisent} \\
\hline
\textbf{Idée directrice} & Ils relient les jeunes par-delà les distances et créent des communautés. & Ils enferment chacun dans son écran et opposent les groupes les uns aux autres. \\
\hline
\textbf{Argument 1} & Un étudiant de Maroua dialogue chaque jour avec ses camarades de Buea grâce à WhatsApp. & Les disputes et insultes entre groupes sur Facebook créent des haines durables. \\
\hline
\textbf{Argument 2} & Les jeunes s'organisent collectivement (groupes d'entraide scolaire, associations). & La comparaison permanente aux vies « parfaites » des autres isole et décourage. \\
\hline
\textbf{Exemple camerounais} & Une campagne de solidarité en ligne pour payer les fournitures d'un élève démuni. & Des rumeurs diffusées en ligne qui dressent des communautés l'une contre l'autre. \\
\hline
\textbf{Limite (transition)} & Mais cette union reste virtuelle et fragile... & Mais n'est-ce pas l'usage, et non l'outil, qui divise ? \\
\hline
\end{tabularx}
\end{center}

\courssub{Exemple guidé : rédiger une antithèse complète}

\begin{exoresolu}[Rédiger l'antithèse d'un sujet camerounais]
\textbf{Énoncé.} Sujet : \emph{Les réseaux sociaux unissent-ils ou divisent-ils la jeunesse ?} La thèse a soutenu que les réseaux sociaux unissent la jeunesse camerounaise. La transition s'achève ainsi : « Cependant, cette union virtuelle ne cache-t-elle pas une division réelle ? » Rédige le paragraphe d'antithèse (environ 10 lignes).
\tcblower
\textbf{Solution détaillée.}
\emph{Cependant}, force est de constater que les réseaux sociaux divisent la jeunesse au moins autant qu'ils ne l'unissent. D'abord, ils rassemblent les jeunes en groupes fermés qui s'opposent : sur les plateformes, les clans se forment autour d'une école, d'un quartier ou d'une équipe de football, et les échanges tournent souvent à l'insulte, comme on l'observe lors des rivalités entre lycées de Yaoundé après un match. Ensuite, les réseaux sociaux installent une comparaison permanente : le jeune qui voit défiler les photos de réussite des autres finit par se sentir exclu et inférieur, ce qui détruit la solidarité au lieu de la construire. Enfin, les rumeurs et les fausses informations qui circulent rapidement en ligne dressent les communautés les unes contre les autres et alimentent des conflits bien réels. Ainsi, \emph{en revanche} de l'union espérée, l'usage quotidien des réseaux semble produire méfiance, jalousie et hostilité. \emph{Néanmoins}, faut-il accuser l'outil lui-même, ou plutôt l'usage que nous en faisons ?

\medskip
\textbf{Commentaire de méthode.} Remarque : (1) le paragraphe s'ouvre sur un connecteur d'opposition (\emph{cependant}) ; (2) il contient trois arguments distincts, chacun illustré par un exemple camerounais ; (3) il se termine par une question qui prépare la synthèse.
\end{exoresolu}

\courssub{La transition vers la synthèse}

L'antithèse ne doit pas se clore brutalement. Sa dernière phrase a une fonction précise : préparer la \cle{synthèse}. Pour cela, on procède en deux temps :

\begin{enumerate}
\item \textbf{Constater l'insuffisance des deux positions prises isolément.} La thèse a vu juste sur un point, l'antithèse sur un autre ; mais chacune est \textbf{partielle} et \textbf{unilatérale}. Aucune ne rend compte, à elle seule, de toute la réalité du problème.
\item \textbf{Formuler la question qui appelle le dépassement.} Cette question annonce la synthèse : « Faut-il alors choisir entre ces deux positions, ou est-il possible de les concilier ? », « Le problème ne vient-il pas de l'usage plutôt que de l'outil ? »
\end{enumerate}

\begin{attention}[Ce que la transition ne doit pas faire]
\begin{itemize}
\item Ne \textbf{tranche pas} encore le débat : dire « l'antithèse a raison » à ce stade ruine la synthèse, qui n'aurait plus rien à apporter.
\item Ne \textbf{répète pas} l'antithèse : la transition est un bilan critique, pas un résumé.
\item Ne fais pas durer la transition : deux ou trois phrases suffisent.
\end{itemize}
\end{attention}

\begin{exercice}[À toi de jouer]
Sujet : \emph{La tradition est-elle un obstacle au progrès ?} La thèse a soutenu que la tradition freine le progrès en enfermant les jeunes dans des coutumes dépassées. Rédige : (a) la phrase de transition qui ouvre l'antithèse ; (b) un paragraphe d'antithèse de trois arguments avec exemples camerounais ; (c) la phrase de transition vers la synthèse.
\end{exercice}

\begin{corrige}[Exercice]
\textbf{(a) Transition d'ouverture.} « Cependant, rejeter la tradition comme un simple frein n'est-ce pas oublier qu'elle peut aussi fonder et orienter le progrès ? »

\textbf{(b) Antithèse.} En revanche, la tradition peut être une force qui soutient le progrès. D'abord, elle transmet des savoir-faire précieux : les techniques agricoles héritées des anciens permettent encore à de nombreux paysans de l'Ouest de cultiver efficacement leurs terres. Ensuite, elle donne une identité et des repères moraux : les valeurs de solidarité transmises dans les familles camerounaises, comme l'entraide lors des tontines, soutiennent concrètement les projets économiques des jeunes. Enfin, un progrès qui ignore la tradition est souvent rejeté : les innovations techniques réussissent mieux lorsqu'elles respectent les couttumes locales. »
\end{corrige}


\coursec{La synthèse et la conclusion}

Nous venons de voir que l'antithèse oppose à la thèse une position contraire, également argumentée. Le devoir se trouve alors dans une situation de tension~: deux discours solides s'affrontent et semblent irréconciliables. Faut-il choisir brutalement l'un contre l'autre~? Non. Le propre de la dissertation philosophique est de \emph{dépasser} cette opposition par un travail de réflexion personnelle~: c'est le rôle de la \cle{synthèse}. Puis le devoir s'achève par la \cle{conclusion}, qui répond clairement à la question posée. Cette dernière étape est décisive~: c'est elle qui montre au correcteur que l'élève n'est pas un simple rapporteur d'opinions, mais un penseur capable de prendre position.

\courssub{Définition et rôle de la synthèse}

Commençons par une image simple. Deux villageois se disputent~: l'un affirme que la pluie est une bénédiction car elle fait pousser le maïs~; l'autre soutient qu'elle est une malédiction car elle rend les routes de latérite impraticables. Qui a raison~? Un sage du village répondrait~: les deux ont raison, chacun à son niveau~; la pluie est bonne pour les champs et mauvaise pour les pistes, et tout dépend du point de vue adopté. Ce sage vient d'accomplir un geste de \emph{synthèse}~: il n'a pas choisi un camp, il a intégré la part de vérité de chaque position dans une réponse plus large.

\begin{definition}[La synthèse]
La \cle{synthèse} est le moment de la dissertation où l'élève prend position de manière réfléchie et personnelle, en \textbf{intégrant la part de vérité des deux positions opposées} (thèse et antithèse) et en dépassant leur conflit apparent. Elle n'est ni une moyenne molle entre les deux thèses, ni une simple addition, mais une position \emph{nuancée et justifiée} qui montre comment les deux points de vue peuvent être conciliés ou dépassés.
\end{definition}

La synthèse répond donc à trois exigences~:
\begin{enumerate}
\item \textbf{Reconnaître le vrai de chaque position}~: la thèse et l'antithèse, si elles ont été bien construites, contiennent chacune une part de vérité qu'il faut conserver.
\item \textbf{Identifier la source du désaccord}~: le plus souvent, les deux positions s'opposent parce qu'elles envisagent la notion sous des angles différents (par exemple, la technique du point de vue du corps et du point de vue de l'esprit).
\item \textbf{Formuler une position personnelle argumentée}~: l'élève montre comment, une fois les angles distingués, le conflit se résout ou se dépasse.
\end{enumerate}

\begin{attention}[Ce que la synthèse n'est pas]
La synthèse n'est pas un compromis paresseux du type «~les deux ont un peu raison~». Dire simplement «~il y a du vrai dans les deux~» sans expliquer \emph{pourquoi} ni \emph{comment} est une synthèse ratée. Il faut montrer précisément sous quelle condition, à quel niveau ou dans quelle limite chaque thèse est vraie.
\end{attention}

\courssub{Les deux formes de synthèse}

Il existe deux grandes manières de construire une synthèse.

\textbf{1. La synthèse conciliatrice (ou de complémentarité).} Elle consiste à montrer que la thèse et l'antithèse ne s'opposent pas réellement, car elles portent sur des \emph{aspects différents} du même problème. Chacune est vraie dans son domaine, et elles se complètent. Schéma type~: «~La thèse a raison de dire X lorsqu'on considère A~; l'antithèse a raison de dire Y lorsqu'on considère B. Le problème était donc mal posé~: il fallait distinguer A et B.~»

\begin{exemplebox}[Synthèse conciliatrice]
Sujet~: «~La liberté consiste-t-elle à faire ce que l'on veut~?~» Thèse~: oui, être libre c'est agir selon ses désirs. Antithèse~: non, faire ce que l'on veut, c'est obéir à ses passions, donc être esclave de soi-même. Synthèse conciliatrice~: les deux positions distinguent en réalité deux sens du mot «~vouloir~»~: vouloir par impulsion (caprice) et vouloir par décision réfléchie (volonté). La vraie liberté est la volonté éclairée, qui réconcilie le désir et la raison.
\end{exemplebox}

\textbf{2. La synthèse de dépassement.} Elle consiste à adopter un \emph{point de vue nouveau}, plus élevé ou plus large, qui rend compte à la fois du vrai et des limites de chaque position. Le conflit n'est pas seulement dissous~: il est dépassé parce qu'on change de perspective (par exemple, on passe du point de vue individuel au point de vue social, ou du point de vue théorique au point de vue pratique).

\begin{exemplebox}[Synthèse de dépassement]
Sujet~: «~L'État doit-il imposer la paix par la force~?~» Thèse~: oui, sans force l'ordre est impossible. Antithèse~: non, la paix imposée par la force est une injustice. Synthèse de dépassement~: au lieu d'opposer force et justice, il faut penser la force \emph{au service du droit}~: la force légitime est celle qui est réglée par des lois justes. Le vrai problème n'est donc plus «~force ou paix~», mais «~quelles conditions rendent la force légitime~?~».
\end{exemplebox}

\begin{popfigure}
\begin{center}
\begin{tikzpicture}[scale=1.0]
% Triangle dialectique
\coordinate (T) at (-4.5,0);
\coordinate (A) at (4.5,0);
\coordinate (S) at (0,4.2);
\draw[line width=1.6pt, popline] (T) -- (A);
\draw[line width=1.2pt, popblue, dashed] (T) -- (S);
\draw[line width=1.2pt, poporange, dashed] (A) -- (S);
% Noeuds
\node[draw, fill=popblueL, rounded corners=3pt, align=center, text width=4.2cm] at (T) {\textbf{THÈSE}\\ \small Première position\\ \small (argumentée)};
\node[draw, fill=poporangeL, rounded corners=3pt, align=center, text width=4.2cm] at (A) {\textbf{ANTITHÈSE}\\ \small Position opposée\\ \small (argumentée)};
\node[draw, fill=popgreenL, rounded corners=3pt, align=center, text width=4.6cm] at (S) {\textbf{SYNTHÈSE}\\ \small Dépassement du conflit~:\\ \small position nuancée et personnelle};
% Flèches de montée
\draw[-{Stealth[length=3mm]}, line width=1.4pt, popgreen] (-3.2,0.9) -- (-1.2,3.1);
\draw[-{Stealth[length=3mm]}, line width=1.4pt, popgreen] (3.2,0.9) -- (1.2,3.1);
\node[popmuted, align=center, text width=4.6cm, font=\small] at (0,-1.1) {Le conflit thèse / antithèse};
\end{tikzpicture}
\end{center}
\legende{Le triangle dialectique~: la thèse et l'antithèse forment la base conflictuelle~; la synthèse, au sommet, intègre le vrai des deux positions et les dépasse (flèches vertes).}
\end{popfigure}

\courssub{Rédiger la conclusion}

La \cle{conclusion} est le dernier paragraphe du devoir. Elle est courte (cinq à dix lignes) mais décisive, car c'est souvent elle que le correcteur lit avec le plus d'attention. Elle remplit trois fonctions.

\begin{methode}[Construire la conclusion en trois étapes]
\begin{enumerate}
\item \textbf{Rappeler brièvement le chemin parcouru}~: résumer en une ou deux phrases le fil du devoir («~Nous avons d'abord montré que..., puis que...~»).
\item \textbf{Répondre clairement à la problématique} posée dans l'introduction, en une ou deux phrases nettes~: c'est la réponse personnelle de l'élève, issue de la synthèse. La réponse doit être explicite~: oui, non, ou une réponse nuancée mais assumée.
\item \textbf{Ouvrir éventuellement} la réflexion vers une question connexe, nouvelle mais liée au sujet (l'\emph{ouverture}). Elle est facultative mais valorisée lorsqu'elle est pertinente.
\end{enumerate}
\end{methode}

\begin{attention}[Les deux erreurs fatales de la conclusion]
\begin{itemize}
\item \textbf{Introduire une idée totalement nouvelle}~: la conclusion ne doit contenir aucun argument inédit. Tout ce qu'elle affirme doit avoir été démontré dans le développement. Une idée neuve en conclusion donne l'impression d'un devoir improvisé.
\item \textbf{Répéter l'introduction mot pour mot}~: la conclusion n'est pas une photocopie de l'introduction. Entre les deux, un travail de réflexion a eu lieu~: la réponse finale doit être plus riche et plus précise que la simple question de départ.
\end{itemize}
\end{attention}

\begin{propriete}[Cohérence globale du devoir]
La conclusion doit \textbf{découler nécessairement du développement}~: elle est le fruit de la synthèse, elle-même issue du dialogue entre la thèse et l'antithèse. Avant de rédiger, l'élève doit vérifier la chaîne~: problématique (introduction) $\rightarrow$ thèse $\rightarrow$ antithèse $\rightarrow$ synthèse $\rightarrow$ réponse (conclusion). Si la réponse finale ne correspond pas à ce qui a été démontré, le devoir est incohérent et perd sa valeur argumentative.
\end{propriete}

\begin{popfigure}
\begin{center}
\begin{tikzpicture}[node distance=0.55cm]
\tikzstyle{etape}=[draw, rounded corners=3pt, align=center, text width=3.4cm, minimum height=1.1cm, font=\small]
\tikzstyle{fleche}=[-{Stealth[length=2.5mm]}, line width=1.2pt, popdark]
\node[etape, fill=popblueL, align=center] (intro){\textbf{Introduction}\\ Accroche, problématique, annonce du plan};
\node[etape, fill=popblueL, below=of intro, align=center] (these){\textbf{Thèse}\\ Première réponse argumentée};
\node[etape, fill=poporangeL, below=of these, align=center] (anti){\textbf{Antithèse}\\ Critique et réponse opposée};
\node[etape, fill=popgreenL, below=of anti, align=center] (synth){\textbf{Synthèse}\\ Dépassement, position personnelle};
\node[etape, fill=popgoldL, below=of synth, align=center] (concl){\textbf{Conclusion}\\ Réponse à la problématique + ouverture};
\draw[fleche] (intro) -- (these);
\draw[fleche] (these) -- (anti);
\draw[fleche] (anti) -- (synth);
\draw[fleche] (synth) -- (concl);
% Vérification de cohérence
\draw[fleche, poppink, dashed] (concl.east) .. controls (7.2,-3.5) .. (intro.east);
\node[poppink, font=\small, align=center, text width=3.2cm] at (7.6,-3.5) {Vérification~: la réponse finale correspond-elle à la question posée~?};
\end{tikzpicture}
\end{center}
\legende{Le parcours complet de la dissertation, de l'introduction à la conclusion. La flèche pointillée symbolise la vérification finale de cohérence~: la conclusion doit répondre exactement à la problématique de l'introduction.}
\end{popfigure}

\courssub{Exemple guidé complet~: «~La technique libère-t-elle l'homme~?~»}

Appliquons maintenant l'ensemble de ces acquis à un sujet classique, proche de la vie des élèves~: \emph{«~La technique libère-t-elle l'homme~?~»}. Supposons le développement déjà rédigé~: la thèse a montré que la technique libère l'homme des contraintes naturelles (la moto réduit la fatigue des longs trajets entre Douala et Yaoundé, le téléphone portable permet de communiquer instantanément, la machine soulage le travail manuel)~; l'antithèse a montré que la technique peut aussi asservir (dépendance au téléphone, chômage provoqué par les machines, pollution, l'homme devenu esclave de ses propres outils).

\textbf{Rédaction de la synthèse.} Il faut d'abord identifier la source du désaccord~: la thèse considère la technique comme un \emph{moyen} au service de l'homme, l'antithèse la considère comme une \emph{puissance autonome} qui échappe à son créateur. Les deux ont raison, mais à des conditions différentes. On peut alors formuler une synthèse de dépassement~:

\begin{exemplebox}[Synthèse rédigée]
«~Ainsi, la technique apparaît à la fois comme un instrument de libération et comme un risque d'asservissement. Mais cette contradiction se comprend si l'on remarque que la technique n'est en elle-même ni bonne ni mauvaise~: tout dépend de l'usage que l'homme en fait. La technique libère l'homme de la nature lorsqu'elle reste un moyen dirigé par une volonté réfléchie~; elle l'asservit lorsqu'elle échappe à toute réflexion éthique et devient une fin en soi. Le véritable problème n'est donc pas la technique elle-même, mais la sagesse de celui qui l'emploie~: libérée par l'outil, l'humanité doit encore apprendre à se libérer de sa propre puissance.~»
\end{exemplebox}

\textbf{Rédaction de la conclusion.} Elle reprend le chemin parcouru, répond nettement à la question, puis ouvre~:

\begin{exemplebox}[Conclusion rédigée]
«~En définitive, notre réflexion nous a conduit de la promesse de libération portée par la technique au danger d'asservissement qu'elle fait courir, avant de découvrir que ce conflit renvoie à l'usage que l'homme fait de ses outils. Nous pouvons donc répondre à la question posée~: la technique libère l'homme, mais à condition qu'il demeure maître d'elle-même par la réflexion et la responsabilité~; sans cette vigilance, le libérateur devient geôlier. Reste alors une question que notre époque rend urgente~: l'éducation peut-elle former des hommes capables de rester maîtres des machines qu'ils ont créées~?~»
\end{exemplebox}

Remarquons la mécanique de cette conclusion~: la première phrase résume le parcours (thèse, antithèse, synthèse)~; la deuxième répond explicitement à la problématique par une réponse nuancée mais ferme («~oui, mais à condition...~»)~; la dernière ouvre sur une question connexe (l'éducation), sans introduire d'argument nouveau.

\begin{savaistu}[La conclusion, un enjeu d'examen]
Au Probatoire comme au Baccalauréat technique, les correcteurs considèrent qu'une conclusion sans réponse explicite à la problématique est un \textbf{devoir inachevé}~: l'élève a discuté, mais n'a pas conclu. Inversement, une conclusion claire, courte et cohérente avec le développement est toujours récompensée, même si la position défendue est discutable~: en philosophie, on n'évalue pas l'opinion, mais la qualité du raisonnement qui la soutient.
\end{savaistu}

\begin{exoresolu}[Repérer et corriger une mauvaise conclusion]
Énoncé. Voici la conclusion d'un élève au sujet «~La technique libère-t-elle l'homme~?~»~: «~En conclusion, la technique est un vaste problème qui préoccupe tous les philosophes depuis l'Antiquité. De plus, le sport est aussi très important pour la santé de l'homme moderne.~» Relevez les erreurs commises et proposez une correction.
\tcblower
Solution détaillée. Cette conclusion comporte trois erreurs majeures.
\begin{enumerate}
\item \textbf{Elle ne répond pas à la problématique}~: à la question «~la technique libère-t-elle l'homme~?~», l'élève ne répond ni oui, ni non, ni de manière nuancée. Dire que le problème est «~vaste~» et «~préoccupe les philosophes~» est une phrase creuse qui évite de prendre position.
\item \textbf{Elle introduit une idée totalement nouvelle}~: le sport et la santé n'ont aucun lien avec le développement~; c'est exactement le piège signalé plus haut.
\item \textbf{Elle ne rappelle pas le chemin parcouru}~: aucune trace de la thèse, de l'antithèse ou de la synthèse.
\end{enumerate}
Correction possible~: «~Nous avons vu que la technique libère l'homme des contraintes naturelles, mais qu'elle peut aussi l'asservir lorsqu'elle devient une fin en soi. Tout dépend donc de l'usage~: la technique libère l'homme à condition qu'elle reste guidée par la réflexion éthique. On peut alors se demander si l'école forme suffisamment les hommes à cet usage responsable.~» Cette version rappelle le parcours, répond clairement et ouvre sans introduire d'argument nouveau.
\end{exoresolu}

\begin{aretenir}[L'essentiel de la section]
\begin{itemize}
\item La \cle{synthèse} dépasse le conflit thèse/antithèse en intégrant la part de vérité de chaque position dans une réponse nuancée et personnelle.
\item Deux formes~: la \textbf{synthèse conciliatrice} (les deux positions se complètent car elles portent sur des aspects différents) et la \textbf{synthèse de dépassement} (un point de vue nouveau résout le conflit).
\item La \cle{conclusion} rappelle le chemin parcouru, répond explicitement à la problématique, et peut s'achever par une ouverture.
\item Interdits~: aucune idée nouvelle en conclusion, jamais de répétition mot pour mot de l'introduction.
\item La conclusion doit découler du développement~: vérifier toujours la chaîne problématique $\rightarrow$ thèse $\rightarrow$ antithèse $\rightarrow$ synthèse $\rightarrow$ réponse.
\end{itemize}
\end{aretenir}

\begin{exercice}[À vous de rédiger]
Sujet~: «~Le travail est-il une contrainte ou un épanouissement~?~» On donne la thèse (le travail est une contrainte pénible imposée par la nécessité de survivre) et l'antithèse (le travail est ce par quoi l'homme se réalise et transforme le monde, comme le menuisier de quartier qui tire fierté de ses meubles). Rédigez~: 1) une synthèse de huit à douze lignes~; 2) une conclusion de cinq à huit lignes comportant une réponse explicite et une ouverture.
\end{exercice}

\begin{corrige}[Exercice~: Le travail, contrainte ou épanouissement~?]
Éléments de correction. \textbf{Synthèse attendue}~: identifier la source du désaccord — la thèse envisage le travail du point de vue de la nécessité (travailler pour vivre), l'antithèse du point de vue de l'activité créatrice (vivre à travers ce que l'on fait). Synthèse de dépassement possible~: le travail est ambivalent~; il est contrainte lorsqu'il est subi, répétitif et imposé sans sens, mais il devient épanouissement lorsque le travailleur y reconnaît le fruit de son effort et y exprime ses capacités. La différence ne tient donc pas au travail lui-même, mais aux conditions dans lesquelles il s'exerce et au sens que l'homme lui donne. \textbf{Conclusion attendue}~: rappel du parcours (contrainte, puis épanouissement, puis ambivalence)~; réponse explicite, par exemple~: «~le travail est à la fois contrainte et épanouissement~: contrainte par sa nécessité, épanouissement possible par le sens qu'on lui donne~»~; ouverture pertinente, par exemple~: «~la société peut-elle organiser le travail pour qu'il soit moins souvent une souffrance et plus souvent une réalisation de soi~?~». Toute conclusion introduisant une idée neuve (par exemple un paragraphe sur les loisirs) ou ne répondant pas à la question doit être signalée comme fautive.
\end{corrige}

\coursec{Application guidée : rédiger une dissertation complète}

Dans la section précédente, nous avons étudié la synthèse et la conclusion, dernières étapes de la construction théorique de la dissertation. Mais connaître les règles d'un jeu ne suffit pas pour y gagner : il faut s'entraîner à jouer. Cette dernière section est donc un véritable \cle{travail pratique} (TP) de philosophie : nous allons, ensemble, rédiger pas à pas une dissertation complète sur le sujet \emph{« Le progrès technique fait-il le bonheur de l'homme ? »}, en appliquant toutes les méthodes acquises dans cette leçon. Ce sujet est particulièrement adapté aux élèves de Première technique, car il touche directement à leur futur métier.

\courssub{Présentation du TP : objectifs et organisation}

\begin{definition}[Hors-sujet]
Le \cle{hors-sujet} est un développement qui ne répond pas à la problématique posée par le sujet. L'élève écrit beaucoup, parfois de belles choses, mais à côté de la question. C'est la \textbf{faute majeure} en dissertation philosophique : une copie hors sujet, même bien rédigée, ne peut obtenir une note satisfaisante au Probatoire ou au Baccalauréat technique. Toute la méthode que nous allons suivre vise à éviter cette faute.
\end{definition}

\textbf{Organisation du TP.} La classe est divisée en groupes de quatre à six élèves. Le travail se déroule en cinq étapes, du brouillon collectif jusqu'à la conclusion individuelle. Le professeur joue le rôle de guide : il corrige les orientations, mais ce sont les élèves qui produisent les idées.

\begin{methode}[Réussir son brouillon avant de rédiger]
\begin{enumerate}
\item \textbf{Noter toutes les idées} qui viennent à l'esprit, sans les juger : mots-clés, exemples, noms d'auteurs, objections, souvenirs de cours.
\item \textbf{Trier} ces idées : éliminer celles qui ne répondent pas à la question posée (risque de hors-sujet).
\item \textbf{Classer} les idées restantes en trois colonnes : celles qui vont dans le sens du « oui » (thèse), du « non » (antithèse), et celles qui permettent de dépasser l'opposition (synthèse).
\item \textbf{Ordonner} chaque colonne : une idée principale par paragraphe, avec son argument et son exemple.
\item Seulement ensuite, \textbf{rédiger au propre}, en suivant le plan établi.
\end{enumerate}
\end{methode}

\begin{attention}[Erreur fréquente : rédiger sans plan]
Beaucoup d'élèves, pressés, commencent à rédiger directement au propre, sans brouillon ni plan. C'est la source de deux catastrophes : le \cle{hors-sujet} (on dérive au fil de la plume) et la \cle{désorganisation} (on répète les mêmes idées, on oublie l'antithèse, la synthèse est bâclée faute de temps). Au Baccalauréat technique, consacrez toujours \textbf{45 minutes à 1 heure} au brouillon et au plan détaillé : c'est un investissement, pas une perte de temps.
\end{attention}

\courssub{Étape 1 : brainstorming et analyse des termes en groupes}

Chaque groupe reçoit le sujet écrit au tableau : \emph{« Le progrès technique fait-il le bonheur de l'homme ? »}. La première tâche consiste à analyser les trois termes essentiels, comme nous l'avons appris dans la section sur l'introduction.

\textbf{Analyse du terme « progrès technique ».} Il désigne l'ensemble des améliorations que l'homme apporte à ses outils, machines et techniques pour mieux maîtriser la nature et satisfaire ses besoins : la machine à vapeur, l'électricité, le téléphone portable, Internet, le motoculteur, Orange Money et MTN MoMo (mobile money). Le mot « progrès » suppose une idée d'\emph{amélioration} : la technique irait de mieux en mieux.

\textbf{Analyse du terme « bonheur ».} Le bonheur est un état de satisfaction durable, de plénitude de l'existence. Il faut le distinguer du simple \emph{plaisir} (sensation passagère) et du \emph{confort} (bien-être matériel). Question clé : le confort matériel suffit-il au bonheur ?

\textbf{Analyse du terme « homme ».} Il s'agit de l'être humain en général, pas seulement du technicien ou du consommateur. L'homme est un être de besoins, mais aussi un être moral, social et spirituel.

\begin{exemplebox}[Brainstorming : récolte des idées d'un groupe]
Un groupe de Première F4 (Génie mécanique) de Yaoundé note au brouillon : téléphone portable et communication facile ; motos-taxis et déplacements rapides ; médicaments et vaccins ; mais aussi : chômage causé par les machines ; accidents de la route ; pollution du Wouri à Douala ; dépendance aux réseaux sociaux ; guerres et armes modernes ; inégalités d'accès à la technologie entre ville et village. Le groupe constate que les idées se répartissent naturellement en deux camps opposés : c'est exactement la matière de la thèse et de l'antithèse.
\end{exemplebox}

La formulation de la problématique découle de cette analyse : le progrès technique apporte indéniablement le confort, mais le confort est-il le bonheur ? La technique ne crée-t-elle pas aussi de nouveaux malheurs ?

\courssub{Étape 2 : rédaction de l'introduction}

L'introduction, rédigée collectivement au tableau, suit les quatre temps étudiés : amorce, analyse des termes, problématique, annonce du plan.

\begin{exoresolu}[Rédiger l'introduction du sujet]
Énoncé : rédigez une introduction complète (8 à 12 lignes) pour le sujet « Le progrès technique fait-il le bonheur de l'homme ? ».
\tcblower
\textbf{Solution détaillée.} \emph{Amorce} : « Depuis la maîtrise du feu jusqu'au téléphone portable, l'homme n'a cessé de transformer son existence par la technique. À Douala comme dans les villages du Nord-Cameroun, le mobile money, les motos et les médicaments modernes ont changé la vie quotidienne. » \emph{Analyse des termes} : « Le progrès technique désigne l'amélioration continue des outils et des machines ; le bonheur, lui, est un état de satisfaction durable qui dépasse le simple confort matériel. » \emph{Problématique} : « Dès lors, une question s'impose : le progrès technique fait-il le bonheur de l'homme ? Le confort qu'il procure suffit-il à rendre l'existence heureuse, ou engendre-t-il de nouveaux maux ? » \emph{Annonce du plan} : « Nous montrerons d'abord que le progrès technique contribue au bonheur en libérant l'homme des besoins ; nous verrons ensuite qu'il peut aussi produire le malheur ; enfin, nous soutiendrons que le bonheur dépend moins de la technique elle-même que de l'usage que l'homme en fait. »
\end{exoresolu}

\courssub{Étape 3 : construction du plan détaillé sur brouillon}

Chaque groupe remplit maintenant le tableau du plan détaillé : pour chaque partie, une idée directrice, un ou deux arguments, un exemple précis et, si possible, une référence à un auteur.

\begin{popfigure}
\begin{center}
\begin{tabularx}{\linewidth}{|l|X|X|X|}
\hline
\rowcolor{popblueL} \textbf{Partie} & \textbf{Idée directrice} & \textbf{Arguments} & \textbf{Exemples / auteurs} \\
\hline
\textbf{Thèse} : le progrès technique fait le bonheur & \dotfill & \dotfill & \dotfill \\
\hline
\textbf{Antithèse} : le progrès technique fait le malheur & \dotfill & \dotfill & \dotfill \\
\hline
\textbf{Synthèse} : tout dépend de l'usage humain & \dotfill & \dotfill & \dotfill \\
\hline
\end{tabularx}
\end{center}
\legende{Tableau vierge du plan détaillé à compléter sur le brouillon : pour chaque partie, noter l'idée directrice, les arguments qui la justifient et les exemples concrets ou références aux auteurs.}
\end{popfigure}

\begin{exemplebox}[Exemple de plan détaillé complété]
\textbf{Thèse} : la technique libère l'homme de la misère. Arguments : elle soulage le travail pénible, soigne les maladies, rapproche les hommes. Exemples : le motoculteur qui remplace la houe épuisante ; les vaccins qui ont fait reculer la mortalité infantile au Cameroun ; Bergson : la technique prolonge l'effort vital de l'homme. \textbf{Antithèse} : la technique crée de nouveaux malheurs. Arguments : elle détruit des emplois, pollue, rend dépendant, fabrique des armes destructrices. Exemples : la pollution industrielle du Wouri ; l'isolement des jeunes absorbés par les écrans ; Rousseau : le progrès des sciences et des arts a corrompu les mœurs. \textbf{Synthèse} : la technique est un moyen, neutre en soi ; c'est l'usage, bon ou mauvais, qui décide. Exemple : le même téléphone portable permet de payer les frais de scolarité par mobile money ou de tricher à l'examen. Le bonheur exige une technique maîtrisée par la sagesse et la morale.
\end{exemplebox}

\courssub{Étape 4 : rédaction d'un paragraphe argumenté et correction croisée}

Chaque groupe rédige \textbf{un seul paragraphe} au propre : le groupe 1 rédige le premier paragraphe de la thèse, le groupe 2 le premier paragraphe de l'antithèse, le groupe 3 le paragraphe de synthèse. Un paragraphe de dissertation obéit à une structure rigoureuse en quatre temps :

\begin{enumerate}
\item \textbf{Annoncer} l'idée directrice du paragraphe (une phrase claire).
\item \textbf{Expliquer} l'idée par un raisonnement (pourquoi est-ce vrai ?).
\item \textbf{Illustrer} par un exemple concret ou une référence à un auteur.
\item \textbf{Conclure} provisoirement et \textbf{transitionner} vers le paragraphe suivant.
\end{enumerate}

\begin{exemplebox}[Paragraphe rédigé par un groupe (thèse)]
« Le progrès technique libère d'abord l'homme du poids du travail pénible, condition d'une vie plus heureuse. En effet, tant que l'homme doit consacrer toutes ses forces à survivre, il n'a ni le temps ni l'énergie de penser, de créer ou de jouir de l'existence. La machine prend en charge les tâches épuisantes et rend possible le loisir. Ainsi, au Cameroun, le motoculteur et le décortiqueur ont allégé le labeur des paysans qui, autrefois, cultivaient à la houe de l'aube au crépuscule. Comme le note Bergson, l'outil prolonge le corps humain et démultiplie sa puissance. Le progrès technique semble donc bien contribuer au bonheur en délivrant l'homme de la fatigue. Mais cette libération est-elle sans contrepartie ? »
\end{exemplebox}

\textbf{Correction croisée.} Les groupes échangent leurs paragraphes. Chaque groupe correcteur vérifie, à l'aide de la grille ci-dessous, que les quatre temps du paragraphe sont présents, que l'exemple illustre réellement l'argument, et que la langue est correcte. Les remarques sont notées en marge, puis discutées collectivement.

\courssub{Étape 5 : rédaction individuelle de la conclusion et relecture méthodique}

La conclusion se rédige \textbf{individuellement} : chacun répond à la problématique en reprenant le résultat de la synthèse, sans introduire d'idée nouvelle. Exemple de formulation attendue : « Le progrès technique ne fait ni le bonheur ni le malheur par lui-même : instrument entre les mains de l'homme, il appelle une maîtrise morale. C'est à l'homme sage, et non à la machine, qu'il appartient de faire le bonheur. »

\begin{methode}[La relecture méthodique en trois passages]
\begin{enumerate}
\item \textbf{Relecture du sens} : chaque paragraphe répond-il à la problématique ? Y a-t-il des répétitions, des passages hors sujet ?
\item \textbf{Relecture de la langue} : orthographe, accords, conjugaison, ponctuation. Vérifier en particulier les mots du vocabulaire philosophique (\emph{problématique}, \emph{antithèse}, \emph{cependant}).
\item \textbf{Relecture de la présentation} : paragraphes aérés, alinéas respectés, écriture lisible, pas de ratures excessives.
\end{enumerate}
\end{methode}

\begin{savaistu}[Le pouvoir de la relecture]
Les correcteurs du Baccalauréat l'affirment : une relecture attentive permet de corriger en moyenne \textbf{5 à 10 fautes} d'orthographe et de grammaire, et peut faire gagner \textbf{1 à 2 points} sur la copie. À l'inverse, une copie pleine de fautes donne au correcteur une impression de négligence qui pénalise l'ensemble, même si les idées sont bonnes. Gardez toujours les dix dernières minutes de l'épreuve pour relire.
\end{savaistu}

\begin{exemplebox}[Le volume d'une copie efficace]
Une copie de dissertation efficace au Baccalauréat technique fait en général \textbf{3 à 5 pages}, avec une écriture lisible et des paragraphes aérés. Il ne sert à rien d'écrire 8 pages confuses : le correcteur évalue la pertinence des idées, la solidité des arguments et la clarté de l'expression, non la quantité d'encre. Une copie de 4 pages bien structurée vaut mieux qu'une copie de 7 pages désorganisée.
\end{exemplebox}

\courssub{Grille d'auto-évaluation}

Pour terminer, chaque élève évalue sa propre production (ou celle de son groupe) à l'aide de la grille suivante, construite sur les critères réels de correction au Baccalauréat technique.

\begin{popfigure}
\begin{center}
\begin{tabularx}{\linewidth}{|l|X|c|}
\hline
\rowcolor{popgreenL} \textbf{Critère} & \textbf{Indicateurs de réussite} & \textbf{Barème} \\
\hline
\textbf{Pertinence} & Le sujet est compris, les termes sont définis, la problématique est claire, aucun hors-sujet & /5 \\
\hline
\textbf{Cohérence} & Plan dialectique respecté (thèse, antithèse, synthèse), transitions présentes, conclusion qui répond à la problématique & /5 \\
\hline
\textbf{Argumentation} & Chaque idée est justifiée par un raisonnement, illustrée par un exemple concret ou une référence à un auteur & /6 \\
\hline
\textbf{Langue et présentation} & Orthographe et grammaire correctes, vocabulaire philosophique approprié, écriture lisible, paragraphes aérés & /4 \\
\hline
\rowcolor{popgoldL} \textbf{Total} & & \textbf{/20} \\
\hline
\end{tabularx}
\end{center}
\legende{Grille d'auto-évaluation de la dissertation : quatre critères (pertinence, cohérence, argumentation, langue) pour un total de 20 points, conforme aux attentes du correcteur au Probatoire et au Baccalauréat technique.}
\end{popfigure}

\begin{aretenir}[L'essentiel de l'application guidée]
Rédiger une dissertation complète exige cinq étapes : analyser les termes et faire un brainstorming ; rédiger une introduction en quatre temps ; construire un plan détaillé thèse-antithèse-synthèse sur brouillon ; rédiger des paragraphes structurés (idée, explication, exemple, transition) ; conclure et relire méthodiquement. Le brouillon n'est pas du temps perdu : c'est la garantie contre le hors-sujet. Une bonne copie fait 3 à 5 pages, est lisible, aérée, et chaque affirmation y est justifiée par un argument et un exemple.
\end{aretenir}

\begin{exercice}[Dissertation complète à rédiger]
En vous inspirant de la démarche suivie dans ce TP, rédigez intégralement, en 2 heures, la dissertation sur le sujet : \emph{« Le travail est-il pour l'homme une contrainte ou un moyen de s'épanouir ? »}. Respectez les cinq étapes : analyse des termes, introduction, plan détaillé au brouillon, développement en trois parties, conclusion et relecture. Auto-évaluez ensuite votre copie avec la grille sur 20.
\end{exercice}

\begin{corrige}[Exercice : éléments de correction]
\textbf{Analyse des termes} : le \emph{travail} est l'activité par laquelle l'homme transforme la nature pour satisfaire ses besoins ; la \emph{contrainte} renvoie à l'obligation subie ; l'\emph{épanouissement} désigne le plein développement des capacités humaines. \textbf{Problématique} : le travail n'est-il qu'une peine imposée par la nécessité, ou peut-il être une source d'accomplissement ? \textbf{Thèse} : le travail est une contrainte — fatigue, exploitation, salaire de misère (exemple : le manœuvre sur les chantiers de Douala ; Marx et l'aliénation du travailleur). \textbf{Antithèse} : le travail épanouit — il humanise, crée la dignité et la reconnaissance sociale (exemple : l'artisan ébéniste fier de son meuble ; Hegel : par le travail, l'homme se forme lui-même). \textbf{Synthèse} : tout dépend des conditions du travail — un travail libre, créatif et justement rémunéré épanouit ; un travail forcé et répétitif aliène. \textbf{Conclusion} : le travail n'est pas une malédiction en soi ; c'est sa forme sociale qui en fait une contrainte ou un bonheur. Barème : pertinence /5, cohérence /5, argumentation /6, langue /4.
\end{corrige}

\newpage

\coursec{Activité d'intégration}

\coursec{Leçon 18 : La dissertation philosophique — Application guidée}

\courssub{Objectifs de l'activité}
Mobiliser l'ensemble des savoir-faire de la leçon (analyse de sujet, rédaction de l'introduction, construction du développement dialectique \emph{thèse / antithèse / synthèse}, rédaction de la conclusion) pour produire une dissertation philosophique complète, cohérente et argumentée, puis la justifier oralement devant la classe à l'aide d'une grille d'évaluation critériée.

\courssub{Situation}
Le club de philosophie du \textbf{Lycée Technique de Bafoussam} organise son concours annuel « \emph{Jeunes Penseurs} ». Le sujet imposé cette année est : \textbf{« L'usage du téléphone portable doit-il être interdit aux élèves au Cameroun ? » }.

Pour préparer la finale, l'enseignant référent constitue des équipes de quatre élèves. Chaque équipe dispose de \textbf{deux séances de 50 minutes} pour rendre une copie collective. Pour ancrer la réflexion dans le réel, l'enseignant fournit les résultats d'un \textbf{mini-sondage anonyme} réalisé la semaine précédente auprès de \textbf{40 élèves de Première} (toutes séries confondues) du même établissement. Voici les données brutes :

\begin{popfigure}
\begin{center}
\begin{tabularx}{\linewidth}{|l|X|}\hline
\textbf{Indicateur} & \textbf{Résultats (sur 40 élèves)} \\ \hline
Possession d'un smartphone personnel & 32 élèves (80\,\%) \\ \hline
Usage quotidien \textgreater 4 heures (hors cours) & 22 élèves (55\,\%) \\ \hline
Usage en classe (réseaux sociaux, jeux) & 18 élèves (45\,\%) \\ \hline
Usage pour réviser / dictionnaire / calculatrice & 26 élèves (65\,\%) \\ \hline
Déjà sanctionné(e) pour usage du téléphone en cours & 14 élèves (35\,\%) \\ \hline
A déjà été victime de cyberharcèlement via téléphone & 6 élèves (15\,\%) \\ \hline
Pense qu'une interdiction totale améliorerait les notes & 10 élèves (25\,\%) \\ \hline
Pense qu'un encadrement (charte d'usage) est préférable & 30 élèves (75\,\%) \\ \hline
\end{tabularx}
\end{center}
\legende{Résultats du mini-sondage « Téléphone portable et vie scolaire » — Lycée Technique de Bafoussam, 40 élèves de 1re, 2025.}
\end{popfigure}

La situation-problème est donc : \emph{Comment, à partir de ces données chiffrées et de votre culture philosophique, construire une dissertation rigoureuse qui ne se contente pas de lister des pour et des contre, mais qui articule une véritable problématique et propose une résolution dialectique ?}

\courssub{Consignes}
Chaque équipe suit scrupuleusement les six étapes ci-dessous. L'enseignant valide chaque étape avant le passage à la suivante.

\begin{enumerate}
    \item \textbf{Analyse du sujet et problématisation (15 min).} Identifiez les termes clés, les présupposés, les tensions. Formulez une \cle{problématique} unique sous forme de question philosophique (ex. : « L'interdiction est-elle la seule réponse éducative aux dérives numériques ? » ou « La liberté de l'élève est-elle compatible avec la discipline scolaire à l'ère numérique ? »).
    \item \textbf{Rédaction collective de l'introduction (15 min).} Produisez une introduction complète : \cle{amorce} (accroche contextualisée), \cle{analyse des termes}, \cle{restitution du sujet}, \cle{problématique}, \cle{annonce du plan dialectique} (Thèse / Antithèse / Synthèse).
    \item \textbf{Construction du plan détaillé au tableau (20 min).} Sur un papier affiché (ou au tableau), dressez le squelette du développement en trois parties. Pour chaque grande partie, indiquez : l'\cle{idée directrice} (sous-thèse), \textbf{deux arguments philosophiques} distincts, et \textbf{un exemple concret} tiré des données du sondage ou de la vie scolaire camerounaise (règlement intérieur, circulaires MINESEC, réalité terrain).
    \item \textbf{Rédaction répartie du développement (40 min).} Chaque membre de l'équipe rédige \emph{une} des trois parties (Thèse, Antithèse, Synthèse) sur une feuille séparée, en respectant la structure : \emph{Idée directrice $\rightarrow$ Argument 1 + Exemple $\rightarrow$ Argument 2 + Exemple $\rightarrow$ Transition}.
    \item \textbf{Assemblage, conclusion et relecture (20 min).} Recopiez l'introduction, assemblez les trois parties en assurant la cohésion (connecteurs logiques), rédigez la \cle{conclusion} (bilan synthétique + réponse à la problématique + \cle{ouverture}) et relisez l'ensemble (orthographe, syntaxe, respect du formalisme).
    \item \textbf{Présentation orale et évaluation par les pairs (20 min).} Un porte-parole présente la démarche (problématique, choix du plan, arguments forts) en 3 minutes. La classe note sur la \cle{grille critériée} (4 critères : Pertinence de la problématique / Structure dialectique / Qualité des arguments et exemples / Qualité de la langue et du formalisme). L'enseignant institutionnalise les corrections.
\end{enumerate}

\courssub{Ressources à mobiliser}
\begin{itemize}
    \item \textbf{Structure de l'introduction} : Amorce $\rightarrow$ Définition des termes (\emph{usage, interdiction, élève, Cameroun}) $\rightarrow$ Reformulation du sujet $\rightarrow$ Problématique $\rightarrow$ Annonce du plan (Thèse / Antithèse / Synthèse).
    \item \textbf{Dialectique} :
    \begin{itemize}
        \item \textbf{Thèse} : L'interdiction est nécessaire / légitime (protection, autorité, concentration).
        \item \textbf{Antithèse} : L'interdiction est contre-productive / injuste (outil pédagogique, liberté, fracture numérique, responsabilité).
        \item \textbf{Synthèse} : Dépasser l'alternative binaire par la \cle{régulation éducative} (charte d'usage, éducation au numérique, responsabilité partagée).
    \end{itemize}
    \item \textbf{Conclusion} : Résumé des étapes $\rightarrow$ Réponse affirmative nuancée à la problématique $\rightarrow$ Ouverture (ex. : « L'école camerounaise face au défi de l'IA générative »).
    \item \textbf{Grille d'évaluation (extrait)} : \textit{Problématique explicite et pertinente (4 pts) — Plan dialectique respecté et transitions fluides (4 pts) — Arguments philosophiques étayés d'exemples camerounais (8 pts) — Formalisme (intro/conclusion) et langue (4 pts)}.
\end{itemize}

\courssub{Démarche et corrigé commenté}
\begin{exoresolu}{Corrigé type et commentaires pédagogiques}
\textbf{Étape 1 — Analyse et problématisation}\\
\textit{Travail de l'équipe :} Le terme « \emph{interdit} » suppose une contrainte légale (règlement intérieur, circulaire ministérielle). « \emph{Doit} » marque l'obligation morale ou juridique. La tension porte sur le conflit entre \textbf{l'autorité éducative} (garante de l'ordre et de la concentration) et \textbf{l'autonomie de l'élève} (sujet de droits, acteur de sa formation).
\\ \textbf{Problématique retenue} : \emph{« L'interdiction totale du téléphone portable au lycée est-elle une réponse pédagogique efficace ou un refus d'éduquer à la responsabilité numérique ? »}
\\ \textit{Commentaire :} Cette problématique est recevable : elle sort du binaire « pour/contre » pour interroger la \emph{finalité éducative} de la mesure.

\textbf{Étape 2 — Introduction (modèle rédigé)}\\
\textit{Amorce :} « Depuis la circulaire n\textsuperscript{o} 00000000/MINESEC du 12 janvier 2018 interdisant l'usage du téléphone portable dans les établissements scolaires publics et privés du Cameroun, le débat oppose tenants de la discipline stricte et partisans de l'intégration pédagogique. »\\
\textit{Analyse :} « L'\emph{usage} désigne l'acte d'employer un outil ; l'\emph{interdiction} est l'acte d'autorité qui le proscrit. L'\emph{élève} est ici un sujet en formation. »\\
\textit{Restitution :} « Le sujet demande si la prohibition légale est la seule voie pour garantir le bon fonctionnement de l'école. »\\
\textit{Problématique :} (Celle ci-dessus).\\
\textit{Annonce du plan :} « Nous montrerons d'abord que l'interdiction s'impose comme condition de l'ordre scolaire (Thèse) ; puis qu'elle ignore les potentialités éducatives de l'outil et risque l'injustice (Antithèse) ; enfin, nous proposerons une voie médiane : la régulation par une charte d'usage responsable (Synthèse). »

\textbf{Étape 3 — Plan détaillé au tableau (Squelette)}

\begin{popfigure}
\begin{center}
\begin{tabularx}{\linewidth}{|l|X|X|}\hline
\textbf{Partie} & \textbf{Idée directrice} & \textbf{Arguments + Exemples (Données sondage / Contexte CM)} \\ \hline
\textbf{THÈSE} & L'interdiction est légitime : elle protège le cadre scolaire. & 1. \textbf{Préserver l'attention et l'autorité.} Le téléphone est une « machine à distraire » (Stiegler). \textit{Exemple :} 45\,\% des sondés l'utilisent en cours (réseaux/jeux) ; 35\,\% déjà sanctionnés.\\
& & 2. \textbf{Protéger les plus vulnérables.} Lutte contre le cyberharcèlement et les inégalités (ceux sans smartphone). \textit{Exemple :} 15\,\% victimes de cyberharcèlement ; interdiction = filet de sécurité. \\ \hline
\textbf{ANTITHÈSE} & L'interdiction est inefficace et contre-productive : elle nie l'outil. & 1. \textbf{Le téléphone est un levier pédagogique.} Accès au savoir, dictionnaire, calculatrice. \textit{Exemple :} 65\,\% l'utilisent pour réviser/calculatrice ; fracture numérique si interdiction totale (élèves sans manuels).\\
& & 2. \textbf{Éduquer à l'autonomie plutôt que réprimer.} L'interdiction infantilise ; le règlement MINESEC est peu appliqué (tolérance tacite). \textit{Exemple :} 75\,\% des sondés réclament une \emph{charte d'usage} et non l'interdiction. \\ \hline
\textbf{SYNTHÈSE} & Dépasser l'interdit par la \textbf{corégulation éducative} (charte + formation). & 1. \textbf{Instituer une « Charte Numérique » cosignée.} Droits/devoirs : zones/temps autorisés (CDI, vie scolaire), interdits (salle de cours, examen). \textit{Exemple :} Modèle « Lycée de Nkolbisson » : casiers à l'entrée, plages « BYOD » (Bring Your Own Device) encadrées.\\
& & 2. \textbf{Former à l'hygiène numérique (EMI).} Cours d'Éducation aux Médias : gestion du temps d'écran, identité numérique, droit à l'image. \textit{Exemple :} Module « Citoyenneté numérique » en 1re Tech (programme MINESEC), impliquant parents (APE). \\ \hline
\end{tabularx}
\end{center}
\legende{Plan dialectique détaillé : Thèse / Antithèse / Synthèse avec arguments philosophiques et ancrage camerounais.}
\end{popfigure}

\textbf{Étape 4 — Extraits de rédaction (Développement)}

\textit{Thèse (extrait) :} « \textbf{La thèse soutient que l'interdiction est une condition nécessaire de l'ordre scolaire.} Premier argument : le téléphone portable fragmente l'attention, condition \emph{sine qua non} de l'apprentissage. Comme l'analyse Bernard Stiegler, les industries de l'attention captent le temps de cerveau disponible. Le sondage le confirme : près de la moitié des élèves (45\,\%) avoue l'utiliser en cours pour des loisirs, et plus d'un tiers (35\,\%) a déjà fait l'objet d'une sanction. Cela prouve que la présence de l'objet, même silencieux, structure une tentation permanente qui mine l'autorité du professeur. Second argument : l'interdiction protège les élèves les plus fragiles. Le cyberharcèlement, facilité par l'anonymat et la viralité des réseaux, touche 15\,\% de nos sondés. L'interdiction agit comme un \emph{pare-feu} institutionnel, réduisant les occasions de nuisance dans l'enceinte de l'école, lieu de droit et de protection. \textit{Transition :} Cependant, si l'interdiction protège, elle risque aussi d'exclure l'élève des compétences numériques indispensables. »

\textit{Antithèse (extrait) :} « \textbf{L'antithèse démontre que l'interdiction totale est une réponse archaïque qui nie la réalité pédagogique de l'outil.} Premier argument : le smartphone est devenu une \emph{prothèse cognitive} (Leroi-Gourhan). Dans un contexte camerounais marqué par la pénurie de manuels et l'accès limité aux laboratoires informatiques, 65\,\% des élèves l'utilisent comme dictionnaire, calculatrice scientifique ou pour consulter des cours en ligne (Khan Academy, YouTube Edu). L'interdire, c'est aggraver la \emph{fracture numérique} entre élèves dotés de ressources à la maison et ceux pour qui le téléphone est l'unique porte d'entrée vers le savoir. Second argument : la prohibition infantilise le futur citoyen. Le droit camerounais (Loi n\textsuperscript{o} 2010/012 sur la cybercriminalité) sanctionne les \emph{actes}, pas la \emph{possession}. 75\,\% des sondés réclament une \emph{charte d'usage} : ils expriment le désir d'être traités en sujets responsables, capables d'autorégulation. L'interdiction générale, souvent contournée (téléphones « de poche »), décrédibilise la règle de droit. \textit{Transition :} Dès lors, ni la tolérance laxiste ni la répression aveugle ne suffisent ; il faut instituer un cadre normatif éducatif. »

\textit{Synthèse (extrait) :} « \textbf{La synthèse propose de substituer à l'interdit binaire une \emph{corégulation éducative} fondée sur la confiance et la compétence.} Premier argument : l'élaboration participative d'une \emph{Charte Numérique d'Établissement}. Contrairement au règlement imposé, cette charte définit des \emph{zones} (interdites : salles de classe, examen ; autorisées : CDI, foyer, vie scolaire) et des \emph{plages horaires} (récréation, temps de travail personnel encadré « BYOD »). L'expérience du Lycée de Nkolbisson (Yaoundé) montre que la mise en place de casiers à l'entrée des salles, couplée à des créneaux « téléphone autorisé pour recherche documentaire », a fait chuter les sanctions de 60\,\% en un an. Second argument : l'inscription de l'\emph{Education aux Médias et à l'Information (EMI)} au cœur du curriculum. Le programme de 1re Technique prévoit des modules de vie scolaire ; y intégrer la gestion de l'attention, le droit à l'image, la protection des données personnelles (Loi n\textsuperscript{o} 2024/016 sur la protection des données à caractère personnel), transforme l'élève en acteur lucide. Cela implique les parents (APE) pour une cohérence école/maison. »

\textbf{Étape 5 — Conclusion (modèle)}\\
\textit{Bilan :} « Au terme de cette analyse, il apparaît que l'interdiction totale, si elle répond à un souci légitime d'ordre (Thèse), se heurte à la réalité pédagogique et sociale de l'élève camerounais (Antithèse). La solution n'est pas dans le choix de l'un contre l'autre, mais dans leur dépassement dialectique : la \emph{régulation éducative}. »\\
\textit{Réponse :} « L'usage du téléphone ne \emph{doit} pas être interdit \emph{absolument}, mais \emph{encadré} par une charte cosignée et une formation au numérique responsable. »\\
\textit{Ouverture :} « Cette réflexion préfigure le défi majeur de l'école technique camerounaise : l'intégration de l'intelligence artificielle générative (ChatGPT, etc.). Si le téléphone pose la question de l'attention, l'IA pose celle de la \emph{pensée déléguée}. La même méthode — régulation plutôt qu'interdiction — devra s'appliquer. »

\textbf{Étape 6 — Grille d'évaluation (Auto-évaluation de l'équipe)}

\begin{popfigure}
\begin{center}
\begin{tabularx}{\linewidth}{|l|X|c|}\hline
\textbf{Critère} & \textbf{Indicateurs observés} & \textbf{Note /4} \\ \hline
Problématique & Explicite, philosophique, issue de l'analyse des termes, guide le plan. & 4 \\ \hline
Structure dialectique & 3 parties équilibrées, transitions logiques (Cependant / Dès lors), synthèse résolutive. & 4 \\ \hline
Arguments / Exemples & Arguments philosophiques (Stiegler, Leroi-Gourhan, Droit) + Données sondage (\%, n) + Contexte CM (Circulaires, Lycée Nkolbisson, Loi 2024). & 8 \\ \hline
Formalisme / Langue & Intro/Conclusion canoniques, connecteurs, vocabulaire précis (fracture numérique, prothèse cognitive, corégulation), syntaxe soignée. & 4 \\ \hline
\textbf{TOTAL} & & \textbf{20/20} \\ \hline
\end{tabularx}
\end{center}
\legende{Grille d'évaluation critériée — Total sur 20 points.}
\end{popfigure}
\tcblower
\textbf{Commentaires pédagogiques de l'enseignant (institutionnalisation) :}
\begin{itemize}
    \item \textbf{Sur la problématique :} Évitez les formulations « Pour ou contre ? ». Préférez une question qui met en tension deux valeurs (ex. \emph{Sécurité vs Liberté}, \emph{Autorité vs Autonomie}, \emph{Interdit vs Éducation}).
    \item \textbf{Sur les exemples :} Les pourcentages du sondage (45\,\%, 65\,\%) sont des \emph{faits sociaux} : ils illustrent, ils ne \emph{prouvent} pas philosophiquement. Il faut les articuler à un concept (ex. : « Ce chiffre de 45\,\% illustre la thèse de Stiegler sur la captation de l'attention »).
    \item \textbf{Sur la synthèse :} C'est la partie la plus faible chez les candidats. Elle ne doit pas être un « mélange » (un peu des deux), mais une \emph{résolution} : un concept nouveau (ici \emph{corégulation / charte}) qui intègre la vérité de la thèse (besoin de cadre) et de l'antithèse (besoin d'autonomie/outil).
    \item \textbf{Sur l'ouverture :} Elle doit être brève (2-3 lignes), liée au sujet, et ouvrir sur une question \emph{proche mais distincte} (ici l'IA), pas sur un sujet général (« La technologie est-elle bonne ou mauvaise ? »).
\end{itemize}
\end{exoresolu}

\courssub{Bilan de l'activité}
Cet atelier d'intégration a permis de vérifier l'acquisition effective des compétences visées par la leçon « La dissertation philosophique » :
\begin{itemize}
    \item \textbf{Savoir-faire 1 (Analyse) :} Les élèves ont appris à transformer un sujet d'actualité scolaire en problème philosophique, en dégageant les concepts opératoires (\emph{autorité, autonomie, régulation, fracture numérique}).
    \item \textbf{Savoir-faire 2 (Structure) :} La contrainte du plan dialectique affiché au tableau a forcé la rigueur : thèse/antithèse/synthèse ne sont pas de simples paragraphes « pour/contre » mais des moments logiques nécessaires.
    \item \textbf{Savoir-faire 3 (Argumentation) :} L'exigence d'adosser chaque argument à une référence philosophique (auteur/concept) \emph{et} à une donnée concrète camerounaise (sondage, texte de loi, expérience locale) ancre la philosophie dans le réel technique.
    \item \textbf{Savoir-faire 4 (Communication) :} La présentation orale et l'évaluation par les pairs (grille critériée) développent l'esprit critique et la capacité à justifier ses choix méthodologiques, compétence transversale essentielle pour le Probatoire.
\end{itemize}
L'activité révèle aussi des points de vigilance pour la suite : la gestion du temps de rédaction (souvent sous-estimée), la qualité des transitions (souvent mécaniques : « En outre », « Par ailleurs » au lieu de liens logiques forts), et la tendance à noyer la synthèse dans des généralités. Ces points feront l'objet d'exercices ciblés en séances suivantes.

\newpage

\coursec{Méthodes et exercices résolus}

\courssub{Méthodes et exercices résolus}

\begin{methode}[Analyser un sujet de dissertation philosophique]
\textbf{Objectif :} Dégager les termes clés, la thèse implicite, le problème et le plan de travail avant d'écrire.
\begin{enumerate}
    \item \textbf{Cercler les mots-clés} : identifier les concepts philosophiques principaux (ex. : \cle{liberté}, \cle{devoir}, \cle{technique}, \cle{bonheur}).
    \item \textbf{Définir chaque terme} : donner le sens philosophique précis (pas le sens courant) et noter les nuances (ex. : liberté = libre arbitre vs autonomie vs indépendance).
    \item \textbf{Repérer la structure logique} : sujet \emph{fermé} (question Oui/Non, ex. « La technique libère-t-elle l'homme ? ») ou \emph{ouvert} (ex. « Qu'est-ce que la technique ? »).
    \item \textbf{Formuler le problème (problématique)} : transformer le sujet en tension entre deux thèses opposées. Question type : « Dans quelle mesure [thèse]... alors que [antithèse]... ? »
    \item \textbf{Esquisser le plan dialectique} : Thèse (arguments pour), Antithèse (arguments contre/nuances), Synthèse (dépassement/hiérarchisation).
\end{enumerate}
\textbf{Réflexe Bac :} Passer 30 à 45 min sur le brouillon pour cette analyse. Ne jamais commencer la rédaction sans problématique écrite noir sur blanc.
\end{methode}

\begin{exoresolu}[Analyse du sujet : « La technique est-elle un danger pour la liberté humaine ? »]
\textbf{Énoncé :} Faites l'analyse complète de ce sujet type Probatoire (termes, définitions, type, problématique, plan détaillé).
\tcblower
\textbf{Solution détaillée :}

\textbf{1. Mots-clés et définitions philosophiques :}
\begin{itemize}
    \item \cle{Technique} : Ensemble de procédés, méthodes, outils rationnels permettant de transformer la nature pour satisfaire des besoins. Ce n'est pas seulement « machines », c'est un \emph{savoir-faire} (technè chez les Grecs). Distinction : technique comme moyen neutre vs technique comme système (Ellul : « système technicien »).
    \item \cle{Danger} : Risque, menace potentielle, ce qui peut nuire à l'intégrité ou à l'essence de quelque chose. Implique une évaluation axiologique (négative).
    \item \cle{Liberté humaine} : Pouvoir de choisir, d'agir par soi-même (libre arbitre), autonomie (se donner sa propre loi - Kant), indépendance à l'égard de la contrainte extérieure. Distinction : liberté \emph{de} (faire ce qu'on veut) vs liberté \emph{pour} (s'accomplir moralement).
\end{itemize}

\textbf{2. Type de sujet :} Sujet \emph{fermé} (interrogation directe « Est-ce que... ? »), appelant une réponse nuancée (ni un « oui » plat, ni un « non » dogmatique). Il suppose une tension : la technique comme puissance d'agir (pour la liberté) vs technique comme aliénation (contre la liberté).

\textbf{3. Problématique :}
« Dans quelle mesure la technique, en augmentant la puissance d'agir de l'homme, ne risque-t-elle pas de devenir une contrainte qui détermine ses fins et aliène sa volonté, transformant le moyen en fin ? »

\textbf{4. Plan dialectique détaillé (brouillon) :}

\textbf{Thèse : La technique libère l'homme (moyen d'émancipation).}
\begin{itemize}
    \item Argument 1 : Maîtrise de la nature $\rightarrow$ satisfaction des besoins vitaux (santé, alimentation, habitat) $\rightarrow$ condition matérielle de la liberté (Marx : « Le développement des forces productives »).
    \item Argument 2 : Gain de temps (automatisation) $\rightarrow$ disponibilité pour loisirs, culture, politique (Aristote : technique libère de la nécessité pour accéder à la vie contemplative).
    \item Argument 3 : Outils de communication/éducation $\rightarrow$ accès au savoir, émancipation intellectuelle (ex. internet au Cameroun : accès aux cours, veille citoyenne).
\end{itemize}

\textbf{Antithèse : La technique menace la liberté (aliénation, hétéronomie).}
\begin{itemize}
    \item Argument 1 : Déterminisme technique (Ellul, Heidegger) : le système technique s'impose comme environnement total, l'homme devient « standing-reserve » (ressource), efficacité = seule valeur. Perte du sens.
    \item Argument 2 : Aléation marxiste : l'ouvrier dépossédé de son travail, de son produit, de son geste. La machine dicte le rythme (taylorisme, algorithmes de gestion).
    \item Argument 3 : Surveillance et contrôle (Foucault, société de contrôle) : traçabilité numérique, notation sociale, manipulation algorithmique des choix (réseaux sociaux, IA). Atteinte à l'autonomie.
\end{itemize}

\textbf{Synthèse : Dépasser l'opposition moyen/fin $\rightarrow$ responsabilité politique et éthique.}
\begin{itemize}
    \item La technique n'est pas neutre (valeurs inscrites dans la conception), mais pas fatale non plus.
    \item Condition de la liberté : \cle{maîtrise politique} de la technique (démocratie technique, éthique de la conception, éducation critique).
    \item Exemple camerounais : Régulation du numérique (loi sur la cybercriminalité, protection données) vs fracture numérique. La liberté passe par l'appropriation critique, pas le refus.
\end{itemize}

\textbf{Conclusion de l'analyse :} Ce plan permet une dissertation structurée, argumentée par des auteurs (Ellul, Marx, Heidegger, Kant) et des exemples concrets (Cameroun), répondant exactement à l'exigence du Probatoire.
\end{exoresolu}

\begin{methode}[Rédiger une introduction « en entonnoir » (norme Bac)]
\textbf{Objectif :} Accrocher, définir, problématiser, annoncer le plan en 4 étapes logiques et fluides.
\begin{enumerate}
    \item \textbf{Amorce (accroche) :} Une phrase large sur le thème (citation d'auteur, fait d'actualité camerounaise, observation générale, étymologie). \emph{Interdit :} « Depuis la nuit des temps... », « De tout temps les hommes se demandent... ».
    \item \textbf{Définition des termes / Analyse du sujet :} Définir rigoureusement les concepts clés du sujet. Montrer qu'on a compris les nuances. Reformuler le sujet pour en montrer la portée.
    \item \textbf{Problématique :} La question centrale qui guide tout le devoir. Elle doit exprimer une \cle{tension} (opposition de deux points de vue légitimes). Souvent formulée : « Dans quelle mesure... ? », « Comment concilier... et... ? ».
    \item \textbf{Annonce du plan :} Présenter les trois grandes parties (Thèse, Antithèse, Synthèse) en utilisant les connecteurs logiques : « Nous montrerons d'abord que... [Thèse] ; nous verrons ensuite que... [Antithèse] ; nous verrons enfin que... [Synthèse] ».
\end{enumerate}
\textbf{Longueur :} Environ 15-20 lignes (une demi-page). Aucune argumentation développée ici.
\end{methode}

\begin{exoresolu}[Rédaction de l'introduction pour : « Le travail est-il une contrainte ou une réalisation ? »]
\textbf{Énoncé :} Rédigez une introduction complète, conforme aux attendus du Bac Technique.
\tcblower
\textbf{Solution détaillée :}

\textit{Amorce :} « Le travail, au sens étymologique du latin \emph{tripalium} (instrument de torture à trois pieux), évoque d'abord la peine, la contrainte physique imposée par la nécessité de survivre. Pourtant, l'expérience moderne, comme la pensée de Hegel, nous invite à voir dans le travail l'activité par excellence où l'homme se forme lui-même en transformant le monde. »

\textit{Définition / Analyse :} « Le \cle{travail} se définit comme une activité humaine consciente et volontaire, médiatisée par des outils, visant à transformer la nature pour produire des biens utiles. La \cle{contrainte} désigne ce qui s'impose à la volonté par une nécessité extérieure (besoins vitaux, lois économiques, hiérarchie). La \cle{réalisation} (ou accomplissement) suppose au contraire un déploiement de ses potentialités, une affirmation de soi dans l'œuvre accomplie. Le sujet oppose ainsi deux visages du travail : l'aliénation subie et l'épanouissement choisi. »

\textit{Problématique :} « Dès lors, \textbf{le travail est-il structurellement une aliénation dont l'homme ne peut s'affranchir, ou bien peut-il devenir, sous certaines conditions sociales et subjectives, le lieu privilégié de sa liberté et de son humanisation ?} »

\textit{Annonce du plan :} « Pour répondre, nous analyserons d'abord le travail comme \textbf{contrainte et aliénation} (I), en nous appuyant sur la condition ouvrière décrite par Marx et la pénibilité physique. Nous examinerons ensuite comment le travail peut être \textbf{source de réalisation et d'autonomie} (II), à travers la reconnaissance sociale (Hegel) et l'acquisition de compétences. Enfin, nous montrerons que la \textbf{transformation du travail en émancipation suppose une maîtrise collective de ses conditions} (III), dépassant l'opposition par l'organisation politique et technique du travail (ex. coopératives, droit du travail au Cameroun). »

\textbf{Commentaire du correcteur :} L'amorce ancre le sujet (étymologie + référence Hegel). Les définitions sont philosophiques (pas le dictionnaire Larousse). La problématique est dialectique (structurellement vs conditions). Le plan annonce le contenu des parties (pas juste « I, II, III »).
\end{exoresolu}

\begin{methode}[Construire et rédiger la Thèse (Partie I)]
\textbf{Objectif :} Défendre la réponse la plus spontanée/évidente au sujet avec des arguments ordonnés, illustrés, philosophiquement fondés.
\begin{enumerate}
    \item \textbf{Identifier la thèse directrice :} Réponse « Oui » ou « Non » au sujet fermé, ou premier angle d'attaque pour sujet ouvert. Ex. : « La technique libère » / « Le travail est contrainte ».
    \item \textbf{Décliner en 2 ou 3 sous-parties (arguments majeurs) :} Chaque sous-partie = une idée force distincte. Ordre logique : du plus évident/accessible au plus profond/structurel.
    \item \textbf{Structurer chaque paragraphe (méthode ARE) :}
        \begin{itemize}
            \item \textbf{A} - \textbf{Affirmation} : Phrase d'introduction claire (idée directrice du paragraphe).
            \item \textbf{R} - \textbf{Raisonnement / Référence} : Explication logique + \textbf{Appui autorisé} (Auteur, Concept, Œuvre). \emph{Ex : « Selon Marx dans \emph{Le Capital}... », « Kant définit l'autonomie comme... »}.
            \item \textbf{E} - \textbf{Exemple / Illustration} : Concret, si possible camerounais ou actuel. \emph{Ex : « L'introduction de la machine à coudre dans les ateliers de Douala... »}.
        \end{itemize}
    \item \textbf{Transition interne :} Lier les sous-parties par des connecteurs logiques (\emph{De plus, En outre, Par ailleurs, Plus profondément}).
    \item \textbf{Mini-conclusion de partie :} Résumer ce que cette thèse a établi et ouvrir sur sa limite (appel de l'antithèse).
\end{enumerate}
\textbf{Règle d'or :} Un auteur cité = une idée utilisée. Pas de « name-dropping » vide.
\end{methode}

\begin{exoresolu}[Rédaction de la Thèse : « Le travail est une contrainte et une aliénation » (Partie I du sujet précédent)]
\textbf{Énoncé :} Rédigez la première partie complète (2 sous-parties minimum) en respectant la méthode ARE.
\tcblower
\textbf{Solution détaillée :}

\textbf{I. Le travail : une contrainte subie, source d'aliénation}

\textbf{A. La contrainte de la nécessité vitale : le travail comme fardeau naturel.}
\textit{Affirmation :} Dans sa dimension la plus immédiate, le travail s'impose à l'homme non comme un choix, mais comme une condition de survie biologique.
\textit{Raisonnement / Référence :} Comme l'analyse \textbf{Aristote} dans \emph{La Politique}, la vie bonne (\emph{eu zên}) suppose l'affranchissement de la nécessité (\emph{anankê}). Les activités relevant de la \cle{poïèsis} (production) sont subordonnées aux besoins ; elles relèvent de l'esclavage à la nature. Le travailleur ne se appartient pas pendant qu'il produit ; il appartient à son besoin.
\textit{Exemple :} Au Cameroun, le paysan dans le Nord (région de l'Extrême-Nord) qui laboure manuellement sous un soleil de plomb pour assurer la récolte de mil ne « réalise » pas sa personne ; il subit la contrainte de la faim et du climat. Son geste est dicté par l'urgence vitale.

\textbf{B. L'aliénation structurelle dans le salariat capitaliste : l'expropriation du geste.}
\textit{Affirmation :} Dans la société industrielle, la contrainte devient sociale : le travailleur vend sa force de travail et perd la maîtrise de son activité.
\textit{Raisonnement / Référence :} \textbf{Karl Marx}, dans les \emph{Manuscrits de 1844} et \emph{Le Capital}, théorise l'\cle{aliénation} quadruple : (1) aliénation du produit (l'objet appartient au capitaliste), (2) aliénation dans l'acte de production (le travail est extérieur à l'ouvrier, « mortification »), (3) aliénation de l'espèce (l'homme réduit à fonction animale), (4) aliénation de l'homme par l'homme (rapport de domination). La machine devient « pouvoir intellectuel » face à l'ouvrier réduit à « appendice » (Manufacture).
\textit{Exemple :} Dans les usines de transformation du cacao ou du coton à Douala ou Garoua, l'ouvrier à la chaîne répète un geste fragmenté, cadencé par la machine, sans voir le produit fini. Il ne développe pas son intelligence ; il l'use. Le salaire n'est que le prix de cette aliénation temporaire.

\textit{Mini-conclusion / Transition :} Ainsi, historiquement et structurellement, le travail apparaît d'abord comme une \cle{servitude} — naturelle ou sociale — où l'homme se nie pour survivre. Mais cette analyse, si elle décrit une réalité massive, ne dit pas le dernier mot : le travail porte-t-il en lui les germes de son propre dépassement ?
\end{exoresolu}

\begin{methode}[Construire et rédiger l'Antithèse (Partie II)]
\textbf{Objectif :} Ne pas dire « le contraire » de la thèse, mais \emph{nuancer}, \emph{complexifier}, \emph{déplacer le regard}. Montrer que la thèse ignore une dimension essentielle.
\begin{enumerate}
    \item \textbf{Changer de niveau d'analyse :} Passer du structurel/économique à l'anthropologique/existentiel, ou du collectif à l'individuel, ou du court terme au long terme.
    \item \textbf{Formuler la thèse adverse :} « Cependant, le travail est aussi... », « À l'inverse, on peut soutenir que... ».
    \item \textbf{Structurer en 2-3 sous-parties (méthode ARE identique) :} Nouveaux arguments, nouveaux auteurs (Hegel, Arendt, Ricœur, Durkheim), nouveaux exemples.
    \item \textbf{Opérer le \cle{tournant dialectique} :} Montrer \emph{pourquoi} la thèse est insuffisante. « La thèse oublie que... », « Elle néglige la dimension... ».
    \item \textbf{Préparer la synthèse :} La dernière sous-partie de l'antithèse doit faire entrevoir la limite de cette seconde vue, appelant un troisième temps.
\end{enumerate}
\textbf{Piège à éviter :} L'antithèse n'est pas une simple liste d'arguments « pour » après des arguments « contre ». C'est une \emph{réponse critique} à la thèse.
\end{methode}

\begin{exoresolu}[Rédaction de l'Antithèse : « Le travail comme réalisation de soi et lien social » (Partie II)]
\textbf{Énoncé :} Rédigez la deuxième partie en réponse critique à la thèse marxiste/aristotélicienne.
\tcblower
\textbf{Solution détaillée :}

\textbf{II. Le travail : médiation nécessaire de l'humanisation et de la reconnaissance}

\textbf{A. Le travail comme formation de l'esprit et appropriation du monde (Hegel).}
\textit{Affirmation :} Loin de n'être que perte de soi, le travail est le processus par lequel la conscience se objectifie, se reconnaît et devient libre.
\textit{Raisonnement / Référence :} Dans la \emph{Phénoménologie de l'Esprit} (dialectique du Maître et de l'Esclave), \textbf{Hegel} montre que l'esclave, \emph{par le travail}, transforme la nature, y imprime sa forme, et se voit dans l'objet. « C'est en travaillant que l'homme se fait homme. » Le travail éduque la patience, la discipline, la rationalité technique. Il passe de la \cle{négativité} (destruction de la nature brute) à la \cle{positivité} (création de l'œuvre).
\textit{Exemple :} L'artisan menuisier à Yaoundé (quartier Mvog-Ada) qui conçoit, choisi le bois, assemble, ponce, vernit une chaise. Dans l'objet fini, il retrouve son intention, son savoir-faire, sa signature. Il s'est \emph{extériorisé} pour se \emph{re-trouver} enrichi.

\textbf{B. Le travail comme lien social et source de reconnaissance (Ricœur / Honneth / Durkheim).}
\textit{Affirmation :} Le travail n'est pas seulement rapport à la nature, c'est un rapport social médiatisé par l'échange et la reconnaissance.
\textit{Raisonnement / Référence :} \textbf{Paul Ricœur} (\emph{Philosophie de la volonté}) distingue le travail comme \cle{puissance d'agir} : « Faire » c'est « pouvoir ». \textbf{Axel Honneth} (\emph{La lutte pour la reconnaissance}) identifie la sphère du travail comme lieu de la \cle{reconnaissance sociale} (estime de soi par la contribution utile). \textbf{Durkheim} (\emph{De la division du travail social}) : la solidarité organique naît de l'interdépendance des métiers.
\textit{Exemple :} Les \cle{Groupements d'Initiative Commune (GIC)} agricoles au Cameroun (ex. groupements de femmes productrices de njansang ou de beurre de karité dans l'Adamaoua). Le travail collectif génère revenu, autonomie financière, statut social reconnu, entraide. Le travail « réalise » la personne \emph{dans} la communauté.

\textit{Mini-conclusion / Transition :} Le travail apparaît ainsi comme le grand \cle{médiateur} de l'humanité : il humanise la nature et socialise l'individu. Mais cette potentialité positive n'est pas automatique ; elle dépend des \emph{conditions concrètes} d'exercice (division du travail, propriété, technologie). C'est pourquoi il faut dépasser l'alternative contrainte/réalisation pour interroger les conditions de possibilité d'un travail émancipateur.
\end{exoresolu}

\begin{methode}[Réaliser la Synthèse et la Conclusion (Partie III + Conclusion)]
\textbf{Objectif :} Dépasser l'opposition Thèse/Antithèse par une \cle{hiérarchisation} ou une \cle{conditionnalité}, pas un simple résumé.
\begin{enumerate}
    \item \textbf{Reprendre le fil :} « Au terme de cette analyse, il apparaît que... » (Constat de la tension).
    \item \textbf{Opérer le dépassement (Aufhebung) :} Conserver la vérité des deux parties, lever leur contradiction.
        \begin{itemize}
            \item \emph{Soit par condition :} « Le travail réalise \textbf{si et seulement si}... » (conditions politiques, techniques, éducatives).
            \item \emph{Soit par distinction de niveaux :} « Au niveau existentiel... mais au niveau structurel... »
            \item \emph{Soit par historicité :} « C'est une tâche historique de transformer le travail... »
        \end{itemize}
    \item \textbf{Structurer en 2 sous-parties :} 1. Le diagnostic lucide (limites des deux vues). 2. La perspective constructive (conditions de la liberté réelle).
    \item \textbf{Rédiger la Conclusion formelle (3 temps) :}
        \begin{enumerate}
            \item \textbf{Réponse directe à la problématique} (une phrase nette).
            \item \textbf{Résumé synthétique du parcours} (Thèse $\rightarrow$ Antithèse $\rightarrow$ Synthèse).
            \item \textbf{Ouverture} : Question nouvelle, prolongement vers un autre thème du programme (ex. : Technique, État, Morale, Religion), ou perspective actuelle (IA, transition écologique). \emph{Pas de banalité.}
        \end{enumerate}
\end{enumerate}
\textbf{Interdit :} « En conclusion, on peut dire que... », introduction de \emph{nouveaux} arguments, simple copier-coller de l'annonce du plan.
\end{methode}

\begin{exoresolu}[Synthèse et Conclusion pour le sujet sur le Travail]
\textbf{Énoncé :} Rédigez la Partie III (Synthèse) et la Conclusion générale.
\tcblower
\textbf{Solution détaillée :}

\textbf{III. Vers un travail émancipateur : la condition politique et technique de la réalisation}

\textbf{A. Le diagnostic : la réalisation n'est pas un don, c'est une conquête.}
\textit{Affirmation :} La thèse (aliénation) et l'antithèse (réalisation potentielle) ne s'excluent pas ; elles décrivent le travail sous des rapports de production différents. La contradiction est \emph{historique}, pas logique.
\textit{Raisonnement / Référence :} \textbf{Marx} lui-même (\emph{L'Idéologie allemande}) précise que le communisme n'abolit pas le travail, mais le \cle{caractère aliéné} du travail : « Le travail devient non seulement un moyen de vie, mais le premier besoin de vie. » \textbf{Hannah Arendt} (\emph{La Condition de l'homme moderne}) distingue \emph{labor} (reproduction biologique, cyclique, contrainte), \emph{work} (fabrication, monde durable, homo faber) et \emph{action} (politique, liberté). La synthèse vise à réduire la part du \emph{labor} (pénibilité) pour étendre le \emph{work} (création) et libérer du temps pour l'\emph{action} (citoyenneté).
\textit{Exemple :} L'automatisation des tâches pénibles (port de charges, tri manuel) dans les ports de Kribi ou Douala \emph{peut} libérer l'humain pour des tâches de contrôle, maintenance, logistique (work), \emph{à condition} que le gain de productivité soit partagé (réduction temps de travail, formation) et non confisqué (licenciement, intensification).

\textbf{B. Les conditions concrètes d'un travail libérateur : éducation, droit, démocratie industrielle.}
\textit{Affirmation :} La transformation de la contrainte en réalisation suppose trois leviers indissociables.
\textit{Raisonnement / Référence :}
\begin{itemize}
    \item \textbf{Éducation / Formation tout au long de la vie :} Maîtriser la technique pour ne pas en être l'esclave (Dewey, \emph{Démocratie et éducation}). Au Cameroun : réforme de l'enseignement technique (FIP, insertion professionnelle) pour passer de l'exécution à la conception.
    \item \textbf{Droit du travail / Protection sociale :} Code du travail camerounais (Loi 92/007), conventions collectives, sécurité sociale (CNPS). Garantir le « travail décent » (OIT) : salaire vital, hygiène, sécurité, liberté syndicale.
    \item \textbf{Démocratie dans l'entreprise / Gouvernance partagée :} Participation aux décisions (comités d'hygiène, sécurité, conditions de travail ; coopératives SCOP). \textbf{Simone Weil} (\emph{La Condition ouvrière}) : « Il faut que l'ouvrier puisse penser son travail. »
\end{itemize}
\textit{Exemple :} Les mutuelles de santé et caisses de retraite gérées paritairement (employeurs/salariés) au Cameroun ; expériences de cogestion dans certaines PME du secteur bois/bois-meuble.

\textbf{CONCLUSION GÉNÉRALE}

\textit{Réponse directe :} « En définitive, le travail n'est ni par essence une malédiction, ni par nature une bénédiction : il est un \textbf{champ de forces} où se joue la condition humaine. Il est contrainte tant que l'homme en est l'instrument ; il devient réalisation lorsqu'il en redevient le sujet et le maître. »

\textit{Résumé du parcours :} « Nous avons vu que la contrainte (I), qu'elle soit naturelle ou capitaliste, structure encore massivement l'expérience du travail. Nous avons montré (II) que le travail porte pourtant la virtualité d'une humanisation par l'appropriation technique et la reconnaissance sociale. Nous concluons (III) que l'actualisation de cette virtualité n'est pas automatique : elle exige une \cle{transformation politique} des conditions de production — éducation, droits, démocratie industrielle — pour que le travail serve l'homme et non l'homme le travail. »

\textit{Ouverture :} « Cette exigence trouve aujourd'hui un nouveau terrain d'épreuve avec l'essor de l'\cle{Intelligence Artificielle} et de la robotique : la délégation des tâches cognitives aux algorithmes rendra-t-elle le travail humain superflu (revenu universel ?) ou redéfinira-t-elle l'humain vers le care, la création, le lien ? La question du travail rejoint ainsi celle de la \cle{technique} et de la \cle{justice sociale} au cœur de la cité. »
\end{exoresolu}

\begin{methode}[Relire et corriger sa copie (Qualité de la langue et cohérence)]
\textbf{Objectif :} Gagner les points de « Qualité de la rédaction » (souvent 4 à 6 pts sur 20) et éviter les contresens.
\begin{enumerate}
    \item \textbf{Relecture 1 : Cohérence logique (10 min).} Lire \emph{seulement} les premières phrases de chaque paragraphe (les affirmations). Le fil conducteur est-il visible ? La problématique est-elle traitée ? Les transitions font-elles sens ?
    \item \textbf{Relecture 2 : Exigence philosophique (10 min).} Chaque nom d'auteur est-il \emph{utilisé} (son concept explique l'argument) ou juste \emph{cité} ? Les définitions sont-elles respectées tout au long ? Pas de glissement de sens (ex. confondre liberté et licence).
    \item \textbf{Relecture 3 : Langue et présentation (5 min).}
        \begin{itemize}
            \item Orthographe des noms propres (Kant, Hegel, Marx, Ricœur, Arendt, Ellul, Heidegger, Descartes, Pascal, Augustin, etc.).
            \item Ponctuation : phrases courtes, pas de « et » en chaîne. Paragraphes sautés (ligne blanche) pour chaque sous-partie.
            \item Connecteurs logiques variés (\emph{cependant, dès lors, par conséquent, en outre, en revanche, corollaire}).
            \item Pas de familier (\emph{« truc, bidule, genre, grave »}), pas de « je pense que » (impersonnel : « il apparaît que », « l'analyse montre »).
        \end{itemize}
    \item \textbf{Format :} Numérotation romaine (I, II, III) pour les parties, chiffres arabes (A, B ou 1, 2) pour sous-parties. Titre de partie centré ou en marge.
\end{enumerate}
\textbf{Astuce :} Écrire au brouillon les \emph{phrases d'amorce} de chaque paragraphe + noms d'auteurs + exemples. La rédaction finale n'est que la mise en forme.
\end{methode}

\begin{exoresolu}[Correction d'une copie type : détection d'erreurs]
\textbf{Énoncé :} Voici un extrait de copie d'élève (Partie I sur « La technique libère-t-elle l'homme ? »). Identifiez 5 erreurs majeures et corrigez-les.
\textit{Extrait :} « I. La technique libère l'homme. A. Elle nous donne plus de temps libre. Par exemple, la machine à laver. Avant les femmes lavaient à la main au bord de la rivière. Maintenant elles ont le temps de regarder la télé ou de faire le commerce. C'est une vraie libération. B. Elle nous soigne. Les hôpitaux ont des scanners, des IRM. On vit plus vieux. C'est grâce à la technique. Donc la technique c'est le progrès. Marx dit que le travail c'est la vie. Donc la technique aide le travail. »
\tcblower
\textbf{Solution détaillée :}

\textbf{Erreur 1 : Absence de structure ARE (Argumentation plate, sans concept).}
\textit{Correction :} Transformer en : « \textbf{A. La libération par le gain de temps : réduction de la nécessité naturelle.} \textit{Affirmation :} La technique libère du temps contraint par les besoins vitaux. \textit{Référence :} Aristote (\emph{Politique}) distingue la \cle{nécessité} (vie biologique) et le \cle{loisir} (\emph{scholê}, condition de la politique et de la philosophie). \textit{Exemple :} L'accès à l'eau courante et aux machines (lave-linge, broyeur) dans les ménages urbains de Yaoundé libère les femmes de la corvée de l'eau et du pilage, ouvrant du temps pour l'éducation des enfants ou l'activité génératrice de revenus. »

\textbf{Erreur 2 : Exemple anecdotique / stéréotypé (« regarder la télé »).}
\textit{Correction :} Remplacer par un exemple signifiant : « activité associative, formation, engagement politique, repos réel ». La libération vise l'\cle{autonomie}, pas la consommation passive.

\textbf{Erreur 3 : Confusion Technique / Progrès / Science (Amalgame).}
\textit{Correction :} Distinguer : La technique (moyens), la science (savoirs), le progrès (valeur). L'argument médical relève de la \cle{technoscience}. Préciser : « L'application technique des savoirs médicaux (imagerie, chirurgie) augmente l'espérance de vie, condition biologique de l'exercice de la liberté. »

\textbf{Erreur 4 : Citation hors sujet / Mal comprise (Marx).}
\textit{Correction :} Marx (\emph{Le Capital}) dit que le travail est « le premier besoin de vie » \emph{dans le communisme}, mais qu'il est « aliéné » \emph{sous le capitalisme}. La technique capitaliste (machinisme) est pour lui un moyen d'exploitation. Ne pas faire dire à Marx que « la technique aide le travail » sans nuance. Mieux : « Selon Marx, le développement des forces productives (technique) crée les \emph{conditions matérielles} de l'émancipation, mais sa mise en œuvre capitaliste les retourne en chaînes. »

\textbf{Erreur 5 : Conclusion de partie manquante / Transition absente.}
\textit{Correction :} Ajouter : « Ainsi, la technique apparaît comme un \cle{moyen} d'élargir le champ des possibles (temps, santé, puissance). Mais ce pouvoir technique devient-il automatiquement liberté ? Rien n'est moins sûr : le temps libéré peut être capté par la consommation, la santé technique peut devenir contrôle biométrique. C'est ce que nous examinerons en Antithèse. »

\textbf{Bilan :} La copie corrigée gagne en rigueur conceptuelle, en références exactes, en exemples pertinents (camerounais) et en logique dialectique.
\end{exoresolu}

\begin{methode}[Gérer son temps à l'épreuve (3h / 4h selon coefficient)]
\textbf{Objectif :} Optimiser le temps imparti pour rendre une copie complète, propre, argumentée.
\begin{enumerate}
    \item \textbf{0h00 - 0h45 : Analyse \& Brouillon structuré.}
        \begin{itemize}
            \item 15 min : Lecture, cerclage, définitions, problématique (écrite en haut de la page de brouillon).
            \item 30 min : Plan détaillé (I, A, B / II, A, B / III, A, B) + Auteurs/Concepts/Exemples pour CHAQUE sous-partie. \emph{Ne pas rédiger les phrases au brouillon, seulement les mots-clés.}
        \end{itemize}
    \item \textbf{0h45 - 2h15 : Rédaction au propre (1h30).}
        \begin{itemize}
            \item Intro (15 min) : Soignée, au net direct ou micro-brouillon marge.
            \item Partie I (25 min) : 2 sous-parties ARE.
            \item Partie II (25 min) : 2 sous-parties ARE.
            \item Partie III (20 min) : Synthèse serrée.
            \item Conclusion (10 min) : 3 temps nets.
        \end{itemize}
    \item \textbf{2h15 - 2h55 : Relecture active (40 min).} Appliquer la \textbf{Méthode 6} (3 relectures ciblées).
    \item \textbf{2h55 - 3h00 : Mise en page finale.} Numérotation, saut de lignes, nom/numéro sur chaque feuille, encre bleue/noire, pas de correcteur (rayer proprement).
\end{enumerate}
\textbf{Règle d'or :} \textbf{Jamais} rendre une copie sans Conclusion. Mieux vaut une synthèse courte qu'une thèse développée sans fin. Surveillez l'horloge à 2h15 : si pas entamé la Conclusion $\rightarrow$ stoppez la rédaction, écrivez la conclusion.
\end{methode}

\begin{exoresolu}[Simulation de gestion du temps : Sujet « L'État a-t-il pour seul but d'assurer la sécurité ? »]
\textbf{Énoncé :} Établissez le rétro-planning précis (minute par minute) pour une épreuve de 3h, en indiquant ce qui est écrit sur le brouillon vs au propre à chaque créneau.
\tcblower
\textbf{Solution détaillée :}

\begin{center}
\begin{tabularx}{\linewidth}{|c|l|X|}\hline
\textbf{Horaire} & \textbf{Support} & \textbf{Action précise} \\ \hline
0h00 - 0h10 & Brouillon & Lecture active. Cerclage : \cle{État}, \cle{seul but}, \cle{sécurité}. Définitions rapides : État (Monopole violence légitime - Weber), Sécurité (protection personnes/biens), But (fin, téléologie), Seul (exclusivité ?). \\ \hline
0h10 - 0h25 & Brouillon & Problématique : « L'État, en garantissant la sécurité (condition nécessaire), ne risque-t-il pas de réduire la politique à la police, oubliant la justice, la liberté, le bien commun qui sont sa \emph{fin} propre ? » \\ \hline
0h25 - 0h45 & Brouillon & \textbf{Plan détaillé (mots-clés only)} : \\
& & \textbf{I. Sécurité : condition \emph{sine qua non} (Fonction police).} A. Hobbes/Léviathan : sortir de l'état de nature (guerre de tous contre tous). Sécurité = vie. B. Fonctions régaliennes (Police, Justice, Armée, Frontières). Ex : Cameroun - lutte Boko Haram, sécurisation élections. \\
& & \textbf{II. Sécurité insuffisante : l'État pour la Justice et la Liberté (Fonction politique).} A. Locke/Rousseau : Contrat social pour \emph{conserver} liberté/propriété, pas juste vie. Droit, Loi générale (Rousseau). B. État social / Providence : Éducation, Santé, Infrastructures (Keynes, Bismarck). Ex : Cameroun - éducation de base, routes, CAN 2019 legs infrastructurels. \\
& & \textbf{III. La Sécurité \emph{au service} de la Justice (Synthèse normative).} A. Sécurité sans droit = Tyrannie (Montesquieu, Arendt). Légalité $\neq$ Légitimité. B. Démocratie / État de droit : Contrôle parlementaire, Droits de l'homme, Séparation pouvoirs. Sécurité \emph{moyen}, Justice \emph{fin}. Ex : Transition démocratique, décentralisation. \\ \hline
0h45 - 1h00 & Propre & \textbf{Introduction} (Amorce : actualité sécuritaire Cameroun / Définitions / Problématique / Plan). \\ \hline
1h00 - 1h25 & Propre & \textbf{Partie I} (2 paragraphes ARE). \\ \hline
1h25 - 1h50 & Propre & \textbf{Partie II} (2 paragraphes ARE). \\ \hline
1h50 - 2h10 & Propre & \textbf{Partie III} (2 paragraphes ARE). \\ \hline
2h10 - 2h20 & Propre & \textbf{Conclusion} (Réponse / Résumé / Ouverture : Sécurité humaine / ODD ONU). \\ \hline
2h20 - 2h55 & Propre & \textbf{Relecture 1 (Logique) :} 10 min. \textbf{Relecture 2 (Auteurs/Concepts) :} 15 min. \textbf{Relecture 3 (Langue/Présentation) :} 10 min. \\ \hline
2h55 - 3h00 & Propre & Vérification numérotation, nom, signature, encre. Soumission. \\ \hline
\end{tabularx}
\end{center}

\textbf{Point de vigilance :} Si à 1h45 la Partie II n'est pas finie $\rightarrow$ accélérer, synthétiser les exemples, passer à la Partie III. La structure complète (I, II, III, Conclu) vaut plus qu'une Partie I parfaite et pas de suite.
\end{exoresolu}

\begin{attention}[Les 5 erreurs qui coûtent des points au Bac (Probatoire / Baccalauréat Technique)]
\begin{enumerate}
    \item \textbf{Sujet « traité » au lieu d'être « problématisé » :} Rédiger un exposé de connaissances (cours récité : « La technique c'est... Marx dit... Hegel dit... ») sans \emph{répondre} à la question précise du sujet. \emph{Sanction :} Note plafond 6-8/20 (Hors sujet déguisé). \textbf{Parade :} Problématique obligatoire + Annonce de plan qui y répond + Transitions qui y reviennent.
    \item \textbf{« Name-dropping » stérile (Citation décorative) :} Citer Kant, Marx, Descartes sans expliquer \emph{en quoi} leur concept éclaire l'argument. Ex : « Comme le dit Kant, la liberté est autonomie. » (Stop). \emph{Sanction :} Points « Argumentation » et « Culture philosophique » perdus. \textbf{Parade :} Méthode ARE : Auteur $\rightarrow$ Concept précis $\rightarrow$ Lien logique avec l'argument.
    \item \textbf{Confusion Thèse / Antithèse / Synthèse (Plan « Pour / Contre / Mélange ») :} Faire une liste d'arguments « pour » (I), « contre » (II), « les deux » (III). L'antithèse doit \emph{critiquer} la thèse (changer de niveau), la synthèse doit \emph{dépasser} (condition, distinction, historicité). \emph{Sanction :} Note « Structure dialectique » faible. \textbf{Parade :} Verbes de liaison : « Cependant », « Plus profondément », « En définitive, c'est à condition que... ».
    \item \textbf{Exemples « de la vie courante » non philosophiques :} « Mon oncle a acheté une moto », « J'ai vu à la télé », « Au quartier on dit que... ». L'exemple doit être \emph{typique}, \emph{signifiant}, \emph{géographiquement situé} (Cameroun) et \emph{lié au concept}. \emph{Sanction :} Perte de points « Application concrète ». \textbf{Parade :} Banque d'exemples préparés : GIC, Code du travail, CAN, Boko Haram, Décentralisation, FIP, Agriculture, Numérique, Diaspora.
    \item \textbf{Copie illisible / incomplète / sans conclusion :} Ratures multiples, pas de saut de paragraphes, pas de numérotation I/II/III, conclusion oubliée par manque de temps. \emph{Sanction :} Points « Présentation » (souvent 2-4 pts) + impression générale négative du correcteur + absence de réponse finale. \textbf{Parade :} Gestion du temps rigoureuse (Méthode 7), écriture soignée, \textbf{Conclusion sacro-sainte à 2h10 max}.
\end{enumerate}
\end{attention}

\newpage

\coursec{Exercices}

\courssub{Je vérifie mes connaissances}

\begin{exercice}[QCM : Structure de la dissertation]
Parmi les propositions suivantes, laquelle correspond à l'ordre chronologique et logique des parties d'une dissertation philosophique standard ?
\begin{enumerate}
\item Thèse $\rightarrow$ Introduction $\rightarrow$ Antithèse $\rightarrow$ Synthèse $\rightarrow$ Conclusion
\item Introduction $\rightarrow$ Thèse $\rightarrow$ Antithèse $\rightarrow$ Synthèse $\rightarrow$ Conclusion
\item Introduction $\rightarrow$ Antithèse $\rightarrow$ Thèse $\rightarrow$ Conclusion $\rightarrow$ Synthèse
\item Thèse $\rightarrow$ Antithèse $\rightarrow$ Introduction $\rightarrow$ Synthèse $\rightarrow$ Conclusion
\end{enumerate}
\end{exercice}

\begin{exercice}[Vrai / Faux justifié : Rôle de l'introduction]
Indiquez si les affirmations suivantes sont vraies ou fausses. Justifiez brièvement votre réponse en vous référant aux attendus du programme de Première technique.
\begin{enumerate}
\item L'accroche (ou amorce) est une étape facultative que l'on peut supprimer pour gagner du temps.
\item La problématique doit être une question fermée (réponse par oui ou non) pour être efficace.
\item L'annonce du plan doit explicitement nommer les parties « Thèse », « Antithèse » et « Synthèse ».
\item La définition des termes du sujet se fait obligatoirement avant l'accroche.
\end{enumerate}
\end{exercice}

\begin{exercice}[Définitions et notions clés]
Donnez une définition précise et concise (2 à 3 lignes maximum) pour chacun des concepts suivants, dans le contexte de la dissertation philosophique :
\begin{enumerate}
\item La \cle{problématique}.
\item La \cle{thèse} (comme partie du développement).
\item La \cle{synthèse} (différencier de la simple conclusion).
\item Le \cle{connecteur logique} (donner 3 exemples classés par fonction : opposition, conséquence, concession).
\end{enumerate}
\end{exercice}

\begin{exercice}[Repérage d'erreurs courantes]
Lisez l'extrait d'introduction suivant pour le sujet : \og La technique libère-t-elle l'homme ? \fg{} \\
\og De tout temps, l'homme a fabriqué des outils. Aujourd'hui, avec les smartphones et l'intelligence artificielle, la technique est partout. Nous allons voir si la technique libère l'homme. D'abord, nous dirons que oui, la technique libère l'homme parce qu'elle facilite le travail. Ensuite, nous dirons que non, elle l'asservit car on devient dépendant. Enfin, nous conclurons. \fg{}

Identifiez \textbf{quatre} défauts majeurs dans cette introduction par rapport aux exigences méthodologiques (accroche, définition, problématisation, annonce de plan). Pour chaque défaut, proposez une correction concrète.
\end{exercice}

\courssub{Je m'entraîne}

\begin{exercice}[Rédaction d'une introduction complète]
Sujet : \og \textbf{La culture dénature-t-elle l'homme ?} \fg{}

Rédigez une introduction complète et soignée (environ 300--350 mots) respectant scrupuleusement les quatre étapes canoniques : \cle{accroche}, \cle{définition des termes} (culture, nature, dénaturer), \cle{problématique} (unique, ouverte, issue de la tension entre les définitions), \cle{annonce du plan} (Thèse / Antithèse / Synthèse avec indication des arguments directeurs).
\end{exercice}

\begin{exercice}[Construction du plan détaillé : Thèse / Antithèse]
Sujet : \og \textbf{L'État est-il un mal nécessaire ?} \fg{}

À partir de ce sujet, construisez le \textbf{plan détaillé} des deux premières parties du développement (Thèse et Antithèse uniquement).
Pour chaque partie :
\begin{itemize}
\item Formulez l'\cle{idée directrice} (la réponse partielle au sujet).
\item Listez \textbf{3 arguments distincts} (philosophiques, historiques ou sociologiques).
\item Associez à chaque argument un \textbf{exemple concret} (référence auteur, fait historique, situation camerounaise ou universelle).
\item Indiquez les \cle{connecteurs logiques} assurant la transition entre les arguments.
\end{itemize}
Présentez le tout sous forme de tableau ou de liste structurée.
\end{exercice}

\begin{exercice}[Rédaction de l'antithèse argumentée]
On donne la thèse suivante pour le sujet \og \textbf{La liberté consiste-t-elle à faire ce que l'on veut ?} \fg{} :

\emph{Thèse : La liberté, c'est l'autonomie. Faire ce que l'on veut, c'est obéir à sa propre raison et non à ses passions ou à des contraintes extérieures. Ainsi, le libre arbitre suppose la maîtrise de soi (ex. : le fumeur qui arrête de fumer par raison est plus libre que celui qui fume \og quand il veut \fg{}).}

Rédigez l'\textbf{antithèse complète} (environ 250--300 mots) qui s'oppose dialectiquement à cette thèse.
Contraintes :
\begin{itemize}
\item Idée directrice claire (ex. : La liberté est absence de contrainte / spontanéité).
\item 3 arguments développés avec exemples.
\item Transitions fluides.
\item Ouverture vers la synthèse en fin de partie.
\end{itemize}
\end{exercice}

\begin{exercice}[Correction et amélioration d'une conclusion faible]
Voici une conclusion proposée pour le sujet \og \textbf{Le bonheur est-il affaire de chance ou de raison ?} \fg{} :

\og En conclusion, nous avons vu que le bonheur dépend de la chance (santé, naissance) mais aussi de la raison (sagesse, vertu). Donc le bonheur est les deux à la fois. Il faut être raisonnable et avoir de la chance. C'est tout. \fg{}

Cette conclusion présente au moins \textbf{trois faiblesses majeures} (absence de réponse tranchée à la problématique, pas de reprise synthétique du parcours, pas d'ouverture).
\begin{enumerate}
\item Identifiez explicitement ces trois faiblesses.
\item Réécrivez une \textbf{nouvelle conclusion} (environ 150--200 mots) rigoureuse, appliquant la méthode : \textbf{Réponse directe à la problématique} $\rightarrow$ \textbf{Résumé articulé du parcours dialectique} (thèse/antithèse/synthèse) $\rightarrow$ \textbf{Ouverture philosophique} (nouvelle question, prolongement, citation pertinente).
\end{enumerate}
\end{exercice}

\courssub{Je me prépare au Bac}

\begin{exercice}[Sujet type Bac : Dissertation complète en temps limité (3h)]
\textbf{Sujet :} \og \textbf{Le bonheur est-il affaire de chance ou de raison ?} \fg{}

Vous disposez de 3 heures pour traiter ce sujet. Les questions ci-dessous balisent les étapes obligatoires de votre copie. \textbf{Rédigez la dissertation complète} en respectant la structure Introduction / Développement (Thèse, Antithèse, Synthèse) / Conclusion.

\textbf{Étape 1 -- Analyse du sujet (brouillon, 20 min) :}
\begin{enumerate}
\item Identifiez les termes clés et définissez-les philosophiquement (bonheur, chance, raison, affaire de).
\item Mettez en évidence l'opposition / la tension (exclusivité ? articulation ?).
\item Formulez 2 problématiques possibles et choisissez la plus féconde. Justifiez votre choix.
\end{enumerate}

\textbf{Étape 2 -- Introduction (30 min) :}
Rédigez l'introduction complète (accroche, définitions, problématique choisie, annonce de plan).

\textbf{Étape 3 -- Développement (1h45) :}
\begin{enumerate}
\item \textbf{Thèse :} Le bonheur relève de la raison (maîtrise, vertu, sagesse stoïcienne, kantienne, aristotélicienne). 3 arguments + exemples.
\item \textbf{Antithèse :} Le bonheur dépend de la chance / conditions extérieures (santé, richesse, naissance, \emph{tyché} grecque, condition sociale). 3 arguments + exemples.
\item \textbf{Synthèse :} Dépasser l'alternative. Le bonheur comme \cle{activité conforme à la vertu dans une vie complète} (Aristote) ou \cle{bonheur comme espérance/construction} (Alain, Ricœur) : la raison \emph{domestique} la chance mais ne la supprime pas. Articulation nécessaire.
\end{enumerate}

\textbf{Étape 4 -- Conclusion (15 min) :}
Réponse finale nuancée, résumé dialectique, ouverture.

\textbf{Étape 5 -- Relecture (10 min) :}
Vérification orthographe, connecteurs, respect du plan annoncé.
\end{exercice}

\begin{exercice}[Analyse et correction d'introductions : Repérage de hors-sujet et défauts]
Trois élèves de Première technique ont rédigé l'introduction du sujet : \og \textbf{La justice est-elle seulement le respect des lois ?} \fg{}

\textbf{Introduction A :}
\og Depuis l'Antiquité, les hommes cherchent la justice. Au Cameroun, le Code pénal définit les infractions. Respecter la loi, c'est être juste. Nous allons voir si la justice est seulement le respect des lois. D'abord, oui, car la loi assure l'ordre. Ensuite, non, car il y a des lois injustes. Enfin, la justice est l'équité. \fg{}

\textbf{Introduction B :}
\og \emph{« La justice est la première vertu des institutions sociales »,} écrit John Rawls. Le terme \og justice \fg{} vient du latin \emph{justitia} (droiture, équité). La \og loi \fg{} est une règle obligatoire édictée par l'autorité légitime. Mais une loi peut-elle être injuste ? Si la justice se réduit à la légalité, comment juger les régimes totalitaires légaux ? \textbf{La justice se réduit-elle à la conformité formelle à la loi ou exige-t-elle une conformité matérielle à l'équité ?} Nous montrerons d'abord que le respect de la loi est la condition formelle de la justice (Thèse), puis que la légalité ne suffit pas sans légitimité morale (Antithèse), pour conclure que la justice est l'harmonie entre légalité et légitimité (Synthèse). \fg{}

\textbf{Introduction C :}
\og La justice, c'est quand tout le monde est traité pareil. Les lois sont faites par les députés. Est-ce que la justice c'est seulement respecter les lois ? Mon plan : 1) La loi protège les faibles. 2) Parfois la loi est méchante. 3) La vraie justice c'est l'amour. \fg{}

\begin{enumerate}
\item Pour chaque introduction (A, B, C), identifiez \textbf{les points forts} et \textbf{les défauts majeurs} (hors-sujet, absence de problématique, définitions manquantes, plan flou, ton familier, etc.).
\item Classez-les de la meilleure à la moins bonne en justifiant.
\item Réécrivez l'introduction la plus faible (C) en une introduction \textbf{modèle} (niveau Bac), en conservant l'idée intuitive de l'élève (l'égalité, la protection des faibles) mais en la formalisant philosophiquement.
\end{enumerate}
\end{exercice}

\begin{exercice}[Sujet d'invention / Application contextuelle : Dissertation à partir d'un dossier]
\textbf{Dossier :}
\begin{itemize}
\item \textbf{Doc 1 (Extrait \emph{Critique de la raison pure}, Kant) :} \og Agis uniquement d'après la maxime qui fait que tu puisses vouloir en même temps qu'elle devienne une loi universelle. \fg{}
\item \textbf{Doc 2 (Article de presse, Yaoundé, 2023) :} \og Corruption : 150 milliards FCFA détournés par an selon la CONAC. L'impunité reste la règle pour les hauts fonctionnaires. \fg{}
\item \textbf{Doc 3 (Proverbe Bamiléké) :} \og \emph{« La loi est comme la toile d'araignée : elle retient les petits insectes mais les gros la traversent. »} \fg{}
\end{itemize}

\textbf{Sujet :} \og \textbf{Le respect de la loi suffit-il à fonder la justice ?} \fg{}

En vous appuyant \textbf{obligatoirement} sur les trois documents du dossier (citations, données, proverbe) et sur vos connaissances philosophiques (Thèse / Antithèse / Synthèse), rédigez une dissertation structurée.

\textbf{Consignes spécifiques :}
\begin{enumerate}
\item L'introduction doit exploiter le proverbe (Doc 3) comme \cle{accroche}.
\item La \textbf{Thèse} doit s'appuyer sur Kant (Doc 1) : la loi comme forme universelle de la justice (légalisme formel).
\item L'\textbf{Antithèse} doit s'appuyer sur le Doc 2 (réalité camerounaise) et le Doc 3 : dissociation entre légalité et justice réelle (injustice structurelle, impunité).
\item La \textbf{Synthèse} doit proposer une articulation : la justice comme \cle{équité} (Rawls) ou \cle{droit naturel} dépassant le positivisme juridique, nécessitant des institutions indépendantes (séparation des pouvoirs, CONAC effective).
\item La conclusion doit ouvrir sur le rôle du \cle{citoyen} et de la \cle{société civile} dans l'effectivité de la justice.
\end{enumerate}
\end{exercice}

\newpage

\coursec{Corrigés des exercices}

\coursec{La dissertation philosophique}

\begin{corrige}[Exercice 1 — QCM : Structure de la dissertation]
\textbf{Réponse correcte : proposition 2.}

L'ordre chronologique et logique d'une dissertation philosophique est :
\[ \text{Introduction} \rightarrow \text{Thèse} \rightarrow \text{Antithèse} \rightarrow \text{Synthèse} \rightarrow \text{Conclusion} \]

\textbf{Justification :}
\begin{itemize}
\item L'\cle{introduction} vient nécessairement en premier : elle présente le sujet, définit les termes, pose la problématique et annonce le plan. Elle ne peut donc pas suivre la thèse (ce qui élimine les propositions 1 et 4).
\item Le \cle{développement} suit l'ordre dialectique classique : on expose d'abord une position (\cle{thèse}), puis on la confronte à la position opposée (\cle{antithèse}), puis on dépasse l'opposition (\cle{synthèse}). La proposition 3 est fausse car elle inverse thèse et antithèse et place la conclusion avant la synthèse, ce qui est absurde : la synthèse appartient au développement, la conclusion clôt le devoir.
\item La \cle{conclusion} vient en dernier : elle répond à la problématique et ouvre éventuellement sur une nouvelle question.
\end{itemize}

\begin{attention}[Point de vigilance]
Au brouillon, on construit souvent le plan \emph{avant} de rédiger l'introduction ; mais sur la copie, l'introduction est toujours rédigée et placée \emph{en premier}. Ne pas confondre ordre de \emph{préparation} et ordre de \emph{rédaction}.
\end{attention}
\end{corrige}

\begin{corrige}[Exercice 2 — Vrai / Faux justifié : Rôle de l'introduction]
\begin{enumerate}
\item \textbf{Faux.} L'\cle{accroche} (ou amorce) est une étape obligatoire de l'introduction philosophique. Elle situe le sujet dans une perspective générale et suscite l'intérêt du lecteur. Une introduction qui commence brutalement par la définition ou la problématique est jugée incomplète à l'examen. On peut la faire brève (une ou deux phrases), mais jamais la supprimer.

\item \textbf{Faux.} La \cle{problématique} doit être une question \emph{ouverte}, appelant une réflexion argumentée et nuancée, et non une question fermée. Une question fermée (\og La technique libère-t-elle l'homme : oui ou non ? \fg{}) ne permet pas de construire un développement dialectique ; c'est précisément parce que la réponse n'est pas évidente que le sujet mérite d'être traité. La problématique exprime une \emph{tension} entre des positions possibles.

\item \textbf{Faux.} L'annonce du plan ne doit \emph{jamais} employer les mots techniques \og thèse \fg{}, \og antithèse \fg{}, \og synthèse \fg{} de manière scolaire et mécanique (\og Dans une première partie nous ferons la thèse... \fg{}). Elle doit annoncer le \emph{contenu} des étapes (\og Nous verrons d'abord que..., nous montrerons ensuite que..., avant de dégager... \fg{}). Le plan doit être visible dans sa logique, non dans son étiquette.

\item \textbf{Faux.} L'ordre canonique est : accroche $\rightarrow$ définition des termes $\rightarrow$ problématique $\rightarrow$ annonce du plan. La définition des termes \emph{suit} l'accroche : c'est l'accroche qui amène naturellement aux termes du sujet, dont l'analyse fait naître la problématique. Inverser accroche et définitions supprime la progression de l'introduction.
\end{enumerate}

\begin{aretenir}[L'ordre canonique de l'introduction]
Accroche $\rightarrow$ Définition des termes du sujet $\rightarrow$ Problématique (question ouverte) $\rightarrow$ Annonce du plan (sans jargon méthodologique).
\end{aretenir}
\end{corrige}

\begin{corrige}[Exercice 3 — Définitions et notions clés]
\begin{enumerate}
\item \textbf{La problématique} : question philosophique \emph{ouverte}, formulée à partir de la tension révélée par l'analyse des termes du sujet. Elle oriente l'ensemble du devoir : toute la dissertation est une réponse argumentée à cette question unique.

\item \textbf{La thèse} : première partie du développement, qui expose et défend une première réponse à la problématique (idée directrice, arguments, exemples, références). Elle est ensuite mise à l'épreuve par l'antithèse.

\item \textbf{La synthèse} : troisième partie du développement, qui \emph{dépasse} l'opposition thèse/antithèse en montrant ce que chacune contient de vrai et en proposant une position plus riche. Elle se distingue de la \cle{conclusion}, qui clôt le devoir en répondant directement à la problématique et en ouvrant éventuellement sur une nouvelle question : la synthèse est une étape du raisonnement, la conclusion est un bilan.

\item \textbf{Le connecteur logique} : mot ou groupe de mots qui articule les idées entre elles et rend visible la progression du raisonnement. Exemples :
\begin{itemize}
\item \emph{Opposition} : \og mais \fg{}, \og cependant \fg{}, \og en revanche \fg{} ;
\item \emph{Conséquence} : \og donc \fg{}, \og par conséquent \fg{}, \og c'est pourquoi \fg{} ;
\item \emph{Concession} : \og certes... mais \fg{}, \og bien que \fg{}, \og il est vrai que... néanmoins \fg{}.
\end{itemize}
\end{enumerate}
\end{corrige}

\begin{corrige}[Exercice 4 — Repérage d'erreurs courantes]
\textbf{Défaut 1 : accroche banale et creuse.} \og De tout temps, l'homme a fabriqué des outils \fg{} est un lieu commun sans portée philosophique. \emph{Correction} : partir d'un constat précis ou d'une référence : \og L'homme, dit Aristote, est un animal raisonnable ; mais il est aussi, selon Bergson, un \emph{homo faber}, un fabricant d'outils qui transforme la nature et se transforme lui-même. \fg{}

\textbf{Défaut 2 : absence de définition des termes.} Ni \og technique \fg{}, ni \og libérer \fg{} ne sont définis ; or c'est de leur analyse que naît la problématique. \emph{Correction} : définir la \cle{technique} (ensemble des moyens artificiels mis en œuvre pour agir sur le monde) et \cle{libérer} (rendre autonome, soustraire à une contrainte), puis montrer la tension : l'outil qui libère d'un effort peut créer une dépendance nouvelle.

\textbf{Défaut 3 : problématique faible et simple reprise du sujet.} \og Nous allons voir si la technique libère l'homme \fg{} répète le sujet sans le problématiser. \emph{Correction} : formuler la tension : \og La technique, en augmentant la puissance humaine, libère-t-elle l'homme des contraintes, ou crée-t-elle de nouvelles servitudes plus subtiles ? \fg{}

\textbf{Défaut 4 : annonce de plan scolaire et réponses déjà tranchées.} \og Nous dirons que oui... nous dirons que non... \fg{} annonce le plan comme une liste et donne les conclusions avant de les démontrer ; de plus, \og enfin, nous conclurons \fg{} n'annonce aucune synthèse. \emph{Correction} : annoncer la \emph{logique} du parcours : \og Nous verrons d'abord que la technique affranchit l'homme du travail pénible ; nous montrerons ensuite qu'elle engendre de nouvelles dépendances ; il s'agira enfin de déterminer à quelles conditions la technique peut rester un instrument de liberté. \fg{}
\end{corrige}

\begin{corrige}[Exercice 5 — Rédaction d'une introduction complète]
\textbf{Modèle d'introduction} (sujet : \og La culture dénature-t-elle l'homme ? \fg{}) :

\medskip
\og Ce qui est naturel n'est point dans l'usage : c'est un fait que l'usage a rendu naturel \fg{}, remarquait déjà Pascal, soulignant combien l'homme est un être transformé par ses propres créations. De l'éducation reçue dans la famille aux techniques apprises à l'école ou à l'atelier, tout en lui semble marqué par la culture. Mais cette transformation est-elle une altération ?

Le terme \cle{culture} désigne l'ensemble des savoirs, des valeurs, des langues et des techniques qu'une société transmet à ses membres et par lesquels l'homme s'élève au-dessus de la simple vie biologique. La \cle{nature} renvoie, elle, à ce qui existe en l'homme indépendamment de toute éducation : ses instincts, ses besoins, son héritage biologique. Quant à \cle{dénaturer}, il signifie étymologiquement \og priver de sa nature \fg{}, c'est-à-dire corrompre ou pervertir une essence originelle. La question prend alors toute sa force : si l'homme possède une nature propre, la culture, en le façonnant, ne la détruit-elle pas ? Mais si, inversement, l'homme n'accède à son humanité que par la culture, celle-ci ne serait-elle pas sa véritable nature ?

On peut donc se demander : \textbf{la culture, en transformant l'homme, le détourne-t-elle de sa nature, ou est-elle précisément le moyen par lequel il accomplit son humanité ?}

Pour répondre à cette question, nous verrons d'abord que la culture peut apparaître comme une dénaturation, en ce qu'elle étouffe la spontanéité naturelle et corrompt l'homme, comme le soutient Rousseau. Nous montrerons ensuite, avec Lévi-Strauss et la tradition humaniste, que l'homme n'a pas de nature accomplie sans la culture, qui est sa seconde nature et sa véritable libération. Nous dégagerons enfin l'idée que la culture ne détruit pas la nature humaine mais la cultive, à condition de rester au service de l'épanouissement de l'homme et non de son aliénation.

\medskip
\textbf{Points de vigilance :} accroche référencée (Pascal) ; trois termes définis ; problématique unique et ouverte ; annonce du plan en trois étapes \emph{sans} les mots \og thèse/antithèse/synthèse \fg{} ; environ 320 mots.
\end{corrige}

\begin{corrige}[Exercice 6 — Construction du plan détaillé : Thèse / Antithèse]
Sujet : \og L'État est-il un mal nécessaire ? \fg{}

\textbf{PARTIE 1 — Thèse. Idée directrice :} l'État est un mal, mais un mal \emph{nécessaire} : il limite les libertés individuelles, cependant sans lui c'est la violence de tous contre tous.

\begin{center}
\begin{tabularx}{\linewidth}{|l|X|X|}
\hline
\rowcolor{popblueL} \textbf{Argument} & \textbf{Exemple / référence} & \textbf{Connecteur} \\
\hline
1. Sans État, règne l'état de nature, guerre de tous contre tous ; l'État garantit la sécurité. & Hobbes, \emph{Léviathan} : \og l'homme est un loup pour l'homme \fg{} ; situations d'effondrement étatique (Somalie des années 1990). & \og D'abord \fg{}, \og en effet \fg{} \\
\hline
2. L'État assure la justice et protège les faibles par la loi commune. & Les tribunaux et le Code pénal au Cameroun sanctionnent les agressions ; sans eux, la vengeance privée (vendetta). & \og Ensuite \fg{}, \og de plus \fg{} \\
\hline
3. L'État rend possible la vie en société : écoles, routes, hôpitaux. & La construction d'écoles publiques et d'hôpitaux de district au Cameroun suppose une autorité organisatrice. & \og Enfin \fg{}, \og ainsi \fg{} \\
\hline
\end{tabularx}
\end{center}

\textbf{PARTIE 2 — Antithèse. Idée directrice :} l'État n'est pas seulement un remède : il peut être lui-même un mal, oppression et violence légitime ; il faut donc le limiter, voire en rêver la disparition.

\begin{center}
\begin{tabularx}{\linewidth}{|l|X|X|}
\hline
\rowcolor{popblueL} \textbf{Argument} & \textbf{Exemple / référence} & \textbf{Connecteur} \\
\hline
1. L'État concentre la violence et peut devenir tyrannique. & Régimes totalitaires du XX\textsuperscript{e} siècle ; La Boétie, \emph{Discours de la servitude volontaire}. & \og Cependant \fg{}, \og mais \fg{} \\
\hline
2. La loi d'État peut être injuste : légalité $\neq$ justice. & L'apartheid en Afrique du Sud était \emph{légal} ; désobéissance civile de Gandhi, Martin Luther King. & \og En outre \fg{}, \og en effet \fg{} \\
\hline
3. La société peut s'organiser sans coercition étatique (anarchisme, entraide). & Kropotkine, \emph{L'Entraide} ; coopératives villageoises et tontines au Cameroun, qui fonctionnent sans intervention de l'État. & \og De surcroît \fg{}, \og par conséquent \fg{} \\
\hline
\end{tabularx}
\end{center}

\textbf{Transition vers la synthèse (à prévoir) :} \og L'État est donc à la fois remède et poison ; la question n'est plus de le subir ou de le supprimer, mais de le limiter par le droit et la séparation des pouvoirs (Montesquieu). \fg{}
\end{corrige}

\begin{corrige}[Exercice 7 — Rédaction de l'antithèse argumentée]
\textbf{Modèle d'antithèse} (sujet : \og La liberté consiste-t-elle à faire ce que l'on veut ? \fg{}) :

\medskip
Certes, définir la liberté par la maîtrise rationnelle de soi semble exigeant et noble. Mais n'est-ce pas confondre la liberté avec la discipline ? À y regarder de plus près, la liberté se reconnaît d'abord à l'\emph{absence de contrainte} : est libre celui qui peut agir selon son désir, sans obstacle extérieur ni règle imposée.

D'abord, l'expérience commune confond liberté et spontanéité. L'enfant qui joue, le voyageur qui choisit sa route au gré de l'envie, l'artiste qui crée sans modèle : tous se sentent libres précisément parce qu'ils ne calculent pas, n'obéissent à aucune règle préalable. Dire au fumeur qu'il n'est libre qu'en s'interdisant de fumer, c'est renverser le sens ordinaire du mot : la liberté serait alors une forme déguisée d'obligation.

Ensuite, la conception rationaliste de la liberté aboutit à un paradoxe redoutable : elle permet de dire qu'on peut \og forcer quelqu'un à être libre \fg{}. Rousseau lui-même, dans le \emph{Contrat social}, admet cette formule inquiétante. Si la vraie liberté consiste à obéir à la raison, alors celui qui refuse cette obéissance pourra être contraint \og pour son bien \fg{} : c'est la porte ouverte au paternalisme, voire à la tyrannie exercée au nom de la raison. L'histoire du XX\textsuperscript{e} siècle montre combien les régimes prétendant libérer les peuples \emph{malgré eux} ont produit les pires servitudes.

Enfin, faire de la raison la condition de la liberté, c'est oublier que le désir n'est pas toujours une passion aveugle : il est aussi une force de vie et d'invention. Nietzsche reproche à la morale rationaliste d'être une morale d'esclave qui mutile les instincts. L'homme qui danse, aime, ose \og parce qu'il en a envie \fg{} n'est pas nécessairement un esclave de ses passions : il affirme sa singularité.

Ainsi, la liberté comme maîtrise de soi risque de se confondre avec une servitude intérieure. Faut-il pour autant renoncer à toute idée d'autonomie ? Il reste à examiner si la véritable liberté ne consiste pas, plutôt, à savoir \emph{ce que} l'on veut et à s'en donner les moyens : ni caprice, ni règle abstraite, mais puissance d'agir éclairée.

\medskip
\textbf{Points de vigilance :} idée directrice opposée à la thèse (liberté = absence de contrainte) ; 3 arguments distincts (sens commun, paradoxe du \og forcer à être libre \fg{}, critique nietzschéenne) ; exemples précis ; connecteurs (\og certes... mais \fg{}, \og d'abord \fg{}, \og ensuite \fg{}, \og enfin \fg{}, \og ainsi \fg{}) ; dernière phrase ouvrant vers la synthèse.
\end{corrige}

\begin{corrige}[Exercice 8 — Correction et amélioration d'une conclusion faible]
\textbf{1. Les trois faiblesses majeures :}
\begin{enumerate}
\item \textbf{Pas de réponse tranchée à la problématique.} \og Le bonheur est les deux à la fois \fg{} juxtapose les deux termes sans les articuler : on ne sait pas \emph{comment} chance et raison se combinent, ni laquelle prime. La synthèse est remplacée par un compromis mou.
\item \textbf{Pas de reprise synthétique du parcours.} La conclusion résume en une phrase ce qui devrait être un rappel structuré du cheminement dialectique (thèse, antithèse, synthèse), montrant ce que chaque étape a établi.
\item \textbf{Pas d'ouverture, ton familier et phrase creuse.} \og C'est tout. \fg{} clôt brutalement le devoir, sans prolongement philosophique, et donne une impression de légèreté incompatible avec l'exercice.
\end{enumerate}

\textbf{2. Nouvelle conclusion (modèle) :}

\medskip
Le bonheur est-il affaire de chance ou de raison ? À l'issue de notre réflexion, nous pouvons répondre : il est affaire de raison \emph{d'abord}, mais d'une raison qui ne prétend pas tout contrôler. Nous avons vu, dans un premier temps, que la sagesse stoïcienne et la vertu aristotélicienne font du bonheur une conquête de la raison : celui qui maîtrise ses désirs et agit selon la vertu dispose d'un bonheur que la fortune ne peut lui ravir. Mais l'expérience contredit cette prétention : la maladie, la pauvreté, l'injustice rappellent que nul ne choisit ni sa naissance ni sa santé, et que la \emph{tyché} gouverne une part de nos vies. Fallait-il alors renoncer à articuler les deux ? Non : la raison ne supprime pas la chance, elle la \emph{domestique} — elle transforme les épreuves en occasions de grandeur et sait recevoir les biens de la fortune sans s'y asservir. Le bonheur est ainsi moins un état donné qu'une œuvre : celle d'une vie conduite avec lucidité au milieu de l'incertain. Reste alors une question plus vaste, que la philosophie politique reprend à son compte : si le bonheur dépend aussi des conditions extérieures, n'est-ce pas à la \emph{cité} elle-même de les rendre possibles pour tous ?

\medskip
\textbf{Barème indicatif :} réponse directe et nuancée (1 pt) ; rappel articulé thèse/antithèse/synthèse (2 pts) ; ouverture pertinente liée au sujet (1 pt) ; langue soutenue, sans \og je \fg{} ni formule familière (1 pt).
\end{corrige}

\begin{corrige}[Exercice 9 — Sujet type Bac : Dissertation complète]
\textbf{Éléments de correction par étapes, avec barème indicatif sur 20.}

\textbf{Étape 1 — Analyse du sujet (attendu au brouillon).}
\begin{itemize}
\item \cle{Bonheur} : état de satisfaction durable et totale, accomplissement de l'existence (à distinguer du plaisir, passager).
\item \cle{Chance} : ce qui advient indépendamment de notre volonté (la \emph{tyché} des Grecs) : naissance, santé, rencontres.
\item \cle{Raison} : faculté de juger, de calculer et de se gouverner soi-même.
\item \og Affaire de \fg{} : ce qui relève de, ce qui dépend de.
\item \textbf{Tension :} le sujet oppose deux sources exclusives du bonheur : ce qui ne dépend pas de nous (chance) et ce qui en dépend (raison).
\item \textbf{Problématiques possibles :} (a) \og Le bonheur dépend-il de nous ou du hasard ? \fg{} ; (b) \og La raison peut-elle suffire à faire le bonheur malgré la fortune ? \fg{}. La (b) est plus féconde car elle permet une synthèse (la raison \emph{face à} la chance) plutôt qu'un simple partage.
\end{itemize}

\textbf{Étape 2 — Introduction (modèle, 4 pts).} Accroche : \og Tous les hommes recherchent le bonheur, constate Aristote dès les premières lignes de l'\emph{Éthique à Nicomaque} ; mais ils ne s'accordent ni sur sa nature ni sur les moyens de l'atteindre. \fg{} Définitions : bonheur (accomplissement durable de la vie), chance (ce qui advient sans dépendre de nous), raison (faculté de se conduire soi-même). Tension : si le bonheur dépend de la chance, il est fragile et injustement distribué ; s'il dépend de la raison, il est à la portée de tout homme vertueux — mais l'expérience semble démentir cette promesse. Problématique : \textbf{la raison suffit-elle à faire le bonheur, ou le bonheur reste-t-il livré à la fortune ?} Annonce du plan : nous verrons que la sagesse fait du bonheur une œuvre de la raison ; nous montrerons ensuite que la fortune impose ses limites ; nous dégagerons enfin l'idée d'un bonheur où la raison domestique la chance sans l'abolir.

\textbf{Étape 3 — Développement (10 pts).}
\begin{itemize}
\item \textbf{Thèse (le bonheur relève de la raison) :} (1) Épicure : le bonheur naît de la maîtrise des désirs — supprimer les désirs vains suffit à être heureux (ex. : se contenter de peu). (2) Le stoïcisme (Épictète, Marc Aurèle) : distinguer ce qui dépend de nous (nos jugements) et ce qui n'en dépend pas ; le sage reste heureux même dans l'épreuve (ex. : Épictète, esclave devenu philosophe). (3) Kant : seule la volonté raisonnable, agissant par devoir, rend digne du bonheur ; la moralité est en notre pouvoir.
\item \textbf{Antithèse (le bonheur dépend de la chance) :} (1) Aristote lui-même reconnaît que le bonheur exige des biens extérieurs : santé, amis, ressources — \og nul ne dirait heureux un homme sur le chevalet \fg{}. (2) Les conditions de naissance pèsent sur le destin : un enfant né dans une région en guerre ou dans l'extrême pauvreté ne part pas avec les mêmes chances (ex. : inégalités d'accès à l'éducation au Cameroun). (3) La maladie, le deuil, les catastrophes détruisent les bonheurs les plus sages : la fortune (\emph{tyché}) gouverne une part irréductible de la vie.
\item \textbf{Synthèse :} l'alternative est à dépasser. Aristote définit le bonheur comme \og activité de l'âme conforme à la vertu, dans une vie complète \fg{} : la raison est le principe directeur, mais elle a besoin d'un minimum de biens extérieurs. Alain ajoute que le bonheur est une \emph{conquête}, une décision de vouloir être heureux malgré l'incertain : \og il faut vouloir être heureux et y mettre du sien \fg{}. La raison ne supprime pas la chance ; elle la \emph{domestique} : elle transforme les épreuves en occasions, et reçoit les biens de la fortune sans s'y attacher.
\end{itemize}

\textbf{Étape 4 — Conclusion (4 pts).} Réponse directe : le bonheur est affaire de raison en son principe, de chance en ses conditions ; il est l'œuvre d'une vie conduite avec lucidité au milieu de l'incertain. Rappel du parcours dialectique en trois temps. Ouverture : si le bonheur dépend aussi des conditions extérieures, la justice sociale devient une condition du bonheur de tous — question qui relève de la philosophie politique.

\textbf{Étape 5 — Relecture (2 pts).} Orthographe et syntaxe ; présence des connecteurs ; conformité au plan annoncé ; aucune partie manquante.

\begin{attention}[Points de vigilance au Bac]
Citer les auteurs avec précision mais brièvement ; ne jamais écrire \og dans ma thèse... \fg{} ; chaque argument doit être suivi d'un exemple ; la synthèse ne doit pas être un simple compromis (\og un peu des deux \fg{}) mais un dépassement argumenté ; gérer le temps : ne pas sacrifier la conclusion.
\end{attention}
\end{corrige}

\begin{corrige}[Exercice 10 — Analyse et correction d'introductions]
\textbf{1. Analyse des trois introductions.}

\textbf{Introduction A.} \emph{Points forts :} une tentative d'accroche ; une annonce de plan en trois temps ; un exemple concret (Code pénal camerounais). \emph{Défauts majeurs :} accroche banale (\og depuis l'Antiquité \fg{}) ; aucune définition des termes (\og justice \fg{}, \og loi \fg{}) ; problématique réduite à la reprise du sujet ; annonce de plan scolaire (\og d'abord oui... ensuite non... \fg{}) qui livre les réponses avant de les démontrer ; la synthèse (\og la justice est l'équité \fg{}) est affirmée sans justification.

\textbf{Introduction B.} \emph{Points forts :} accroche référencée (Rawls) ; étymologie et définitions précises des deux termes clés ; mise en tension réelle (légalité vs équité, exemple des régimes totalitaires) ; problématique ouverte et formulée ; annonce de plan claire avec arguments directeurs. \emph{Défaut mineur :} l'annonce du plan nomme explicitement \og Thèse / Antithèse / Synthèse \fg{}, ce qui est à éviter ; la formulation gagnerait à rester dans le contenu.

\textbf{Introduction C.} \emph{Points forts :} une intuition juste (l'égalité de traitement) ; un souci de concret. \emph{Défauts majeurs :} pas d'accroche ; définitions enfantines et fausses (\og c'est quand tout le monde est traité pareil \fg{} confond justice et stricte égalité) ; problématique orale et familière (\og Est-ce que... \fg{}) ; plan numéroté comme une liste, avec une troisième partie hors sujet philosophique (\og la vraie justice c'est l'amour \fg{}, affirmation gratuite) ; ton familier incompatible avec une dissertation.

\textbf{2. Classement :} \textbf{B} (la meilleure : méthode complète, seule l'étiquette du plan est à corriger) $>$ \textbf{A} (structure présente mais étapes creuses) $>$ \textbf{C} (la plus faible : aucune exigence méthodologique respectée).

\textbf{3. Réécriture modèle de l'introduction C :}

\medskip
\og Rendre à chacun son dû \fg{} : telle est, depuis le \emph{Corpus} du droit romain, la définition classique de la justice. Mais qui décide de ce qui est \og dû \fg{} ? Le plus souvent, c'est la loi, règle écrite et obligatoire, édictée par l'autorité légitime et identique pour tous. La \cle{justice} désigne, elle, l'idéal d'équité qui donne à chacun ce qui lui revient, et protège en priorité les plus faibles. Or ces deux notions coïncident-elles toujours ? Si la loi est la même pour tous, elle semble garantir l'égalité ; mais une loi peut être injuste, et l'histoire offre des exemples de persécutions parfaitement légales. On peut alors se demander : \textbf{la justice se réduit-elle au respect des lois, ou exige-t-elle que la loi elle-même soit conforme à l'équité ?} Nous verrons d'abord que le respect de la loi est la condition première de la justice, car sans règle commune c'est la loi du plus fort qui s'impose aux plus faibles. Nous montrerons ensuite que la légalité ne suffit pas : une loi injuste demeure une injustice, fût-elle votée. Nous dégagerons enfin l'idée que la justice véritable est l'accord entre la loi et l'équité, accord que seule une réflexion morale sur le droit peut instituer.

\medskip
\textbf{Points de vigilance :} on a conservé les intuitions de l'élève (égalité, protection des faibles) en les formalisant ; problématique unique ; plan annoncé par son contenu.
\end{corrige}

\begin{corrige}[Exercice 11 — Sujet d'invention : Dissertation à partir d'un dossier]
\textbf{Corrigé détaillé (plan et passages rédigés attendus).}

\textbf{Introduction.} L'accroche exploite obligatoirement le Doc 3 : \og \emph{La loi est comme la toile d'araignée : elle retient les petits insectes mais les gros la traversent}, dit un proverbe bamiléké. Cette sagesse populaire exprime un soupçon universel : la loi, égale en principe, serait inégale en fait. \fg{} Définitions : la \cle{loi} (règle générale, obligatoire, édictée par l'autorité légitime) ; la \cle{justice} (idéal d'équité : rendre à chacun son dû) ; \og fonder \fg{} (constituer le fondement suffisant). Tension : si la loi est la règle commune, son respect semble garantir la justice ; mais si la loi épargne les puissants et frappe les faibles, le respect de la loi fonde-t-il encore la justice ? Problématique : \textbf{le respect de la loi suffit-il à fonder la justice, ou la justice exige-t-elle que la loi soit elle-même juste et également appliquée ?} Annonce du plan : la loi comme forme universelle de la justice ; les écarts entre légalité et justice réelle ; les conditions d'une justice effective.

\textbf{Thèse — la loi fonde la justice (appui sur le Doc 1).} (1) La loi est universelle : elle s'impose à tous de la même manière, ce qui réalise l'égalité formelle. Kant le formule dans l'impératif catégorique : \og Agis uniquement d'après la maxime qui fait que tu puisses vouloir en même temps qu'elle devienne une loi universelle \fg{} (Doc 1) : l'exigence d'universalité est le cœur même de la justice. (2) Sans loi commune, c'est l'arbitraire et la vengeance privée : la loi substitue le jugement public à la force privée (Hobbes). (3) Le respect de la loi protège les faibles : là où règne l'État de droit, le citoyen ordinaire peut faire valoir ses droits contre plus puissant que lui.

\textbf{Antithèse — le respect de la loi ne suffit pas (appui sur les Docs 2 et 3).} (1) L'application de la loi peut être inégale : selon la CONAC, 150 milliards de FCFA sont détournés chaque année au Cameroun, et \og l'impunité reste la règle pour les hauts fonctionnaires \fg{} (Doc 2) : la loi existe, mais elle ne frappe pas également. (2) Le proverbe bamiléké (Doc 3) exprime exactement cette dissociation : la toile retient les petits, les gros la traversent — la légalité formelle coexiste avec l'injustice réelle. (3) Une loi peut être injuste en son contenu : l'apartheid ou les lois discriminatoires étaient des lois ; les respecter ne fondait aucune justice (sophiste Calliclès, Antigone chez Sophocle, désobéissance civile).

\textbf{Synthèse — articuler légalité et équité.} La justice exige la loi, mais une loi juste et également appliquée. Rawls définit la justice comme \og la première vertu des institutions sociales \fg{} et la conçoit comme équité : les règles doivent être telles que tous puissent les accepter, y compris les plus défavorisés. La tradition du \cle{droit naturel} (Cicéron, Thomas d'Aquin) ajoute qu'au-dessus des lois positives existe une exigence morale qui permet de les juger. Concrètement, cela suppose des institutions indépendantes : séparation des pouvoirs (Montesquieu), justice réellement impartiale, organismes de contrôle effectifs comme la CONAC dotée de moyens réels de sanction.

\textbf{Conclusion.} Réponse directe : non, le respect de la loi ne suffit pas à fonder la justice ; il en est la condition nécessaire mais non suffisante — il faut encore que la loi soit juste et appliquée à tous sans exception. Rappel du parcours : la loi comme forme universelle (Kant), ses défaillances réelles (CONAC, proverbe), l'exigence d'équité institutionnelle (Rawls). Ouverture : la justice n'est pas seulement l'affaire des juges ; elle suppose la vigilance du \cle{citoyen} et de la \cle{société civile} — presse libre, associations anticorruption, éducation civique — sans lesquelles les meilleures lois restent des toiles d'araignée.

\medskip
\textbf{Barème indicatif (sur 20) :} introduction exploitant le proverbe et problématisant (4 pts) ; thèse fondée sur Kant, correctement citée et commentée (4 pts) ; antithèse mobilisant les données chiffrées du Doc 2 et le sens du proverbe (4 pts) ; synthèse articulée (Rawls, droit naturel, institutions) (4 pts) ; conclusion avec ouverture citoyenne (2 pts) ; cohérence, connecteurs, langue (2 pts).

\begin{attention}[Points de vigilance]
Les trois documents doivent être \emph{cités et exploités}, pas seulement mentionnés ; les données chiffrées (150 milliards FCFA) doivent être reprises exactement ; éviter le commentaire de documents successif (Doc 1, puis Doc 2, puis Doc 3) : les documents doivent être intégrés à une \emph{démonstration dialectique}.
\end{attention}
\end{corrige}

\newpage

\coursec{Fiche bilan}

\begin{popfigure}
\begin{center}\tcbox[colback=popink,colframe=popink,coltext=white,arc=9pt,boxsep=4pt,left=6pt,right=6pt]{\bfseries\small La Dissertation Philosophique}\end{center}
\vspace{-2pt}
\begin{tcbraster}[raster columns=3,raster equal height=rows,raster column skip=6pt,raster row skip=6pt,enhanced,blankest]
\begin{tcolorbox}[enhanced,colback=popblue,colframe=popblue,coltext=white,boxrule=0.9pt,arc=3pt,left=4pt,right=4pt,top=3pt,bottom=3pt,boxsep=1pt,halign=left]\bfseries\scriptsize Généralités Définition, But, Types de sujets\end{tcolorbox}
\begin{tcolorbox}[enhanced,colback=poporange,colframe=poporange,coltext=white,boxrule=0.9pt,arc=3pt,left=4pt,right=4pt,top=3pt,bottom=3pt,boxsep=1pt,halign=left]\bfseries\scriptsize Introduction Accroche, Analyse, Problématique, Plan\end{tcolorbox}
\begin{tcolorbox}[enhanced,colback=popgreen,colframe=popgreen,coltext=white,boxrule=0.9pt,arc=3pt,left=4pt,right=4pt,top=3pt,bottom=3pt,boxsep=1pt,halign=left]\bfseries\scriptsize Thèse Arguments POUR, Exemples, Auteurs\end{tcolorbox}
\begin{tcolorbox}[enhanced,colback=poppurple,colframe=poppurple,coltext=white,boxrule=0.9pt,arc=3pt,left=4pt,right=4pt,top=3pt,bottom=3pt,boxsep=1pt,halign=left]\bfseries\scriptsize Antithèse Arguments CONTRE, Nuances, Limites\end{tcolorbox}
\begin{tcolorbox}[enhanced,colback=poppink,colframe=poppink,coltext=white,boxrule=0.9pt,arc=3pt,left=4pt,right=4pt,top=3pt,bottom=3pt,boxsep=1pt,halign=left]\bfseries\scriptsize Synthèse Dépassement, Réponse, Ouverture\end{tcolorbox}
\begin{tcolorbox}[enhanced,colback=popgold,colframe=popgold,coltext=white,boxrule=0.9pt,arc=3pt,left=4pt,right=4pt,top=3pt,bottom=3pt,boxsep=1pt,halign=left]\bfseries\scriptsize Conclusion Bilan, Réponse directe, Perspective\end{tcolorbox}
\begin{tcolorbox}[enhanced,colback=popdark,colframe=popdark,coltext=white,boxrule=0.9pt,arc=3pt,left=4pt,right=4pt,top=3pt,bottom=3pt,boxsep=1pt,halign=left]\bfseries\scriptsize Méthodologie Brouillon, Temps, Relecture, Style\end{tcolorbox}
\end{tcbraster}

\legende{Carte mentale : structure et étapes clés de la dissertation philosophique en Première Technique.}
\end{popfigure}

\begin{bilan}

\textbf{Définitions clés}
\begin{itemize}
  \item \cle{Dissertation} : Exercice écrit de réflexion argumentée répondant à une question philosophique (sujet) par une démarche dialectique (thèse / antithèse / synthèse).
  \item \cle{Sujet} : Question explicite ou citation à discuter. Types : \emph{sujet fermé} (réponse oui/non), \emph{sujet ouvert} (``Qu'est-ce que... ?''), \emph{citation} (à expliquer et discuter).
  \item \cle{Problématique} : Tension philosophique mise en lumière par l'analyse du sujet ; elle oppose au moins deux thèses légitimes et guide le plan.
  \item \cle{Argument} : Raison logique (concept, principe, référence d'auteur) qui soutient une thèse. Distinguer \emph{argument d'autorité} (référence) et \emph{argument rationnel} (logique).
  \item \cle{Exemple} : Illustration concrète (fait social, historique, littéraire, camerounais) qui ancre l'argument dans le réel. Ne prouve pas seul, il \emph{illustre}.
  \item \cle{Plan dialectique} : Structure en trois temps : \textbf{Thèse} (réponse spontanée/majoritaire), \textbf{Antithèse} (critique/limites/opposée), \textbf{Synthèse} (dépassement/réponse nuancée).
\end{itemize}

\textbf{Méthodes express (fiches mémo)}

\begin{center}
\begin{tabularx}{\linewidth}{|l|X|}\hline
\rowcolor{popblueL}
\textbf{Étape} & \textbf{Actions clés} \\ \hline
\textbf{1. Analyse (brouillon)} & \begin{itemize}\itemsep0pt \parskip0pt
  \item Souligner mots-clés, connecteurs, présupposés.
  \item Identifier le type de sujet et le domaine (éthique, politique, épistémologie...).
  \item Lister thèses possibles (POUR / CONTRE) + auteurs + exemples.
  \item Formuler la \cle{problématique} en question tensionnelle.
\end{itemize} \\ \hline
\textbf{2. Plan détaillé} & \begin{itemize}\itemsep0pt \parskip0pt
  \item \textbf{Thèse} : 2--3 arguments forts + exemples + auteurs.
  \item \textbf{Antithèse} : 2--3 objections radicales + exemples + auteurs.
  \item \textbf{Synthèse} : Critère de dépassement (hiérarchie, condition, distinction conceptuelle) $\rightarrow$ réponse définitive.
  \item Rédiger \emph{transitions} et \emph{sous-conclusions}.
\end{itemize} \\ \hline
\textbf{3. Introduction} & \begin{itemize}\itemsep0pt \parskip0pt
  \item \textbf{Accroche} : citation, fait d'actualité (Cameroun), définition, paradoxe.
  \item \textbf{Analyse du sujet} : définition termes, reformulation, enjeu.
  \item \textbf{Problématique} : question unique, claire, issue de l'analyse.
  \item \textbf{Annonce du plan} : ``Nous verrons d'abord que... Ensuite que... Enfin...''
\end{itemize} \\ \hline
\textbf{4. Développement} & \begin{itemize}\itemsep0pt \parskip0pt
  \item Un paragraphe = un argument = (Idée directrice + Explication + Exemple + Auteur + Mini-conclusion).
  \item Respecter l'ordre : Thèse $\rightarrow$ Antithèse $\rightarrow$ Synthèse.
  \item Transitions explicites : ``Cependant...'', ``Toutefois...'', ``Au terme de cette analyse...''
\end{itemize} \\ \hline
\textbf{5. Conclusion} & \begin{itemize}\itemsep0pt \parskip0pt
  \item \textbf{Bilan synthétique} : résumé articulé des 3 parties (pas liste).
  \item \textbf{Réponse directe} à la problématique (affirmative, nuancée, conditionnelle).
  \item \textbf{Ouverture} : autre domaine, auteur, question nouvelle, actualité.
\end{itemize} \\ \hline
\end{tabularx}
\end{center}

\textbf{Repères auteurs incontournables (programme 1re Tech)}
\begin{itemize}
  \item \textbf{Conscience / Sujet} : Descartes (\emph{Cogito}), Kant (\emph{Moi transcendantal}), Bergson (\emph{Durée}), Sartre (\emph{Liberté}).
  \item \textbf{Science / Vérité} : Descartes (\emph{Règles}), Popper (\emph{Falsifiabilité}), Kuhn (\emph{Paradigme}), Bachelard (\emph{Obstacle épistémologique}).
  \item \textbf{Politique / Justice} : Hobbes (\emph{Léviathan}), Locke (\emph{Droits naturels}), Rousseau (\emph{Contrat social}), Rawls (\emph{Justice comme équité}), Sen (\emph{Capabilités}).
  \item \textbf{Moral / Devoir} : Kant (\emph{Impératif catégorique}), Bentham/Mill (\emph{Utilitarisme}), Aristote (\emph{Vertu}), Levinas (\emph{Visage d'autrui}).
  \item \textbf{Technique / Travail} : Marx (\emph{Travail aliéné}), Heidegger (\emph{Essor de la technique}), Ellul (\emph{Système technicien}), Simondon (\emph{Individuation technique}).
\end{itemize}

\end{bilan}

\begin{aretenir}[Checklist avant le Bac (Probatoire)]
\begin{itemize}
  \item $\square$ Sujet \textbf{analysé} (mots-clés, type, présupposés) et \textbf{problématisé} (tension réelle) ?
  \item $\square$ Introduction \textbf{complète} : accroche pertinente + analyse rigoureuse + problématique unique + plan annoncé ?
  \item $\square$ Plan \textbf{dialectique} respecté : Thèse (arguments POUR) / Antithèse (arguments CONTRE) / Synthèse (dépassement) ?
  \item $\square$ Arguments \textbf{philosophiques} (concepts, distinctions, auteurs du programme) et non seulement opinions ?
  \item $\square$ Exemples \textbf{concrets et variés} dont au moins un \textbf{contextualisé Cameroun} (coutumes, actualité, droit OHADA, biodiversité, urbanisation...) ?
  \item $\square$ Transitions \textbf{explicites} entre parties et paragraphes (liaisons logiques : concession, opposition, conséquence) ?
  \item $\square$ Synthèse \textbf{réelle} : critère de dépassement clair (hiérarchie, condition, distinction) $\neq$ simple résumé ?
  \item $\square$ Conclusion \textbf{efficace} : bilan articulé + réponse directe à la problématique + ouverture pertinente ?
  \item $\square$ Langage \textbf{philosophique} : vocabulaire précis (\cle{subjectivité}, \cle{intersubjectivité}, \cle{aliénation}, \cle{finalité}...), syntaxe soignée, pas de familier ?
  \item $\square$ \textbf{Gestion du temps} : 30--40 min brouillon / 2h30--3h rédaction / 15--20 min relecture ?
  \item $\square$ \textbf{Relecture} : orthographe, accord, cohérence du fil conducteur, respect du sujet (hors-sujet éliminé) ?
  \item $\square$ Présentation \textbf{soignée} : paragraphes visibles, ratures propres, lisibilité pour le correcteur ?
\end{itemize}
\end{aretenir}

\findecours

\end{document}
